% =====================================================================
%  第 12 章  Drell--Yan 过程与强子—强子碰撞            （ch12_dy.tex，writer W14）
%  约定：本书约定（WRITER_SPEC_v2 §3）；n_A=(1^+,0^-,0_T)，n_B=(0^+,1^-,0_T)；
%  D_mu = d_mu - i g A_mu^a t^a；D = 4-2 eps，S_eps=(4pi)^eps e^{-eps gamma_E}；a=alpha_s/(2pi)。
% =====================================================================
\chapter{Drell--Yan 过程与强子—强子碰撞}\label{ch:dy}

在 Drell--Yan 过程 $A(P_A)+B(P_B)\to\gamma^*(q)+X\to\ell^+\ell^-+X$ 中，两个高能强子各提供一个部分子，二者湮灭成一个大质量的虚光子，我们只测量其衰变产生的轻子对。这是两个强子初态的硬过程中最简单的一个，也是因子化理论中最微妙的一个。它与前几章的过程有两点本质区别。第一，它不是 OPE 可及的量：截面不是单个强子态上流算符乘积的前向矩阵元，\cref{thm:dis-moments}（\tagR，$R_{\mathrm{OPE}}$）所依赖的色散关系与矩方法在这里都不适用，证明只能在动量空间中逐图进行（“因子化就是 OPE”是一个常见的误解）。第二，两个强子的旁观者（spectator）在硬碰撞之前与之后都存在，它们之间可以交换 Glauber 胶子；在单个割图中这些交换的积分围道被钳住（钳夹，pinch），软近似在实轴上失效。全阶论证的核心正是如何处理这一 Glauber 区。

本章在本书的标签体系中给出如下结果。在重整化微扰论中严格（$R_{\mathrm{PT}}$）的有：Born 截面（\cref{prop:dy-born}）、单个割图中的单圈 Glauber 钳夹（\cref{prop:dy-glauber-pinch}）、光锥序中间态的取向与 $k^-$ 极点（\cref{lem:dy-poles-orientation}）、两条程函线（eikonal line）之间的双胶子交换及其 $\order{g^4}$ 幺正性（\cref{prop:dy-eikonal-twogluon}）、$\order{\as}$ 的 $\MSbar$ 系数函数与 $\mu_F$ 无关性（\cref{prop:dy-nlo,prop:dy-scale}）。由定义直接推出（$R_{\mathrm{def}}$）的有：单个割图给出指向未来的 Wilson 线（\cref{ex:dy-future-vs-past}），以及 DY 与 DIS 的共线部分子分布是同一个算符矩阵元（\cref{prop:dy-universality}）。在阿贝尔程函模型中，我们完整地证明了旁观者 Glauber 相位的相消（\cref{prop:dy-abelian-model}，作为 QCD 陈述为 \tagP）。对单举 DY 的共线因子化（\cref{thm:dy-factorization}）以及色单态末态、极化 DY 与 $1/Q^2$ 处的广义因子化，本章给出全阶论证（\tagA），并以四个层次说明其地位：文献中的论证；\cref{app:appO}与\cref{app:appN}中把它写成的精确条件定理链；以标签列出的、尚未证明的输入；以及“没有已知的领头幂次反例”。有色末态的因子化只是猜想（\tagC），$1/Q^2$ 之外的幂次修正上因子化被破坏（\tagV）。

\begin{tcolorbox}[enhanced,breakable,title=本章证明了什么／没有证明什么,fonttitle=\bfseries,
  colback=gray!3!white,colframe=gray!55!black,boxrule=0.6pt,arc=2pt,
  left=5pt,right=5pt,top=3pt,bottom=3pt,before skip=6pt,after skip=6pt]
\textbf{证明了：}\cref{prop:dy-born,prop:dy-glauber-pinch,lem:dy-poles-orientation,prop:dy-eikonal-twogluon,prop:dy-nlo,prop:dy-scale}（\tagR，$R_{\mathrm{PT}}$）；\cref{prop:dy-universality,ex:dy-future-vs-past}（\tagR，$R_{\mathrm{def}}$）；\cref{thm:appO-main,prop:appN-dy}（\tagR，$R_{\mathrm{cond}}$，证明在\cref{app:appO}与\cref{app:appN}，前提为\cref{asm:appO-lcopt}等编号假设）；\cref{prop:dy-abelian-model}（\tagP：阿贝尔程函模型中的完整证明）。\par
\textbf{给出全阶论证：}\cref{prop:dy-regions,lem:dy-fsi,prop:dy-deformation,lem:dy-soft,thm:dy-factorization,thm:dy-colour-singlet,thm:dy-polarized,thm:dy-higher-twist}（\tagA；瓶颈\cref{asm:appO-mixing}(c)、\cref{asm:appO-deformation}(b)--(d)、\cref{asm:appO-rewritten-cs}）。\par
\textbf{没有证明：}非程函（复合、相互作用）旁观者与活动部分子的 Glauber--软--共线混合（\cref{asm:appO-mixing}(c)）；余项的幂次 $p$（\tagC）；有色末态（\cref{conj:dy-coloured-final}，\tagC）；$1/Q^2$ 以外的幂次（\cref{prop:dy-twist6-failure}，\tagV）。
\end{tcolorbox}

本章的结构如下。\cref{sec:dy-born}固定运动学与 Born 截面。\cref{sec:dy-regions}列出领头区域，并在单圈上显式地展示 Glauber 钳夹。\cref{sec:dy-orientation}把“极点在哪一侧”化为光锥序中间态的一个性质。\cref{sec:dy-glauber}给出单举 DY 中 Glauber 问题的解决：割图求和的幺正性、单个割图中指向未来的 Wilson 线、同时变形、软因子与部分子分布的普适性。\cref{sec:dy-model}在阿贝尔模型中把这一切算到底，并计算两条程函线之间的双胶子交换。\cref{sec:dy-theorem}陈述主定理并说明其地位。\cref{sec:dy-nlo}推导 $\order{\as}$ 系数函数。\cref{sec:dy-other}与\cref{sec:dy-twist}讨论其他强子—强子过程与高扭度。章末的\cref{tab:dy-correspondence}列出本章陈述与\cref{app:appO}、\cref{app:appN}中精确形式的对应。

% =====================================================================
\section{运动学与 Born 截面}\label{sec:dy-born}
% =====================================================================

Drell 与 Yan 在部分子模型中提出了这一过程的基本图像\cite{DrellYan1970}：强子 $A$ 中的一个夸克与强子 $B$ 中的一个反夸克（或反过来）湮灭成虚光子，其余部分子作为旁观者飞走而不参与硬过程。在强子质心系中，两个强子都被 Lorentz 收缩并且时间膨胀，硬湮灭只持续 $\sim1/Q$ 的时间，在这段时间里部分子之间的相互作用来不及改变它们的分布；于是截面应当是对部分子对的非相干求和，每个部分子以它在单个强子中的分布出现。本章要回答的问题是：这一图像在 QCD 的相互作用之下是否仍然成立，在什么意义下成立。最敏感的一点是两个强子的旁观者：它们在硬碰撞之前就已经彼此接近，部分子模型假定它们之间的相互作用不影响湮灭过程。本节先固定运动学与归一化。

\begin{definition}[DY 运动学]\label{def:dy-kin}
考虑 $A(P_A)+B(P_B)\to\gamma^*(q)+X$，$\gamma^*\to\ell^+\ell^-$（只计虚光子交换，一种无质量轻子）。在强子质心系中强子 $A$ 沿 $+z$、$B$ 沿 $-z$ 运动，
$P_A=(P_A^+,\frac{M_A^2}{2P_A^+},\mathbf 0_T)$，$P_B=(\frac{M_B^2}{2P_B^-},P_B^-,\mathbf 0_T)$（光锥分量按\cref{def:prelim-lightfront}）。
\begin{enumerate}[label=(\roman*),leftmargin=2.2em,itemsep=1pt,topsep=2pt]
\item $s=(P_A+P_B)^2$，$Q^2=q^2$，$\tau=Q^2/s$；轻子对的快度 $y=\frac12\ln(q^+/q^-)$。
\item $x_A=\sqrt\tau\,\ee^{y}$，$x_B=\sqrt\tau\,\ee^{-y}$，$0\le\tau\le1$，$|y|\le\frac12\ln\frac1\tau$。
\item 部分子动量 $p_1=\xi_AP_A$、$p_2=\xi_BP_B$（无质量部分子，忽略强子质量），$\hat s=(p_1+p_2)^2=\xi_A\xi_Bs$，$z=Q^2/\hat s$。
\item 部分子截面的归一化：记 $H^\mu$ 为部分子过程 $p_1+p_2\to\gamma^*(\mu)+X$ 的振幅，其中去掉光子传播子与轻子电荷，保留夸克电荷 $e_q$ 与全部 QCD 耦合；$\langle\cdot\rangle$ 表示对初态的自旋与颜色求平均（夸克自旋 $\frac12$、颜色 $\frac1{N_c}$；胶子极化 $\frac{1}{D-2}$、颜色 $\frac1{N_c^2-1}$）并对末态 $X$ 求和；$\dd\Phi_{\gamma^*X}$ 是含 $(2\pi)^D\delta^D(p_1+p_2-q-p_X)$、并固定 $q^2=Q^2$ 的 $D$ 维相空间。令 $\sigma_0\equiv\frac{4\pi\alpha^2}{3N_cQ^2}$，部分子道 $ij$ 的裸系数定义为
\begin{equation}\label{eq:dy-bare}
\Delta^{\rm bare}_{ij}(z,\varepsilon)\equiv\frac{N_c}{2\pi e_q^2(1-\varepsilon)}\int\dd\Phi^{(D)}_{\gamma^*X}\,\big\langle-g_{\mu\nu}H^\mu H^{\nu*}\big\rangle .
\end{equation}
% VF-DY-2
\end{enumerate}
\uses{def:prelim-lightfront}
\end{definition}

定义中的几点说明。由 $q^+=x_1P_A^+$、$q^-=x_2P_B^-$（$x_{1,2}$ 为两个湮灭部分子的动量分数，最低阶时 $q=p_1+p_2$）得 $Q^2=2q^+q^-=x_1x_2\,(2P_A^+P_B^-)$，而 $2P_A^+P_B^-=s\,[1+\order{M^2/s}]$；在 $P_A^+=P_B^-$ 的系中又有 $y=\frac12\ln(x_1/x_2)$。所以最低阶时 $x_1=x_A$、$x_2=x_B$ 正是 (ii) 中由测量量 $(Q^2,y)$ 确定的值，强子质量只以 $x_Ax_B=\tau[1+\order{M^2/s}]$ 的形式出现。\cref{eq:dy-bare}中除以的因子 $(1-\varepsilon)$ 是 $D$ 维 Born 项的自旋迹因子，它使 $\Delta^{(0)}_{q\bar q}=\delta(1-z)$（\cref{prop:dy-born}）；质量因子化的抵消项乘的正是 $D$ 维的 Born 截面，因此只有这样归一，$\MSbar$ 系数函数才不含人为的 $\pm2P_{qq}$（\cref{sec:dy-nlo}）。

\paragraph{主公式。}轻子张量把强子部分与轻子部分分开。对轻子自旋求和，$|M|^2=\frac{e^4}{Q^4}L_{\mu\nu}H^\mu H^{\nu*}$，$L_{\mu\nu}=\Tr[\sla\ell_1\gamma_\mu\sla\ell_2\gamma_\nu]$。在四维中对轻子的两体相空间积分时，Lorentz 协变与 $q^\mu L_{\mu\nu}=0$ 只允许 $\int\dd\Phi_2L_{\mu\nu}=A\,(q_\mu q_\nu-q^2g_{\mu\nu})$；用 $g^{\mu\nu}$ 缩并，左边为 $\int\dd\Phi_2(-8\ell_1\cdot\ell_2)=-\frac{4Q^2}{8\pi}$（无质量两体相空间的体积为 $\frac1{8\pi}$），右边为 $-3Q^2A$，故
\begin{equation}\label{eq:dy-leptontensor}
\int\dd\Phi_2\,L_{\mu\nu}=\frac{1}{6\pi}\big(q_\mu q_\nu-q^2g_{\mu\nu}\big).
\end{equation}
% VF-DY-2
轻子张量对所有部分子道相同，并且在系数函数中消去，所以始终取在四维中。再用 $\frac{\dd^4q}{(2\pi)^4}=\frac{\dd Q^2}{2\pi}\,\frac{\dd^3q}{(2\pi)^32q^0}$、通量 $1/(2\hat s)$ 与流守恒 $q_\mu H^\mu=0$（因而只剩 $-Q^2g_{\mu\nu}$ 一项），得
\[
\frac{\dd\hat\sigma}{\dd Q^2}=\frac{1}{2\hat s}\,\frac{e^4}{Q^4}\,\frac{1}{2\pi}\,\frac{Q^2}{6\pi}\int\dd\Phi_{\gamma^*X}\big\langle-g_{\mu\nu}H^\mu H^{\nu*}\big\rangle ,
\]
以 $e^2=4\pi\alpha$ 代入即
\begin{equation}\label{eq:dy-master}
\hat s\,\frac{\dd\hat\sigma}{\dd Q^2}=\frac{2\alpha^2}{3Q^2}\int\dd\Phi_{\gamma^*X}\,\big\langle-g_{\mu\nu}H^\mu H^{\nu*}\big\rangle ,\qquad
\Delta^{\rm bare}_{ij}=\frac{\hat s\,\dd\hat\sigma_{ij}/\dd Q^2}{\sigma_0\,e_q^2\,(1-\varepsilon)} .
\end{equation}
% VF-DY-2

\begin{proposition}[Born 截面 \tagR]\label{prop:dy-born}
（$R_{\mathrm{PT}}$：在重整化微扰论中对红外正规化的部分子态严格。）本命题是最低阶（$\as^0$）的陈述。
\begin{enumerate}[label=(\alph*),leftmargin=2.2em,itemsep=1pt,topsep=2pt]
\item 对 $q(p_1)\bar q(p_2)\to\gamma^*$，$\hat s\,\dd\hat\sigma^{(0)}/\dd Q^2=\sigma_0\,e_q^2\,(1-\varepsilon)\,\delta(1-z)$，即 $\Delta^{(0)}_{q\bar q}=\delta(1-z)$；在 $D=4$ 中对 $Q^2$ 积分给出 $\hat\sigma(q\bar q\to\ell^+\ell^-)=\dfrac{4\pi\alpha^2e_q^2}{3N_c\hat s}$。
\item 若强子截面具有\cref{def:intro-factorization}的因子化形式，则其最低阶为
\begin{align}
\frac{\dd\sigma}{\dd Q^2}&=\frac{4\pi\alpha^2}{3N_cQ^2s}\sum_qe_q^2\int_0^1\!\dd x_1\,\dd x_2\,\delta(x_1x_2-\tau)\,\big[q(x_1)\bar q(x_2)+\bar q(x_1)q(x_2)\big],\label{eq:dy-LO}\\
\frac{\dd^2\sigma}{\dd Q^2\,\dd y}&=\frac{4\pi\alpha^2}{3N_cQ^2s}\sum_qe_q^2\,\big[q(x_A)\bar q(x_B)+\bar q(x_A)q(x_B)\big],\label{eq:dy-LOy}
\end{align}
% VF-DY-1
其中 $q(x)=f_{q/A}(x)$（\cref{def:pdf-quark}），$\bar q(x)=-f_{q/A}(-x)$（\cref{def:pdf-antiquark}），第一个因子属于强子 $A$、第二个属于 $B$；从 $(x_1,x_2)$ 到 $(x_1x_2,y)$ 的变量变换的 Jacobi 行列式等于 $1$。
\item （Drell--Yan 标度律）$Q^4\,\dd\sigma/\dd Q^2=\frac{4\pi\alpha^2}{3N_c}\,\tau\sum_qe_q^2\int_\tau^1\frac{\dd x}{x}\big[q(x)\,\bar q(\tfrac\tau x)+\bar q(x)\,q(\tfrac\tau x)\big]$，在不计部分子分布的标度依赖时只是 $\tau$ 的函数。
\end{enumerate}
（(a) 只涉及部分子态；(b)(c) 是在因子化形式之下的代数推论，不涉及强子态的非微扰性质。）
\uses{def:dy-kin,def:pdf-quark,def:pdf-antiquark,def:intro-factorization}
\end{proposition}

\begin{proof}
(a) Born 振幅为 $H^\mu=e_q\,\bar v(p_2)\gamma^\mu u(p_1)\,\delta_{ij}$（$i,j$ 为色指标）。在 $D$ 维中 $\gamma^\mu\sla a\gamma_\mu=-(D-2)\sla a$，故对自旋求和
$-g_{\mu\nu}\sum H^\mu H^{\nu*}=-e_q^2\TrD[\sla p_2\gamma^\mu\sla p_1\gamma_\mu]=(D-2)\,e_q^2\,\TrD[\sla p_2\sla p_1]=(D-2)\,e_q^2\,2\hat s$。
颜色求和 $\sum_{ij}\delta_{ij}\delta_{ij}=N_c$，初态平均 $\frac14\cdot\frac1{N_c^2}$，于是 $\langle-gHH^*\rangle=e_q^2(1-\varepsilon)\hat s/N_c$（$D-2=2(1-\varepsilon)$）。$2\to1$ 的相空间为 $\int\dd\Phi_{\gamma^*}=2\pi\,\delta(\hat s-Q^2)$（在任意 $D$ 中成立）。代入\cref{eq:dy-master}，并用 $\delta(\hat s-Q^2)=\delta(1-z)/\hat s$，
$\hat s\,\dd\hat\sigma^{(0)}/\dd Q^2=\frac{2\alpha^2}{3Q^2}\cdot2\pi\cdot\frac{e_q^2(1-\varepsilon)}{N_c}\,\delta(1-z)=\sigma_0e_q^2(1-\varepsilon)\delta(1-z)$，再由\cref{eq:dy-bare}得 $\Delta^{(0)}_{q\bar q}=\delta(1-z)$。$D=4$ 时 $\dd\hat\sigma^{(0)}/\dd Q^2=\sigma_0(Q^2)\,e_q^2\,\delta(\hat s-Q^2)$，对 $Q^2$ 积分即得 $\sigma_0(\hat s)\,e_q^2=4\pi\alpha^2e_q^2/(3N_c\hat s)$。

(b) 因子化形式的最低阶就是部分子截面对两个强子的部分子分布的卷积：$q$ 来自 $A$、$\bar q$ 来自 $B$ 与反过来两种排列，
$\frac{\dd\sigma}{\dd Q^2}=\sum_qe_q^2\int\dd x_1\dd x_2\,[q(x_1)\bar q(x_2)+\bar q(x_1)q(x_2)]\,\sigma_0\,\delta(x_1x_2s-Q^2)$，而 $\delta(x_1x_2s-Q^2)=\delta(x_1x_2-\tau)/s$，即\cref{eq:dy-LO}。对\cref{eq:dy-LOy}，把 $(x_1,x_2)$ 换成 $(u,y)=(x_1x_2,\frac12\ln\frac{x_1}{x_2})$，其 Jacobi 行列式
$\big|\partial(u,y)/\partial(x_1,x_2)\big|=\big|x_2\cdot(-\tfrac{1}{2x_2})-x_1\cdot\tfrac{1}{2x_1}\big|=1$，所以 $\int\dd x_1\dd x_2\,\delta(x_1x_2-\tau)\,\delta(y-\frac12\ln\frac{x_1}{x_2})F(x_1,x_2)=F(\sqrt\tau\ee^{y},\sqrt\tau\ee^{-y})$。

(c) 在\cref{eq:dy-LO}中先对 $x_2$ 积分，$\delta(x_1x_2-\tau)=\delta(x_2-\tau/x_1)/x_1$，乘以 $Q^4=\tau sQ^2$ 即得所述形式；右边除了部分子分布之外只含 $\tau$。
\end{proof}

\cref{prop:dy-born}(b)(c) 只断言从部分子截面到强子截面的代数步骤；“强子截面确实具有因子化形式”本身就是因子化定理在最低阶的内容，其地位是\cref{thm:dy-factorization}的地位（全阶论证，\tagA）。在 $D=4$ 中，用显式旋量计算还可以得到轻子的角分布 $\propto1+\cos^2\theta$（$\theta$ 为轻子对静止系中轻子相对于碰撞轴的极角），这是自旋 $\frac12$ 的部分子湮灭的特征。标度律 (c) 在 $\order{\as}$ 被对数地破坏：系数函数含 $\ln(Q^2/\mu_F^2)$，与部分子分布的演化合在一起才使截面与 $\mu_F$ 无关（\cref{prop:dy-scale}）。

% =====================================================================
\section{领头区域：初态与末态相互作用}\label{sec:dy-regions}
% =====================================================================

DY 的约化图（\cref{fig:reduced-dy}）与 DIS、$e^+e^-$ 的约化图形式相同：硬子图、两个共线子图与一个软子图。差别在于软胶子所连接的线的“时间性质”。在 $e^+e^-$ 湮灭中，全部共线线都属于末态；在 DIS 中，靶子图 $J_A$ 既含初态线（硬顶点之前的活动部分子）也含末态线（飞向末态的旁观者），但另一个共线子图——被击出的夸克喷注——只含末态线，于是一个软动量的 $k^+$ 分量只看到同一半平面的极点，围道可以变形。在 DY 中，$J_A$ 与 $J_B$ 都同时含有初态线与末态线：两个强子的旁观者之间交换的软胶子在 $k^-$ 上被 $J_A$ 钳住，在 $k^+$ 上被 $J_B$ 钳住。这就是 DY 特有的 Glauber 问题；本节先把领头区域列出来，再在单圈上把钳夹显式算出。

\begin{proposition}[DY 的领头区域 \tagA]\label{prop:dy-regions}
在 Feynman 规范中，单举 DY 截面（对全部割图求和）的领头幂次区域由\cref{thm:ir-leading-regions}（全阶论证，\tagA）的约化图给出（\cref{fig:reduced-dy}）：
\begin{enumerate}[label=(\roman*),leftmargin=2.2em,itemsep=1pt,topsep=2pt]
\item 末态割线两侧各一个硬子图 $H$、$H^\dagger$，其中的线离壳 $\order{Q^2}$；
\item 共线子图 $J_A$、$J_B$，各自含强子顶点与全部旁观者并穿过末态割线，在割线每一侧经一个物理部分子（夸克的好分量或横向胶子）加任意多个纵向极化的共线胶子连到 $H$（$H^\dagger$）；
\item 软子图 $S$，只经胶子连到 $J_A$、$J_B$，不连到 $H$、$H^\dagger$；
\item 作为 (iii) 的子区域，Glauber 动量 $k\sim(\lambda^2,\lambda^2,\lambda)Q$（\cref{def:ir-glauber}），特别是两个强子的旁观者之间的交换，它们与软动量一样是领头幂次。
\end{enumerate}
在单个割图中，$J_A$、$J_B$ 都既含硬顶点之前的初态线，也含硬顶点之后、割线之前的末态线；旁观者的软相互作用因而既有初态相互作用，也有末态相互作用。（本命题只涉及部分子态。）
% errtype: E-S
\uses{thm:ir-leading-regions,def:dy-kin,def:ir-glauber}
\end{proposition}

\begin{proofsketch}
(1) 由 Landau 方程与 Coleman--Norton 图像得到钳夹面，由幂次计数挑出 $\omega=0$ 的那些，得到\cref{thm:ir-leading-regions}的约化图；对 DY 该定理已把旁观者之间的 Glauber 子区域列入（\cite[第~5.8--5.11~节]{Collins2011}，\cite[第~8.7~节]{CollinsSoperSterman1989}）。
(2) Glauber 区确为领头幂次：相对于不含该胶子的图，测度 $\dd k^+\dd k^-\dd^2\mathbf k_T\sim\lambda^6Q^4$，胶子传播子 $\simeq-1/\mathbf k_T^2\sim\lambda^{-2}Q^{-2}$，两条新的旁观者传播子各离壳 $\order{\lambda^2Q^2}$，共 $\lambda^{-4}Q^{-4}$，分子中与大分量的耦合给出 $(2p_A^+)(2p_B^-)\sim Q^2$，乘积为 $\lambda^0$，与 $Q$ 无关\cite[第~5.6.2~节]{Collins2011}（此步只给出 E-S 型陈述）。
(3) 旁观者穿过割线，软胶子在割线同一侧既可附着在硬顶点 $V_H$ 之前的旁观者段，也可附着在其后的段上，于是两种相互作用在同一个割图中出现。
完整论证见\cite[第~5~章与第~14.3~节]{Collins2011}。
\end{proofsketch}

\paragraph{地位与缺口}\tagA。最弱的一步是 (1)–(2) 中的幂次计数：它们是标度陈述（E-S），积分范数下的形式是\cref{asm:appN-envelope}、\cref{asm:appN-separation}（\tagA），在所有阶都没有以积分形式写出。Feynman 规范中领头区域还不完备：额外的不连通硬散射只在对图求和后相消\cite{Labastida:1984gy}，Collins 把这一点列为 DY 证明中“尚存的问题”之一\cite[第~14.4.8~节]{Collins2011}。对强子态，“微扰论找到的领头区域就是全部领头区域”是桥接假设\cref{asm:intro-bridge}(Br1)，不属于本命题。

\cref{fig:glauber-exchange}把 DY 中旁观者之间的初态 Glauber 交换与 DIS 中被击出夸克与靶残余之间的末态交换并列。图中说明了单个割图中极点分居两侧、围道被钳住的情形；需要补充的是：离开 Glauber 区须沿 $e_G$ 同时形变两个分量（\cref{fig:method-glauber-direction}）——只推动 $k^-$ 一个分量，即使推到 $|k^-|\sim|\mathbf k_T|$，也仍有 $|k^+k^-|\sim\lambda^3\ll\mathbf k_T^2$，动量并没有离开 Glauber 区（\cref{rem:method-two-component}）。

%%% BEGIN fig:glauber-exchange
\begin{figure}[htbp]
\centering
\begin{tikzpicture}[fd,
  t2 hadron/.style={draw=black!70,line width=2.2pt,postaction={decorate,decoration={markings,
      mark=at position 0.6 with {\arrow[black!70]{Stealth[length=7pt,width=7.5pt]}}}}},
  t2 lab/.style={font=\footnotesize},
  t2 time/.style={lcaxis,gray!60!black}]
  % ======================= (a) Drell-Yan: initial-state exchange between spectators =======================
  \begin{scope}
    \node[jet blob,circle,minimum size=0.66cm,inner sep=0pt] (A) at (-2.45,1.0) {$A$};
    \node[jetB blob,circle,minimum size=0.66cm,inner sep=0pt] (B) at (-2.45,-1.25) {$B$};
    \node[hard blob,circle,minimum size=0.6cm,inner sep=0pt] (H) at (1.65,1.5) {$H$};
    \draw[t2 hadron] (-3.6,1.0) -- (A.west);
    \draw[t2 hadron] (-3.6,-1.25) -- (B.west);
    % active partons: quark from A, antiquark from B (fermion-number arrow points back into B)
    \draw[quark] (A.12) -- (H.185);
    \draw[quark=0.45] (H.235) -- (B.42);
    \draw[photon] (H.east) -- (2.85,1.5);
    % spectators (remnants) of A and B, continuing into the final state
    \coordinate (SA) at (A.-28); \coordinate (XA) at (2.8,0.25);
    \coordinate (SB) at (B.-8);  \coordinate (XB) at (2.8,-1.2);
    \draw[quark=0.25] (SA) -- (XA);
    \draw[quark=0.25] (SB) -- (XB);
    % Glauber exchange before the hard collision (left of H)
    \coordinate (Gt) at (0.78,2.0); \coordinate (Gb) at (0.78,-2.0);
    \coordinate (GA) at (intersection of SA--XA and Gt--Gb);
    \coordinate (GB) at (intersection of SB--XB and Gt--Gb);
    \draw[glauber] (GA) -- (GB);
    \node[vertex] at (GA) {}; \node[vertex] at (GB) {};
    \node[t2 lab,text=red!55!black,anchor=west] at ($(GA)!0.5!(GB)+(0.16,0)$) {Glauber};
    \node[t2 lab,anchor=west] at (2.85,1.5) {$\gamma^*$};
    \node[t2 lab,anchor=west] at (XA) {$X_A$};
    \node[t2 lab,anchor=west] at (XB) {$X_B$};
    \node[t2 lab,anchor=south] at (-3.25,1.1) {$P_A$};
    \node[t2 lab,anchor=north] at (-3.25,-1.35) {$P_B$};
    \draw[t2 time] (-3.5,-1.95) -- (3.15,-1.95) node[above left,t2 lab,text=gray!50!black] {时间};
    \node at (-0.2,-2.5) {(a)\ DY：旁观者之间的初态相互作用};
  \end{scope}
  % ======================= (b) DIS: final-state exchange between struck quark and remnant =======================
  \begin{scope}[xshift=7.6cm]
    \node[jet blob,circle,minimum size=0.66cm,inner sep=0pt] (A) at (-2.0,-0.95) {$A$};
    \node[hard blob,circle,minimum size=0.6cm,inner sep=0pt] (H) at (0.15,0.45) {$H$};
    \draw[t2 hadron] (-3.5,-0.95) -- (A.west);
    \draw[photon] (-1.35,1.75) -- (H.135);
    \draw[quark] (A.35) -- (H.215);
    \draw[quark] (H.20) -- (2.45,1.45);
    \draw[quark=0.42] (A.-15) -- (2.45,-1.25);
    \coordinate (Gq) at ($(H.20)!0.52!(2.45,1.45)$);
    \coordinate (Gr) at ($(A.-15)!0.735!(2.45,-1.25)$);
    \draw[glauber] (Gq) -- (Gr);
    \node[vertex] at (Gq) {}; \node[vertex] at (Gr) {};
    \path[mom={$q$},mom sep=0.28cm] (-1.35,1.75) -- (H.135);
    \node[t2 lab,anchor=west] at (2.45,1.45) {被击出的夸克};
    \node[t2 lab,anchor=west] at (2.45,-1.25) {靶残余};
    \node[t2 lab,text=red!55!black,anchor=west] at ($(Gq)!0.5!(Gr)+(0.17,0)$) {Glauber};
    \node[t2 lab,anchor=south] at (-3.15,-0.85) {$P$};
    \draw[t2 time] (-3.5,-1.95) -- (2.7,-1.95) node[above left,t2 lab,text=gray!50!black] {时间};
    \node at (-0.45,-2.5) {(b)\ DIS：被击出夸克与靶残余之间的末态相互作用};
  \end{scope}
\end{tikzpicture}
\caption{Glauber 胶子交换（暗红色）。时间自左向右，硬碰撞 $H$ 之前的交换为初态相互作用，之后为末态相互作用。
(a)~Drell--Yan：强子 $A$、$B$ 的旁观者（强子残余，进入末态 $X_A$、$X_B$）在硬碰撞之前交换 Glauber 胶子（$|k^+k^-|\ll k_T^2$）；来自 $B$ 的入射反夸克线与 $A$ 的旁观者线的交叉不是顶点。
在单个割图中，$A$ 一侧（旁观者线与活动部分子线）的 $k^-$ 极点与 $B$ 一侧的 $k^+$ 极点都分居实轴两侧，围道被钳在 Glauber 区（\cref{fig:glauber-contour}），只有对旁观者的末态切割求和（幺正性）之后才能把围道变形出 Glauber 区；
(b)~DIS：被击出的夸克与靶残余在硬碰撞之后交换 Glauber 胶子；对单举 DIS 它们由幺正性相消，而在 SIDIS 的 TMD 中它们组合成未来指向的 Wilson 线，
与 Drell--Yan 中由初态相互作用产生的过去指向 Wilson 线相对照（Sivers 函数因而在两过程间变号）。}
\label{fig:glauber-exchange}
\end{figure}
%%% END fig:glauber-exchange

钳夹的机制可以在单圈上完全显式地看到。下面的命题把\cref{prop:ir-glauber-trapping}（\tagR，$R_{\mathrm{PT}}$）的一般判据用到 DY 的旁观者上；同一个单圈例子见\cite[第~5.6.2~节]{Collins2011}与\cite[第~9.3~节，式~(166)]{CollinsSoperSterman1989}。

\begin{proposition}[单个割图中的 Glauber 钳夹 \tagR]\label{prop:dy-glauber-pinch}
（$R_{\mathrm{PT}}$：在重整化微扰论中对红外正规化的部分子态严格。）本命题是单圈的显式陈述，只涉及 Feynman 积分。设强子 $A$ 在顶点处分裂为活动部分子 $a$（动量 $l$，$l^+>0$）与旁观者 $s_A$（末态动量 $p'$，在壳 $p'^2=m_s^2$，$p'^+>0$），$l=P_A-p'$。一个动量为 $k$ 的胶子附着于 $s_A$，按标准路径 $k$ 经 $s_A$ 进入 $J_A$、经 $a$ 流向硬顶点，于是 $s_A$ 在附着点之前的一段携带 $p'-k$，$a$ 携带 $l+k$。设 $k$ 具有 Glauber 标度，$k^\pm\sim\lambda^2Q$，$|\mathbf k_T|\sim\lambda Q$，并且 $l^-,p'^-,|\mathbf l_T|^2/Q,\dots$ 具有 $J_A$ 的共线标度。则：
\begin{enumerate}[label=(\alph*),leftmargin=2.2em,itemsep=1pt,topsep=2pt]
\item $D_s\equiv(p'-k)^2-m_s^2+\ii0=-2p'^+(k^--c)+\ii0+\order{\lambda^4Q^2}$，$D_a\equiv(l+k)^2-m_a^2+\ii0=2l^+(k^--d)+\ii0+\order{\lambda^4Q^2}$，其中
$c=\frac{2\mathbf p'_T\cdot\mathbf k_T-\mathbf k_T^2}{2p'^+}$，$d=\frac{(\mathbf l_T+\mathbf k_T)^2+m_a^2-2l^+l^-}{2l^+}$ 为实数，均为 $\order{\lambda^2Q}$；$D_s$ 的零点 $k^-=c+\ii0$ 在上半平面，$D_a$ 的零点 $k^-=d-\ii0$ 在下半平面；
\item 在领头幂次（其余因子不依赖 $k^-$）下，$\displaystyle\int\frac{\dd k^-}{2\pi}\,\frac{\ii}{D_s}\,\frac{\ii}{D_a}=\frac{\ii}{4p'^+l^+(c-d)+\ii0}$，并且把积分围道整体移到 $\Imag k^-=\kappa$，只要 $|\kappa|$ 大于两个零点的虚部（$|\kappa|>0$ 的任何固定值即可），积分值都变为 $0$：围道被钳在 $|k^-|\lesssim|c-d|\sim\mathbf k_T^2/p'^+$ 处，不能离开 Glauber 区；
\item 若胶子的另一端附着于 $B$ 的旁观者 $s_B$（旁观者—旁观者交换），则对 $J_B$ 中的 $s_B$ 与活动部分子 $b$ 有同样的结论，其中 $+\leftrightarrow-$、$k\to-k$：$k^+$ 的两个零点分居实轴两侧，$k^+$ 也被钳住；两个分量同时被钳住，即 Glauber 钳夹。
\end{enumerate}
\uses{prop:ir-glauber-trapping,def:ir-glauber}
\end{proposition}

\begin{proof}
(a) 用 $p'^2=m_s^2$ 与 $v\cdot w=v^+w^-+v^-w^+-\mathbf v_T\cdot\mathbf w_T$，
$(p'-k)^2-m_s^2=-2p'^+k^--2p'^-k^++2\mathbf p'_T\cdot\mathbf k_T+2k^+k^--\mathbf k_T^2$。按 Glauber 与共线标度，$2p'^+k^-$、$2\mathbf p'_T\cdot\mathbf k_T$、$\mathbf k_T^2$ 都是 $\order{\lambda^2Q^2}$，而 $2p'^-k^+$、$2k^+k^-$ 是 $\order{\lambda^4Q^2}$；保留前三项即 $-2p'^+(k^--c)$。同样，$(l+k)^2-m_a^2=(2l^+l^--\mathbf l_T^2-m_a^2)+2l^+k^-+2l^-k^+-2\mathbf l_T\cdot\mathbf k_T+2k^+k^--\mathbf k_T^2=2l^+(k^--d)+\order{\lambda^4Q^2}$。零点位置由 $p'^+>0$、$l^+>0$ 与 $+\ii0$ 读出。
(b) 被积函数随 $|k^-|\to\infty$ 衰减为 $|k^-|^{-2}$，可在上半平面闭合围道，只包围 $k^-=c+\ii0$：
\[
\int\frac{\dd k^-}{2\pi}\frac{\ii}{-2p'^+(k^--c)+\ii0}\,\frac{\ii}{2l^+(k^--d)+\ii0}
=\ii\cdot\frac{\ii}{-2p'^+}\cdot\frac{\ii}{2l^+(c-d)+\ii0}=\frac{\ii}{4p'^+l^+(c-d)+\ii0}.
\]
若把实轴整体上移到 $\Imag k^-=\kappa>0$（$\kappa$ 大于零点 $c+\ii0$ 的虚部），新的围道上方没有奇点，在上半平面闭合给出 $0$；若下移到 $\Imag k^-=-\kappa$，在下半平面闭合同样给出 $0$。两种形变都改变了积分值，因此不存在离开实轴附近的形变：被积函数的全部贡献来自两零点之间 $|k^--c|\lesssim|c-d|$ 的区域，而 $|c-d|\sim\mathbf k_T^2/p'^+$。这正是\cref{prop:ir-glauber-trapping}中一条入射线与一条出射线的钳夹判据。
(c) 对 $J_B$，$s_B$ 在发射前携带 $p'_B+k$，$(p'_B+k)^2-m^2=2p'^-_B(k^+-c_B)+\order{\lambda^4Q^2}$，零点 $k^+=c_B-\ii0$；$b$ 携带 $l_b-k$，$(l_b-k)^2-m_b^2=-2l_b^-(k^+-d_B)+\order{\lambda^4Q^2}$，零点 $k^+=d_B+\ii0$。论证同 (b)。
\end{proof}

我们用数值积分核对了 (b)：取 $p'^+=0.6$、$l^+=0.4$、$c=0.013$、$d=-0.021$（任意单位），在实轴上以 $\ii\varepsilon$ 代替 $\ii0$ 求积分，与 $\ii/(4p'^+l^+(c-d))$ 的相对偏差在 $\varepsilon=10^{-4}$、$10^{-6}$ 时分别为 $6\times10^{-3}$、$6\times10^{-5}$，即 $\order{\varepsilon/|c-d|}$，随 $\varepsilon\to0$ 线性趋于零；而把围道移到 $\Imag k^-=\pm0.05$ 时积分在 $10^{-32}$ 的精度内为零。还应注意 (a) 的另一面：在 Glauber 区里，$D_a$ 中被程函近似（只保留 $2l^+k^-$）丢掉的 $-2\mathbf l_T\cdot\mathbf k_T-\mathbf k_T^2$ 与保留的项同阶，程函近似的误差是 $\order{1}$（\cref{def:ir-glauber}）。所以，在实轴上软近似失效，而在实轴附近又无法形变：这两件事合在一起，就是 DY 因子化证明必须专门处理的“Glauber 问题”。

% =====================================================================
\section{光锥序中间态的取向与 \texorpdfstring{$k^-$}{k-} 极点}\label{sec:dy-orientation}
% =====================================================================

\cref{prop:dy-glauber-pinch}是逐条线算出来的：旁观者的极点在上半平面，活动部分子的极点在下半平面。这种“逐线极点表”在多圈图中很快变得无法管理，而且它掩盖了真正起作用的结构。正确的表述在光锥序微扰论（light-cone ordered perturbation theory，LCOPT）中：对 $J_A$ 按光锥时间 $\tau=x^+$ 排序，软动量的负分量 $k^-$ 是注入的“光锥能量”；它出现在哪些能量分母中，由中间态位于硬顶点之前还是之后决定，与它附着在哪条 Feynman 线上无关。本节给出这一表述；附录中的精确形式与完整证明见\cref{sec:appO-orientation}。

先回顾所需的记号（\cref{def:appO-lcopt}）。固定 $J_A$ 中全部线的正分量与横分量，且每条线 $l^+\neq0$。在壳光锥能量 $E(l)\equiv(\mathbf l_T^2+m_l^2-\ii0)/(2l^+)$ 使 $l^2-m_l^2+\ii0=2l^+\big(l^--E(l)\big)$；对 $l^+>0$，$E(l)$ 位于下半平面。单线的时间表示
\begin{equation}\label{eq:dy-lc-fourier}
\int\frac{\dd l^-}{2\pi}\,\ee^{-\ii l^-\tau}\,\frac{\ii}{2l^+\big(l^--E(l)\big)}=\frac{\theta(l^+)\theta(\tau)+\theta(-l^+)\theta(-\tau)}{2|l^+|}\,\ee^{-\ii E(l)\tau}
\end{equation}
表明 $l^+>0$ 的线沿 $x^+$ 向前传播。把 $J_A$ 的全部负分量积掉，得到对与线的方向相容的顶点排序 $\pi$ 的求和，每个排序给出能量分母之积 $\prod_i\ii/D^\pi_i$，$D^\pi_i=Q^\pi_i-\sum_{l\in S^\pi_i}E(l)$，其中 $S^\pi_i$ 是第 $i$ 个中间态的线集，$Q^\pi_i$ 是到第 $i$ 个顶点为止注入的负分量之和（\cref{lem:appO-lcopt}，\tagR，$R_{\mathrm{PT}}$，在有限的 $\ii\eta_0$ 下成立，并可取 $\eta_0\to0^+$ 的分布极限）。在 DY 的割图中，振幅一侧硬顶点 $V_H$ 之前的中间态称为\emph{初态}，$V_H$ 之后、割线之前的称为\emph{末态}，共轭一侧为镜像（\cref{def:appO-orientation}）。软动量 $k_j$ 的\emph{标准流}是：$k_j$ 在其附着点进入 $J_A$，沿同一侧的路径流到 $V_H$（或 $V_H^\dagger$），经活动线离开。

\begin{lemma}[光锥序中间态的取向与 $k^-$ 极点 \tagR]\label{lem:dy-poles-orientation}
（$R_{\mathrm{PT}}$：在重整化微扰论中对红外正规化的部分子态严格。）本引理只涉及 Feynman 积分。考虑 $J_A$ 在有限 $\eta_0$ 下的光锥序表示（\cref{lem:appO-lcopt}），其活动动量在一个顶点 $V_H$ 处离开振幅部分，软动量取标准流。
\begin{enumerate}[label=(\alph*),leftmargin=2.2em,itemsep=1pt,topsep=2pt]
\item 取向是光锥序中间态的性质，而不是 Feynman 线的性质。振幅一侧的每个初态分母形如 $D=A_D+\sum_{j\in I(D)}k_j^-+\ii0$，每个末态分母形如 $D=A'_D-\sum_{j\in F(D)}k_j^-+\ii0$；这里 $A_D$、$A'_D$ 与全部软动量的负分量无关，$I(D)$ 是在该中间态之前注入的软动量的集合，$F(D)$ 是在其之后、割线之前注入的集合。除此之外，$k_j^-$ 只以多项式形式出现在顶点与接触因子中，不出现在 $1/(2l^+)$、在壳分子与割线的 $\delta$ 函数中。同一条旁观者线可以既属于初态又属于末态。
\item 因此，作为 $k_j^-$ 的函数，初态分母的零点位于 $\Imag k_j^-<0$，末态分母的零点位于 $\Imag k_j^->0$。耦合到软动量 $k$ 的线的 $k^-$ 极点全在同一侧，当且仅当 $k$ 所经过的全部中间态具有同一取向。共轭一侧镜像；对 $J_B$，交换 $+\leftrightarrow-$ 并以 $-k_j$ 为注入动量，$k_j^+$ 有同样的结论。
\end{enumerate}
\uses{lem:appO-lcopt,lem:coll-final-poles,prop:dy-glauber-pinch,def:appO-lcopt,def:appO-orientation,lem:appO-orientation}
\end{lemma}

\begin{proof}
这是\cref{lem:appO-orientation}（\tagR，$R_{\mathrm{PT}}$）的章内表述，完整证明（包括瞬时项与接触顶点）见\cref{sec:appO-orientation}；核心如下。在 $V_H$ 之前的第 $i$ 个中间态，注入的负分量是强子的 $P_A^-$ 加上此前注入的软能量，$D_i=P_A^-+\sum_{j\in I(D_i)}k_j^--\sum_{l\in S_i}E(l)$。在 $V_H$ 之后，振幅部分负分量的总守恒 $P_A^-+K_{\rm tot}-u^-=\mathcal E$（$u$ 为离开 $V_H$ 的活动动量，$\mathcal E=\sum_{\rm cut}E(p_c)$ 为割线上各线在壳能量之和）给出第 $i$ 个末态的注入能量 $\mathcal E-K_{>i}$（$K_{>i}$ 为该态之后、割线之前注入的软能量）。$E(l)$ 只依赖正分量与横分量，割线的 $\delta$ 函数只约束 $\mathcal E$，都不含软负分量；两种情形中 $\Imag D=\sum_{l\in S}\eta_0/(2l^+)>0$。这就是 (a)；(b) 中零点的位置由此直接读出。若 $k$ 只经过初态，$k^-$ 的围道可以整体推入上半平面（这只增加每个初态分母的虚部）；只经过末态的情形就是\cref{lem:coll-final-poles}（\tagR，$R_{\mathrm{PT}}$）；\cref{prop:dy-glauber-pinch}中的 $s_A$ 在一个光锥序项中属于初态（胶子在 $V_H$ 之前附着）、在另一项中属于末态，因而极点分居两侧。
\end{proof}

\cref{fig:dy-pole-orientation}画出了同一个费曼图的两个光锥序项：软胶子在 $V_H$ 之前附着于 $s_A$ 时只进入初态分母，在 $V_H$ 之后附着时只进入末态分母，二者之和使 $k^-$ 围道被钳住。还应注意 (b) 中“不同项之间的钳夹”的意义：Ward 恒等式只在对附着位置求和之后成立，所以全部光锥序项必须共用一条围道；位于向上与向下两种形变之间的极点（例如纵向胶子项在 $k^-=0$ 处的极点）的留数正是 Glauber 相位（\cref{prop:method-glauber-oneloop}，\tagR，$R_{\mathrm{PT}}$）。

%%% BEGIN fig:dy-pole-orientation
\begin{figure}[htbp]
\centering
\begin{tikzpicture}[fd,
  fc hadron/.style={draw=black!70,line width=1.5pt,line cap=butt},
  fc soft/.style={gluon,draw=green!45!black},
  fc state/.style={draw=gray!70!black,line width=0.6pt,dash pattern=on 0.6pt off 1.6pt},
  fc stateB/.style={draw=gray!62,line width=0.6pt,dash pattern=on 0.6pt off 1.6pt},
  fc ax/.style={lcaxis,gray!60!black},
  fc lab/.style={font=\small},
  fc note/.style={font=\footnotesize},
  fc tiny/.style={font=\scriptsize},
  fc pole/.pic={\draw[line width=0.9pt,line cap=round] (-0.075,-0.075)--(0.075,0.075) (-0.075,0.075)--(0.075,-0.075);},
  fc contour/.style={draw=blue!65!black,line width=1.05pt,postaction={decorate,decoration={markings,
      mark=at position #1 with {\arrow{Stealth[length=5pt,width=4.2pt]}}}}},
  fc contour/.default=0.5,
  fc band/.style={fill=red!9,draw=none},
  fc bandA/.style={fill=blue!5,draw=none},
  fc bandB/.style={fill=orange!8,draw=none}]
  % =====================================================================
  % (a) the soft gluon is attached before V_H (tau<0) on both spectators
  % =====================================================================
  \begin{scope}
    % jet subgraphs J_A (top, ordered in x^+) and J_B (bottom, ordered in x^-)
    \fill[fc bandA] (-2.75,0.27) rectangle (2.22,1.95);
    \fill[fc bandB] (-2.75,-1.95) rectangle (2.22,-0.27);
    \node[fc note,text=blue!45!black,anchor=north east] at (2.17,1.5) {$J_A$};
    \node[fc note,text=orange!45!black,anchor=south east] at (2.17,-1.5) {$J_B$};
    % light-cone time axes: x^+ for J_A (top), x^- for J_B (bottom)
    \draw[fc ax] (-3.55,2.75) -- (2.75,2.75) node[right,fc note,text=black] {$x^+$};
    \draw[fc ax] (-3.55,-2.75) -- (2.75,-2.75) node[right,fc note,text=black] {$x^-$};
    \draw[gray!60!black,line width=0.5pt] (0,2.68) -- (0,2.82) (0,-2.68) -- (0,-2.82);
    \node[fc tiny,above] at (0,2.78) {$0$};  \node[fc tiny,below] at (0,-2.78) {$0$};
    \node[fc tiny,above,text=gray!40!black] at (-1.75,2.78) {初态};
    \node[fc tiny,above,text=gray!40!black] at (1.2,2.78) {末态};
    \node[fc tiny,below,text=gray!40!black] at (-1.75,-2.78) {初态};
    \node[fc tiny,below,text=gray!40!black] at (1.2,-2.78) {末态};
    % intermediate states (vertical dotted lines)
    \draw[fc state] (-2.0,1.05) -- (-2.0,1.85);
    \draw[fc state] (-0.7,0.22) -- (-0.7,1.85);
    \draw[fc stateB] (-2.0,-1.05) -- (-2.0,-1.85);
    \draw[fc stateB] (-0.7,-0.22) -- (-0.7,-1.85);
    \node[fc tiny,text=gray!40!black] at (-2.0,2.08) {$D_1$};
    \node[fc tiny] at (-0.15,2.27) {$\dfrac{1}{D_2}=\dfrac{1}{k^-+A+\ii0}$};
    % soft gluon k from s_B to s_A (drawn first: it passes behind a and b)
    \draw[fc soft] (-1.5,-1.6) -- (-1.5,1.6);
    \path[mom'={$k$},mom pos=0.5,mom sep=0.27cm] (-1.5,-1.6) -- (-1.5,1.6);
    % hadrons, splitting vertices, spectators, active partons
    \draw[fc hadron] (-3.6,1.6) -- (-2.5,1.6);
    \draw[fc hadron] (-3.6,-1.6) -- (-2.5,-1.6);
    \draw[plain] (-2.5,1.6) -- (2.75,1.6);
    \draw[plain] (-2.5,-1.6) -- (2.75,-1.6);
    % a and b pass in front of the soft gluon (gap in the background colour of the band)
    \draw[blue!5,line width=3.8pt] (-1.71,1.095) -- (-1.29,0.825);
    \draw[orange!8,line width=3.8pt] (-1.71,-1.095) -- (-1.29,-0.825);
    \draw[quark=0.62] (-2.5,1.6) -- (-0.17,0.11);
    \draw[quark=0.38] (-0.17,-0.11) -- (-2.5,-1.6);
    \draw[photon] (0.2,0) -- (0.9,0);
    \node[hard blob,circle,minimum size=0.4cm,inner sep=0pt] (H) at (0,0) {};
    \foreach \p in {(-2.5,1.6),(-2.5,-1.6),(-1.5,1.6),(-1.5,-1.6)} \node[vertex] at \p {};
    \draw[cut] (2.3,-2.05) -- (2.3,2.05);
    % labels
    \node[fc lab,above] at (-3.25,1.62) {$P_A$};
    \node[fc lab,below] at (-3.25,-1.62) {$P_B$};
    \node[fc lab,right] at (2.75,1.6) {$s_A$};
    \node[fc lab,right] at (2.75,-1.6) {$s_B$};
    \node[fc lab] at (-2.25,0.98) {$a$};
    \node[fc lab] at (-2.25,-0.98) {$b$};
    \node[fc note,anchor=north west] at (0.12,-0.12) {$V_H$};
    \node[fc note,above] at (0.62,0.07) {$\gamma^*$};
    \node[fc lab] at (-0.4,-3.6) {(a)\ $k$在$V_H$之前附着：只进入初态};
  \end{scope}
  % =====================================================================
  % (b) the same Feynman graph, soft gluon attached after V_H (tau>0)
  % =====================================================================
  \begin{scope}[xshift=7.45cm]
    \fill[fc bandA] (-2.75,0.27) rectangle (2.22,1.95);
    \fill[fc bandB] (-2.75,-1.95) rectangle (2.22,-0.27);
    \node[fc note,text=blue!45!black,anchor=north east] at (2.17,1.5) {$J_A$};
    \node[fc note,text=orange!45!black,anchor=south east] at (2.17,-1.5) {$J_B$};
    \draw[fc ax] (-3.55,2.75) -- (2.75,2.75) node[right,fc note,text=black] {$x^+$};
    \draw[fc ax] (-3.55,-2.75) -- (2.75,-2.75) node[right,fc note,text=black] {$x^-$};
    \draw[gray!60!black,line width=0.5pt] (0,2.68) -- (0,2.82) (0,-2.68) -- (0,-2.82);
    \node[fc tiny,above] at (0,2.78) {$0$};  \node[fc tiny,below] at (0,-2.78) {$0$};
    \node[fc tiny,above,text=gray!40!black] at (-1.75,2.78) {初态};
    \node[fc tiny,above,text=gray!40!black] at (1.2,2.78) {末态};
    \node[fc tiny,below,text=gray!40!black] at (-1.75,-2.78) {初态};
    \node[fc tiny,below,text=gray!40!black] at (1.2,-2.78) {末态};
    \draw[fc state] (-1.25,0.58) -- (-1.25,1.85);
    \draw[fc state] (0.62,1.35) -- (0.62,1.85);
    \draw[fc stateB] (-1.25,-0.58) -- (-1.25,-1.85);
    \draw[fc stateB] (0.62,-1.35) -- (0.62,-1.85);
    \node[fc tiny,text=gray!40!black] at (-1.25,2.08) {$D_1$};
    \node[fc tiny] at (0.85,2.27) {$\dfrac{1}{D_2}=\dfrac{-1}{k^--A'-\ii0}$};
    \draw[fc soft] (1.25,-1.6) -- (1.25,1.6);
    \path[mom'={$k$},mom pos=0.5,mom sep=0.27cm] (1.25,-1.6) -- (1.25,1.6);
    \draw[fc hadron] (-3.6,1.6) -- (-2.5,1.6);
    \draw[fc hadron] (-3.6,-1.6) -- (-2.5,-1.6);
    \draw[plain] (-2.5,1.6) -- (2.75,1.6);
    \draw[plain] (-2.5,-1.6) -- (2.75,-1.6);
    \draw[quark=0.7] (-2.5,1.6) -- (-0.17,0.11);
    \draw[quark=0.3] (-0.17,-0.11) -- (-2.5,-1.6);
    \draw[photon] (0.2,0) -- (0.9,0);
    \node[hard blob,circle,minimum size=0.4cm,inner sep=0pt] (H) at (0,0) {};
    \foreach \p in {(-2.5,1.6),(-2.5,-1.6),(1.25,1.6),(1.25,-1.6)} \node[vertex] at \p {};
    \draw[cut] (2.3,-2.05) -- (2.3,2.05);
    \node[fc lab,above] at (-3.25,1.62) {$P_A$};
    \node[fc lab,below] at (-3.25,-1.62) {$P_B$};
    \node[fc lab,right] at (2.75,1.6) {$s_A$};
    \node[fc lab,right] at (2.75,-1.6) {$s_B$};
    \node[fc lab] at (-2.25,0.98) {$a$};
    \node[fc lab] at (-2.25,-0.98) {$b$};
    \node[fc note,anchor=north west] at (0.12,-0.12) {$V_H$};
    \node[fc note,above] at (0.55,0.07) {$\gamma^*$};
    \node[fc lab] at (-0.4,-3.6) {(b)\ $k$在$V_H$之后附着：只进入末态};
  \end{scope}
  % =====================================================================
  % (c), (d): complex k^- plane
  % =====================================================================
  \begin{scope}[shift={(1.2,-5.45)}]
    \fill[fc band] (-0.42,-0.85) rectangle (0.42,0.85);
    \draw[fc ax] (-1.75,0) -- (1.8,0) node[below,fc tiny,text=black,xshift=-3pt] {$\Real k^-$};
    \draw[fc ax] (0,-0.9) -- (0,1.0) node[left,fc tiny,text=black] {$\Imag k^-$};
    \draw[fc contour=0.85] (-1.65,0) -- (1.65,0);
    \pic at (-0.16,-0.26) {fc pole};  \pic at (0.16,0.26) {fc pole};
    \node[fc tiny,anchor=east] at (-0.24,-0.34) {(a)};
    \node[fc tiny,anchor=west] at (0.24,0.34) {(b)};
    \node[fc note,align=center] at (0,-1.52) {(c)\ (a)+(b)：极点分居两侧，\\$k^-$围道被钳住};
  \end{scope}
  \begin{scope}[shift={(5.7,-5.45)}]
    \fill[fc band] (-0.42,-0.85) rectangle (0.42,0.85);
    \draw[fc ax] (-1.75,0) -- (1.8,0) node[below,fc tiny,text=black,xshift=-3pt] {$\Real k^-$};
    \draw[fc ax] (0,-0.9) -- (0,1.0) node[left,fc tiny,text=black] {$\Imag k^-$};
    \draw[fc contour=0.5] (-1.65,0) -- (-1.2,0) .. controls (-0.95,0) and (-0.85,0.62) .. (0,0.62)
        .. controls (0.85,0.62) and (0.95,0) .. (1.2,0) -- (1.65,0);
    \pic at (-0.2,-0.22) {fc pole};  \pic at (0.12,-0.3) {fc pole};
    \node[fc note,align=center] at (0,-1.52) {(d)\ 只经过初态：极点同侧，\\$k^-$围道可移入上半平面};
  \end{scope}
\end{tikzpicture}
\caption{光锥序中间态的取向与$k^-$极点（\cref{lem:dy-poles-orientation}）。Drell--Yan 割图的振幅一侧（割线右侧的共轭振幅为镜像，未画出）按混合光锥序排列：$J_A$（上，浅蓝）按$x^+$、$J_B$（下，浅橙）按$x^-$排序，硬顶点$V_H$在两者中都位于$\tau=0$；
竖直点线标出中间态，$a$、$b$为活动部分子，$s_A$、$s_B$为旁观者，软胶子（绿色）动量$k$由$J_B$流入$J_A$（标准流）。
(a)、(b) 是同一费曼图的两个光锥序项：$k$在$V_H$之前附着于$s_A$时，$k^-$只出现在初态分母$D_2=k^-+A+\ii0$中，零点在下半平面；在$V_H$之后附着时，$k^-$只出现在末态分母$D_2=-(k^--A'-\ii0)$中，零点在上半平面
（$A$、$A'$为与$k^-$无关的实数；$J_B$中对$k^+$有同样的结论，$+\leftrightarrow-$）。取向因而是中间态而不是线的性质：(c)~$k$流经两种取向的中间态时，极点分居实轴两侧，$k^-$围道被钳在 Glauber 区（阴影）内，旁观者之间的交换在$J_A$、$J_B$中同时被钳住，即 Glauber 钳夹；
(d)~$k$只流经初态（例如只附着于$a$）时极点同侧，$k^-$围道可移入上半平面（离开 Glauber 区还须同时形变$k^+$，见\cref{fig:method-glauber-direction}）。}
\label{fig:dy-pole-orientation}
\end{figure}
%%% END fig:dy-pole-orientation

\begin{remark}[每个喷注一个光锥时间]\label{rem:dy-one-time-per-jet}
上面的构造对 $J_A$ 用光锥时间 $x^+$（能量为负分量），对 $J_B$ 用 $x^-$（能量为正分量），硬顶点在两者中都位于 $\tau=0$：这是 Aybat 与 Sterman 引入的“混合”光锥序\cite{AybatSterman2009}。选择的理由在于：$A$ 共线线的大分量是正分量，以 $x^+$ 为时间时，Glauber 与软动量只以小的能量 $k^-$ 进入其能量分母；在 Glauber 区内连接线的传播子 $\simeq-1/\mathbf k_T^2$ 与 $k^\pm$ 无关，因而在坐标空间中正比于 $\delta(x^+-y^+)\delta(x^--y^-)$，而 $J_A$ 不感知 $k^+$（附着点在 $x^-\simeq0$）、$J_B$ 不感知 $k^-$（附着点在 $y^+\simeq0$）——Glauber 交换发生在硬顶点的光锥时刻，即两个被收缩的强子彼此穿过的瞬间。后文所说的“测量与光锥演化相容”，精确的意思是：测量算符在割线两侧的末态之间作为恒等算符（或与演化对易的算符）出现；其一般判据见\cref{prop:break-measurement}（\tagR，$R_{\mathrm{cond}}$）。
\end{remark}

% =====================================================================
\section{Glauber 问题及其在单举 DY 中的解决}\label{sec:dy-glauber}
% =====================================================================

上两节表明：在单个割图中，旁观者之间交换的 Glauber 胶子在 $k^-$ 与 $k^+$ 上同时被钳住，软近似在实轴上又不成立。CSS 与 Bodwin 的全阶论证\cite{CollinsSoperSterman1985b,CollinsSoperSterman1988,Bodwin1985}与 Collins 的教科书处理\cite[第~14.4~节]{Collins2011}用下面的五步绕开这一障碍：(i) 对旁观者的末态割线求和，末态相互作用由幺正性相消；(ii) 余下的极点只来自初态线，位于同一侧；(iii) 把软动量沿 $e_G$ 同时形变出 Glauber 区；(iv) 在变形后的围道上做程函近似，用 Ward 恒等式把软胶子收成指向过去的 Wilson 线；(v) 对 $q_T$ 积分的截面，软因子由幺正性等于 $1$，共线因子就是 DIS 中的部分子分布。\cref{app:appO}把这五步写成一个逐步或者被证明、或者被列为编号假设的条件定理（\cref{thm:appO-main}，步骤 (O-1)--(O-8)）。本节给出各步的章内陈述，指明它们在附录中的精确形式；\cref{tab:dy-status}（\cref{sec:dy-theorem}）汇总每一步的地位与误差类型。

\paragraph{末态相互作用的幺正性相消。}

\begin{lemma}[旁观者末态相互作用的幺正性相消 \tagA]\label{lem:dy-fsi}
对旁观者系统的全部末态割线求和（截面对 $X$ 单举，测量只涉及轻子对），硬散射之后的相互作用相消：在每个光锥序的末态链中，把割线移过最后一个内部顶点只改变该顶点的权重；对割线位置求和后，末态相互作用的贡献等于去掉这些相互作用后的未割子图。（本引理只涉及部分子态。）
% errtype: E-S
\uses{prop:method-fsi,lem:appO-unitarity,asm:appO-lcopt,lem:appO-cut-position,lem:appJ-cut-hermiticity}
\end{lemma}

\begin{proofsketch}
这是\cref{prop:method-fsi}（全阶论证，\tagA）对 DY 旁观者的情形，精确形式是\cref{lem:appO-unitarity}（\tagR，$R_{\mathrm{cond}}$，前提\cref{asm:appO-lcopt}）。(1) 由\cref{asm:appO-lcopt}，$J_A$ 可写成光锥序形式；固定两侧初态的排序，把 $V_H$ 之后的顶点按闭时路径排列为 $v_1,\dots,v_n$，相邻中间态为 $s_0,\dots,s_n$，能量 $E_i=\sum_{l\in s_i}E(l)$。割线穿过 $s_j$ 时，振幅一侧给出 $\prod_{i<j}(\mathcal E-E_i+\ii0)^{-1}$，割线给出 $2\pi\delta(\mathcal E-E_j)$，共轭一侧给出 $\prod_{i>j}(\mathcal E-E_i-\ii0)^{-1}$，其中 $\mathcal E$ 是末态的总光锥能量，测量不约束它。(2) 记 $a_i=(\mathcal E-E_i+\ii\eta)^{-1}$、$b_i=(\mathcal E-E_i-\ii\eta)^{-1}$，则 $2\pi\delta(\mathcal E-E_j)\to2\eta/[(\mathcal E-E_j)^2+\eta^2]=\ii(a_j-b_j)$，令 $P_j=\prod_{i\le j}a_i\prod_{i>j}b_i$，得逐项伸缩求和的精确恒等式
\[
\sum_{j=0}^{n}\Big[\prod_{i<j}a_i\Big]\ii(a_j-b_j)\Big[\prod_{i>j}b_i\Big]=\sum_{j=0}^n\ii\,(P_j-P_{j-1})=\ii\prod_{i=0}^na_i-\ii\prod_{i=0}^nb_i ,
\]
% VF-GLA-5
对 $\mathcal E$ 积分时（$\int\dd\mathcal E/2\pi$），$n\ge1$ 时两个乘积都至少衰减为 $\mathcal E^{-2}$ 且全部极点各在同一半平面，故积分为零；$n=0$ 时积分为 $1$。这里用到内部顶点在割线两侧的权重相等，即割图规则的 Hermite 性，对鬼场与三胶子顶点由\cref{lem:appJ-cut-hermiticity}（\tagR，$R_{\mathrm{def}}$）保证。(3) 有软胶子附着的顶点在两侧权重不同，不能这样配对：割线求和后只剩下这些插入（\cref{lem:appO-cut-position}），其闭时路径形式为 $\langle P_A|\bar{\rm T}\{V_H^\dagger\cdots\}{\rm T}\{V_H\cdots\}|P_A\rangle$，$J_A$ 内部的末态相互作用不出现。文献：\cite{CollinsSoperSterman1985b,CollinsSoperSterman1988,Bodwin1985}与\cite[第~14.4~节]{Collins2011}；光锥微扰论的版本见\cite[第~4~节]{Diehl:2015bca}。完整论证见\cref{sec:appO-unitarity}。
\end{proofsketch}

\paragraph{地位与缺口}\tagA。割线位置恒等式本身是精确的代数恒等式；在\cref{asm:appO-lcopt}（\tagA：零模、瞬时圈与端点贡献在领头幂次不出现）之下，闭时路径形式是严格的（\cref{lem:appO-unitarity}）。全阶论证的等级来自前提：领头区域的标度陈述（\cref{prop:dy-regions}，E-S）与光锥序微扰论对全部图的适用性，后者在所有阶都没有作为定理证明，文献中只有零模与等价性的部分结果。唯一的物理输入是“测量不约束旁观者的末态”：若测量末态强子或加快度间隙，测量投影不与光锥演化对易，相消失效（\cref{prop:break-measurement}）。“单个割图本身就是幺正的”是错误的说法：相消只在对割线求和后发生。

\paragraph{单个割图中的 Wilson 线指向未来。}在讨论形变之前，先看一个常被误解的事实：在单个割图里，软胶子“看到”的强子 $A$ 的色荷沿什么方向延伸。

\begin{example}[单个割图给出未来指向的 Wilson 线 \tagR]\label{ex:dy-future-vs-past}
（$R_{\mathrm{def}}$：由算符定义直接推出。）在单个割图的振幅一侧，设强子 $A$ 的活动部分子 $a$（色荷 $Q_a$）与旁观者 $s$（色荷 $Q_s$）耦合到软胶子 $k$（$k$ 流入 $J_A$），在程函近似下只保留依赖于 $n_A\cdot k=k^-$ 的因子（规则见\cref{prop:prelim-wilson-rules}（\tagR，$R_{\mathrm{def}}$），去掉耦合 $\ii g\,n_A\cdot\tilde A$）。活动部分子在 $V_H$ 处被湮灭，只存在于 $\tau<0$，给出指向过去的因子 $\ii Q_a/(k^-+\ii0)$；旁观者从强子顶点一直存在到末态，其 $V_H$ 之前与之后的两段分别给出 $\ii Q_s/(k^-+\ii0)$ 与 $-\ii Q_s/(k^--\ii0)$。对色单态强子，作用在强子态上有 $Q_s=-Q_a$（阿贝尔情形为电荷，非阿贝尔情形为色荷算符），于是
\begin{equation}\label{eq:dy-future}
\frac{\ii Q_a}{k^-+\ii0}+\ii Q_s\Big[\frac1{k^-+\ii0}-\frac1{k^--\ii0}\Big]=\frac{\ii Q_a}{k^--\ii0}=Q_s\cdot\frac{-\ii}{k^--\ii0}.
\end{equation}
% VF-REF-11
右边是一条从硬顶点指向未来、携带旁观者色荷 $Q_s$ 的半无限程函线 $W[\infty n_A,0]$ 的单胶子因子，而不是指向过去的线 $W[0,-\infty n_A]$。
\uses{prop:prelim-wilson-rules}
\end{example}

推导只是代数：左边 $=(Q_a+Q_s)\frac{\ii}{k^-+\ii0}+Q_s\frac{-\ii}{k^--\ii0}$，第一项因 $Q_a+Q_s=0$ 而消失。由 $\frac1{k^-+\ii0}-\frac1{k^--\ii0}=-2\pi\ii\,\delta(k^-)$，旁观者的项等于 $Q_s\,2\pi\delta(k^-)$：一条完整的直程函线只吸收 $k^-=0$ 的胶子，这正是\cref{prop:method-glauber-oneloop}中的 Glauber（Coulomb）相位。物理图像很清楚：在硬碰撞之前，色单态强子对波长大于其尺度的软胶子是不可见的；硬碰撞取走活动部分子之后，留下的旁观者系统带有净色荷 $-Q_a$，像一个从 $\tau=0$ 出发、飞向未来的色源那样辐射（这一半经典图像见\cite[第~9.3~节]{CollinsSoperSterman1989}）。因子化公式中出现的却是指向过去的线：它们只有在对割图求和、把旁观者的耦合局域到 $V_H$ 的光锥时刻、并把余下的依赖改写为只含初态分母之后才出现（\cref{prop:dy-deformation}，精确形式\cref{prop:appO-past-wilson}，\tagA）。对共线部分子分布，Wilson 线的方向无关紧要（\cref{prop:dy-universality}）；对 TMD 分布，它决定订书钉形 Wilson 线的方向与 Sivers 函数的符号（\cref{thm:tmd-sivers-sign}，\tagR，$R_{\mathrm{def}}$）。“单个割图中的软 Wilson 线已经是指向过去的”是一个常见的误解。

\paragraph{同时形变与指向过去的 Wilson 线。}

\begin{proposition}[移除末态极点后的围道变形 \tagA]\label{prop:dy-deformation}
在对割图求和之后，耦合到软胶子的 $A$ 共线线的剩余 $k^-$ 极点可以全部移到同一侧（$B$ 共线线的 $k^+$ 极点亦然），软动量可以沿 $e_G=n_B-n_A$ 同时形变出 Glauber 区；在变形后的围道上 Grammer--Yennie 近似成立，软胶子耦合的程函因子取指向过去的处方。（本命题只涉及部分子态。）
% errtype: E-W
\uses{prop:method-glauber-deformation,lem:dy-fsi,lem:dy-poles-orientation,lem:appO-unitarity,lem:appO-equal-time,lem:appO-joint-localization,prop:appO-one-sided,prop:appO-target,prop:appO-eikonal,prop:appO-past-wilson,asm:appO-mixing,asm:appO-deformation,asm:appO-envelope,asm:appO-rewritten-cs,prop:appO-ss-cancel}
\end{proposition}

\begin{proofsketch}
按\cref{app:appO}的链条，共六步。(1) 割图求和后，$J_A$、$J_B$ 取闭时路径形式，各喷注内部的末态相互作用消失，只剩软插入（\cref{lem:dy-fsi}，\cref{lem:appO-unitarity}）。(2) 等光锥时刻局域化：在 Glauber 窗口内，外因子 $R(k^-)$ 在 $|k^-|<\Lambda$ 内解析而喷注的能量标度 $\lambda_J\ll\Lambda$，于是一个插入的 $k^-$ 积分等于 $R(0)$ 乘以插入在 $\tau=0^\pm$ 处的等时值，误差 $\order{\lambda_J/\Lambda}$（\cref{lem:appO-equal-time}）；旁观者插入的两个等时值相等，因为 $V_H$ 与作用于旁观者的算符对易，这一插入只传递 $\mathbf k_T$；同一部分子上的多个非对易插入由带次序权重的联合局域化处理（\cref{lem:appO-joint-localization}，\tagR，$R_{\mathrm{PT}}$）。(3) 单侧改写（\cref{prop:appO-one-sided}，\tagR，$R_{\mathrm{cond}}$）：把每个可局域插入的喷注因子换成其等时值乘以整函数载子 $\phi(k^-/w)/w$，$\phi(x)=2\sqrt\pi\,\ee^{-x^2}$；改写后每个未割的软 $k^-$ 在 $J_A$ 中或者只出现在初态分母中，或者只出现在载子中，与原被积函数弱等价（E-W），所需的纤维积分包络是\cref{asm:appO-envelope}；非重合的旁观者—旁观者交换在族求和后相消（\cref{prop:appO-ss-cancel}，\tagR，$R_{\mathrm{cond}}$），重合交换（非阿贝尔情形）需要\cref{asm:appO-mixing}(c)。(4) 同时形变：$k^-\to k^-+\ii\sigma\kappa\psi|\mathbf k_T|$、$k^+\to k^+-\ii\sigma\kappa\psi|\mathbf k_T|$，即沿 $e_G$ 形变；由
$(l+\ii t\sigma fe_G)^2-m^2=(l^2-m^2+2t^2f^2)+2\ii t\sigma f(l^+-l^-)$（\cref{prop:method-glauber-deformation}，\tagA），形变只使实部更类时，初态分母的虚部单调增大；对类 $\mathcal G_0$ 与分级构型，变形后 $|k^+||k^-|\ge c\,\mathbf k_T^2$，$c=\min(\kappa^2,\delta/2)$（\cref{prop:appO-target}，\tagR，$R_{\mathrm{PT}}$），一般情形是\cref{asm:appO-deformation}(a)。只推一个分量不够：$|k^+k^-|$ 仍是 $\order{\lambda^3}$。(5) 在变形后的围道上程函近似的相对误差为 $\order{\lambda/\kappa}$（逐点，E-P），而在实轴上的 Glauber 区内它是 $\order{1}$（\cref{prop:appO-eikonal}，\tagR，$R_{\mathrm{PT}}$）。(6) 在变形围道上 $\sigma\Imag(n_A\cdot K)\ge0$，唯一相容的处方是 $+\ii\sigma0$，即指向过去的 $W[0,-\infty n_A]$；对附着求和、用 Ward 恒等式即得 Wilson 线（\cref{prop:appO-past-wilson}，\tagA，逐项匹配归入\cref{asm:appO-rewritten-cs}）。Collins 的步骤是：区域分解、割线求和、“away from initial-state singularities”的形变（末态奇点被越过，但在割线求和后相消）、指向过去的 Grammer--Yennie 近似、Ward 恒等式\cite[第~14.4.1~节]{Collins2011}；混合光锥序与“钳夹变为初、末态形变方向之间的不匹配”这一看法来自\cite{AybatSterman2009}（这是一个纲要，不是证明）。完整论证见\cref{sec:appO-rewriting}与\cref{sec:appO-wilson}。
\end{proofsketch}

\paragraph{地位与缺口}\tagA。各步的地位（\cref{tab:dy-status}）：(1)(2)(4)(5) 在所述假设下有完整证明（误差类型依次为精确、E-F、精确、E-P）；最弱的是 (3) 与 (6)。它们依赖：\cref{asm:appO-mixing}（(a$'$)(b) 处理一般软子图，(c) 处理重合连接的完整图集相消）、\cref{asm:appO-deformation}（(b)(c)(d)：与区域集合代数的相容性、软角区、H-发射端点与重合的 SA/AS 插入；\emph{即使对阿贝尔图与类 $\mathcal G_0$ 也需要}）、\cref{asm:appO-envelope}、\cref{asm:appO-rewritten-cs}，全部为 \tagA。因此在截面层次上，DY 的 Glauber 相消没有任何一部分是无条件的。Collins 在“Issues remaining”中承认了这一步的状况\cite[第~14.4.8~节]{Collins2011}。要提升本命题的等级，需要证明这四个假设（特别是重合连接的完整图集相消与软角区），以及单侧改写所用的\cref{asm:appO-lcopt}、\cref{asm:appN-envelope}、\cref{asm:appN-separation}；在此之前，本命题的精确内容是附录中以它们为前提的条件定理（\cref{thm:appO-main}）。

第 (2) 步的最简单实例——一个旁观者模型中两个光锥序项之和恰好等于两个协变传播子之积、其 $k^-$ 积分只是一次等时刻的横动量转移——见\cref{ex:dy-spectator-lcopt}。

\begin{lemma}[软因子为 1 \tagA]\label{lem:dy-soft}
在\cref{prop:dy-deformation}的形变之后，软胶子退耦为指向过去的 Wilson 线 $W[0,-\infty n_A]$ 与 $W[0,-\infty n_B]$；对 $q_T$ 积分的单举 DY 截面，软因子
\[
S=\sum_{X_s}\langle0|\bar{\rm T}\big[\mathcal U^{(-)\dagger}\big]|X_s\rangle\langle X_s|{\rm T}\big[\mathcal U^{(-)}\big]|0\rangle,\qquad
\mathcal U^{(-)}=W[0,-\infty n_A]\,W[0,-\infty n_B],
\]
等于 $1$。（本引理只涉及部分子态。）
\uses{prop:dy-deformation,lem:method-soft-unitarity,prop:appO-past-wilson}
\end{lemma}

\begin{proofsketch}
(1) 由\cref{prop:dy-deformation}与\cref{prop:appO-past-wilson}，对软胶子在 $J_A$（$J_B$）上的全部附着求和，得到沿 $n_A$（$n_B$）、在 $V_H$ 处终止的指向过去的 Wilson 线，色荷在活动部分子的表示中。(2) 对 $q_T$ 积分的截面不约束软动量，对软末态 $X_s$ 求完备和，$S=\langle0|\bar{\rm T}[\mathcal U^{(-)\dagger}]\,{\rm T}[\mathcal U^{(-)}]|0\rangle$；对指向过去的类光线（$v^0>0$），编时乘积就是相互作用绘景的演化算符 ${\rm T}[\mathcal U^{(-)}]=U_I(0,-\infty)$，它是幺正的，故 $S=1$（\cref{lem:method-soft-unitarity}，\tagR，$R_{\mathrm{def}}$）；在维数正规化中，树图以上的每个图都是无标度的。(3) Collins 用时间反演不变性把过去与未来指向的软因子等同起来\cite[第~13.17~节与第~14.4.6~节]{Collins2011}；Bodwin--Brodsky--Lepage 的初态相位\cite{BodwinBrodskyLepage1981}“do not break factorization; rather they get moved into the parton densities”\cite[第~14.4.7~节]{Collins2011}。
\end{proofsketch}

\paragraph{地位与缺口}\tagA。等级继承自\cref{prop:dy-deformation}：一旦软耦合退耦为 Wilson 线，(2) 是严格的（\cref{lem:method-soft-unitarity}，$R_{\mathrm{def}}$）。“软因子为 1”只对共线单举截面成立：对 $q_T$ 微分的 DY，软因子 $\tilde S(\mathbf b_T)\neq1$，含快度发散，并通过定义吸收进 TMD 分布（\cref{def:tmd-soft}；\cref{thm:tmd-factorization}，\tagA）。逻辑次序上，幺正性只能在 Glauber 区处理之后使用：在实轴上的 Glauber 区内，软耦合还不是 Wilson 线。

\begin{proposition}[DY 与 DIS 的共线 PDF 是同一对象 \tagR]\label{prop:dy-universality}
（$R_{\mathrm{def}}$：由算符定义直接推出。）对 $w=(0^+,w^-,\mathbf 0_T)=w^-n_B$，
\begin{equation}\label{eq:dy-universality}
W[w,-\infty n_B]\,W[-\infty n_B,0]=W[w,0]=W[w,\infty n_B]\,W[\infty n_B,0] .
\end{equation}
因此由指向过去的 Wilson 线（DY）与指向未来的 Wilson 线（DIS）构成的共线夸克分布算符都等于\cref{def:pdf-quark}中的同一个算符 $\bar\psi(w)\gamma^+W[w,0]\psi(0)$，胶子分布同理。（本命题是算符恒等式。）
\uses{prop:prelim-wilson-covariance,def:pdf-quark}
\end{proposition}

\begin{proof}
三个点 $0$、$w$、$\pm\infty n_B$ 都在同一条类光直线 $\{\lambda n_B\}$ 上，半无限线理解为有限线段 $L\to\infty$ 的极限（带绝热正规子）。对同一直线上的点，路径排序的复合律与沿同一路径往返的线互逆（\cref{prop:prelim-wilson-covariance}，\tagR，$R_{\mathrm{def}}$）给出
\[
W[-Ln_B,0]=W[-Ln_B,w]\,W[w,0],\qquad W[w,-Ln_B]\,W[-Ln_B,w]=1
\]
（从 $0$ 经 $w$ 到 $-Ln_B$ 是同一直线上的路径，无论 $w^-$ 的正负；沿直线来回的部分由第二式消去）。两式相乘得 $W[w,-Ln_B]\,W[-Ln_B,0]=W[w,0]$；令 $L\to\infty$ 即得左边的等式。把 $-L$ 换成 $+L$ 得右边的等式。代入共线分布的定义即得结论。
\end{proof}

这一恒等式说明，Wilson 线方向的问题对共线因子化是无害的：DY 中的“指向过去”与 DIS 中的“指向未来”给出同一个部分子分布，这就是部分子分布普适性的算符根据（在单圈上的显式检验见\cref{rem:dy-dis-scheme}）。对 TMD 分布，算符的两个端点有横向间隔 $\mathbf b_T$，订书钉形的线不能收缩成直线段，方向成为物理的（\cref{thm:tmd-sivers-sign}）。在强子态之间取矩阵元时，还需要复合算符的存在性（\cref{asm:intro-composite}），这对 DIS 与 DY 是同一个假设。

% =====================================================================
\section{阿贝尔模型、光锥序表示与程函线间的双胶子交换}\label{sec:dy-model}
% =====================================================================

上一节的论证有许多步骤只在附录中以编号假设为前提而完成。为了看清每一步在做什么，本节先在一个可以算到底的模型中检验它们：阿贝尔规范理论、有质量的矢量玻色子（质量 $m$ 只作红外正规子）、程函（eikonal）旁观者。然后转向非阿贝尔情形中最常被误解的一点——两次 Glauber 交换的色结构——并精确计算两条程函线之间的双胶子交换。

\begin{proposition}[阿贝尔模型中旁观者 Glauber 相位的相消 \tagP]\label{prop:dy-abelian-model}
本命题是关于如下阿贝尔程函模型的陈述，在领头幂次下成立。模型：质量为 $m$ 的阿贝尔矢量玻色子，Feynman 规范传播子 $-\ii g_{\mu\nu}/(k^2-m^2+\ii0)$；带电标量部分子，$D_\mu=\partial_\mu-\ii g_jA_\mu$，顶点 $+\ii g_j(p+p')^\mu$；强子 $A$（$B$）在点顶点处分裂为活动部分子 $a$（$b$）与旁观者 $s_A$（$s_B$），$a$、$b$ 在硬顶点湮灭为一个质量为 $Q$ 的中性态；需要时把旁观者近似为沿 $n_A$（$n_B$）的无限长类光程函线 $x_j(\lambda)=\boldsymbol b_j+\lambda n_j$，荷为 $g_A$（$g_B$），横向间距 $\boldsymbol b=\boldsymbol b_A-\boldsymbol b_B$。
\begin{enumerate}[label=(G-\roman*),leftmargin=3em,itemsep=1pt,topsep=2pt]
\item 程函旁观者之间全部交换之和为 $\langle{\rm T}\,W_AW_B\rangle=\ee^{\ii\chi(\boldsymbol b)}$，$W_j=\exp\big(\ii g_j\int_{-\infty}^\infty\dd\lambda\,n_j\cdot A(x_j(\lambda))\big)$，
\begin{equation}\label{eq:dy-chi}
\chi(\boldsymbol b)=-\frac{g_Ag_B}{2\pi}\,K_0(m|\boldsymbol b|)\qquad(\text{实数}).
\end{equation}
% VF-DY-13
对末态动量为 $p'_A$、$p'_B$ 的在壳有限能量旁观者，单圈 Glauber 区的贡献为
$\mathcal M_1(p'_A,p'_B)=\int\frac{\dd^2\mathbf k_T}{(2\pi)^2}\,\frac{-\ii g_Ag_B}{\mathbf k_T^2+m^2}\,\mathcal M_0(p'_A-\mathbf k_T,p'_B+\mathbf k_T)$，即在与旁观者相对横动量共轭的冲击参数空间中 $\tilde{\mathcal M}_1=\ii\chi\,\tilde{\mathcal M}_0$。
\item 对任何关于旁观者相对横动量单举的截面，这一相位逐阶相消（$\order{g^{2n}}$ 的各割线位置之和 $\sum_{j=0}^n\frac{(\ii\chi)^j}{j!}\frac{(-\ii\chi)^{n-j}}{(n-j)!}=0$，$n\ge1$），也全阶相消（$|\ee^{\ii\chi}|^2=1$），包括领头幂次下的 $q_T$ 谱。
\item 旁观者—旁观者交换的 Glauber 区被钳住：$k^-$ 被 $s_A$ 与 $a$ 钳住，$k^+$ 被 $s_B$ 与 $b$ 钳住，不能用围道形变代替 (G-ii)。
\item $B$ 的活动部分子 $b$ 与 $A$ 的旁观者 $s_A$ 之间的交换（玻色子动量 $k$ 由 $s_A$ 流入 $b$，$b$ 携带 $l_b+k$，$l_b^->0$）在 $k^+$ 上不被钳住：
\begin{equation}\label{eq:dy-Ias}
I_{as}(k^-)\equiv\int\dd k^+\,\frac{1}{(2l_b^-k^++c_b+\ii0)(2k^+k^--K^2+\ii0)}=\frac{2\pi\ii\,\theta(-k^-)}{2l_b^-K^2+2k^-c_b},
\end{equation}
其中 $K^2=\mathbf k_T^2+m^2$，
它由 $k^+=K^2/(2k^-)$、即 $|k^+|\sim Q$ 处的玻色子极点给出；这一交换属于 $A$ 的共线子图，$b$ 可以换成指向过去的程函线 $\ii/(2l_b^-k^++\ii0)$，相对误差为 $\order{\lambda^2}$。
\end{enumerate}
本模型没有给出：非阿贝尔情形、全阶的区域分析、幂次修正，以及测量末态强子的可观测量。
\uses{prop:method-glauber-oneloop,lem:prelim-eikonal-identity,prop:dy-glauber-pinch}
\end{proposition}

\begin{proof}
(G-i) 对自由阿贝尔场，${\rm T}$ 乘积的 Gauss 积分给出 $\langle{\rm T}\,\ee^{\ii\int J\cdot A}\rangle=\ee^{-\frac12\int\!\!\int J^\mu J^\nu D_{\mu\nu}}$，其中 $J^\mu(x)=\sum_jg_j\int\dd\lambda\,n_j^\mu\delta^4(x-x_j(\lambda))$。自身项正比于 $n_j^2=0$，在 Feynman 规范中为零。交叉项为
\[
-g_Ag_B\,n_A\cdot n_B\int\frac{\dd^4k}{(2\pi)^4}\frac{-\ii\,\ee^{-\ii k\cdot(\boldsymbol b_A-\boldsymbol b_B)}}{k^2-m^2+\ii0}\,2\pi\delta(k^-)\,2\pi\delta(k^+)
=-\ii g_Ag_B\int\frac{\dd^2\mathbf k_T}{(2\pi)^2}\frac{\ee^{\ii\mathbf k_T\cdot\boldsymbol b}}{\mathbf k_T^2+m^2}=\ii\chi(\boldsymbol b),
\]
这里 $n_A\cdot n_B=1$，$\int\dd\lambda\,\ee^{-\ii\lambda n_A\cdot k}=2\pi\delta(k^-)$，最后用了 $\int\frac{\dd^2\mathbf k}{(2\pi)^2}\frac{\ee^{\ii\mathbf k\cdot\boldsymbol b}}{\mathbf k^2+m^2}=\frac{K_0(m|\boldsymbol b|)}{2\pi}$。无限长线迫使 $k^\pm=0$，交换是纯 Glauber 的。对有限能量的旁观者，取\cref{prop:dy-glauber-pinch}的路径：$k^-$ 的被积函数随 $|k^-|^{-2}$ 衰减，在上半平面闭合只包围 $s_A$ 的极点，$\int\frac{\dd k^-}{2\pi}\frac{\ii}{-2p'^+_A(k^--k_1)+\ii0}F(k^-)=F(k_1)/(2p'^+_A)$，即把内线的旁观者置于在壳，$F$ 中余下的因子成为动量平移后的树图振幅；$k^+$ 同理给出 $1/(2p'^-_B)$。两个顶点在领头幂次为 $\ii g_A\,2p'^+_A$ 与 $\ii g_B\,2p'^-_B$，Glauber 区的玻色子传播子为 $-\ii/(-\mathbf k_T^2-m^2)$，合起来
$\frac{1}{2p'^+_A}\frac{1}{2p'^-_B}(\ii g_A2p'^+_A)(\ii g_B2p'^-_B)\frac{-\ii}{-\mathbf k_T^2-m^2}=\frac{-\ii g_Ag_B}{\mathbf k_T^2+m^2}$，即所述的 $\mathcal M_1$。交换只把旁观者的相对横动量 $\boldsymbol p=\frac12(\mathbf p'_A-\mathbf p'_B)$ 平移 $\mathbf k_T$；以 $\mathcal M(\boldsymbol p)=\int\dd^2\boldsymbol b\,\ee^{-\ii\boldsymbol p\cdot\boldsymbol b}\tilde{\mathcal M}(\boldsymbol b)$ 表示，卷积变为乘积，$\tilde{\mathcal M}_1(\boldsymbol b)=\tilde{\mathcal M}_0(\boldsymbol b)\int\frac{\dd^2\mathbf k}{(2\pi)^2}\frac{-\ii g_Ag_B\,\ee^{\ii\mathbf k\cdot\boldsymbol b}}{\mathbf k^2+m^2}=\ii\chi(\boldsymbol b)\tilde{\mathcal M}_0(\boldsymbol b)$。$n$ 次交换的梯形图与全部交叉梯形图之和为 $(\ii\chi)^n/n!$：对程函线这是 Gauss 公式的展开，对有限能量的旁观者则由程函恒等式 $\sum_{\pi}\prod_j(w_{\pi(1)}+\dots+w_{\pi(j)})^{-1}=\prod_jw_j^{-1}$（\cref{lem:prelim-eikonal-identity}）保证。

(G-ii) 交换不改变旁观者的总横动量 $\boldsymbol P=\mathbf p'_A+\mathbf p'_B=-\mathbf q_T$，$q^\pm$ 的改变是 $\order{\lambda^2}$。对不依赖于 $\boldsymbol p$ 的可观测量，用 Parseval 恒等式，
$\int\dd^2\boldsymbol p\,|\mathcal M|^2=(2\pi)^2\int\dd^2\boldsymbol b\,|\tilde{\mathcal M}_0|^2\,|\ee^{\ii\chi}|^2=\int\dd^2\boldsymbol p\,|\mathcal M_0|^2$。逐阶地，$\order{g^{2n}}$ 中 $j$ 次交换在割线左侧（振幅）、$n-j$ 次在右侧（共轭振幅），其和为 $\sum_j\frac{(\ii\chi)^j}{j!}\frac{(-\ii\chi)^{n-j}}{(n-j)!}=\frac{(\ii\chi-\ii\chi)^n}{n!}=0$（二项式定理）。

(G-iii) 这是\cref{prop:dy-glauber-pinch}(b)(c) 的特例：$s_A$ 与 $a$ 的 $k^-$ 零点分居两侧，$s_B$ 与 $b$ 的 $k^+$ 零点分居两侧，围道的任何离开实轴附近的形变都改变积分值。

(G-iv) 按所述路径，$b$ 的分母为 $(l_b+k)^2-m_b^2+\ii0=2l_b^-k^++c_b+\ii0+\order{\lambda^4Q^2}$，$c_b=l_b^2-m_b^2-2\mathbf l_{bT}\cdot\mathbf k_T-\mathbf k_T^2$ 为实数、$\order{\lambda^2Q^2}$，零点 $k^+=-c_b/(2l_b^-)-\ii0$ 在下半平面；$s_A$ 与 $a$ 只在次领头幂次依赖 $k^+$。玻色子极点 $k^+=(K^2-\ii0)/(2k^-)$ 对 $k^->0$ 在下半平面，对 $k^-<0$ 在上半平面。被积函数衰减为 $|k^+|^{-2}$，在上半平面闭合：$k^->0$ 时上半平面无极点，积分为零；$k^-<0$ 时只包围玻色子极点，得 $2\pi\ii/[(2l_b^-K^2/(2k^-)+c_b)\,2k^-]=2\pi\ii/(2l_b^-K^2+2k^-c_b)$。该极点在 $|k^+|=K^2/(2|k^-|)\sim\lambda^2Q^2/(\lambda^2Q)=Q$ 处：动量 $k$ 是 $A$ 共线的，不是 Glauber 的。把 $b$ 的分母换成 $2l_b^-k^++\ii0$，结果变为 $2\pi\ii\theta(-k^-)/(2l_b^-K^2)$，相对改变为 $k^-c_b/(l_b^-K^2)=\order{\lambda^2}$；$\ii/(2l_b^-k^++\ii0)=\frac{1}{2l_b^-}\,\ii/(n_B\cdot k+\ii0)$ 正是指向过去的 Wilson 线 $W[0,-\infty n_B]$ 的程函因子（\cref{prop:method-glauber-oneloop}的规则），它属于 $f_{q/A}$ 的 Wilson 线。
\end{proof}

\paragraph{地位与缺口}\tagP。作为 QCD 陈述，“旁观者之间的 Glauber 相位在单举截面中相消”只在这个模型中被证明（\cref{lem:dy-fsi}、\cref{prop:dy-deformation}给出一般情形的全阶论证）。模型与 QCD 的差别正是论证的难点：非阿贝尔情形中相位成为不与硬过程的颜色流、也彼此不对易的色矩阵；有限能量的复合旁观者不是程函线；区域之间有重叠与减除。阿贝尔理论中、只含非重合连接的阶上，对任意的 $J_A$、$J_B$ 与任意多个旁观者，族求和后的相消在\cref{asm:appO-lcopt}之下成立（\cref{prop:appO-ss-cancel}，本模型的推广）；非阿贝尔的重合连接需要\cref{asm:appO-mixing}(c)。这些结果是已知结果的推导，见\cite[第~14.3~节，式~(14.11)--(14.16)，第~14.3.3--14.3.4~节]{Collins2011}；一个显式的阿贝尔相位见\cite[式~(14)]{AybatSterman2009}，SCET 中旁观者—旁观者 Glauber 相位在对 $\Delta p_\perp$ 积分后相消见\cite[式~(11.12)--(11.13)]{RothsteinStewart2016}。

\cref{fig:dy-abelian-model}画出了 (G-i)(G-ii)：$n$ 次交换之和是冲击参数空间中的相位，两种割线位置在对相对横动量单举的截面中相消。

%%% BEGIN fig:dy-abelian-model
\begin{figure}[htbp]
\centering
\begin{tikzpicture}[fd,
  fc lab/.style={font=\small},
  fc note/.style={font=\footnotesize},
  fc amp/.style={hadron blob,ellipse,minimum width=0.95cm,minimum height=1.9cm,inner sep=0pt,font=\footnotesize}]
  % =====================================================================
  % (a) n exchanges between the two eikonal spectator lines
  % =====================================================================
  \begin{scope}
    \node[fc lab] at (-3.15,0) {$\displaystyle\sum_{n\ge0}$};
    \draw[wilson=0.1] (-2.5,0.75) -- (1.9,0.75);
    \draw[wilson=0.1] (-2.5,-0.75) -- (1.9,-0.75);
    \foreach \x in {-1.35,-0.55,1.05} {
      \draw[glauber] (\x,-0.75) -- (\x,0.75);
      \node[vertex] at (\x,0.75) {};  \node[vertex] at (\x,-0.75) {}; }
    \node[fc lab] at (0.25,0) {$\cdots$};
    \node[fc note,anchor=south west] at (-2.5,0.86) {$s_A$：沿$n_A$，荷$g_A$};
    \node[fc note,anchor=north west] at (-2.5,-0.86) {$s_B$：沿$n_B$，荷$g_B$};
    \draw[gray!60!black,line width=0.5pt,{Stealth[length=3.5pt,width=3pt]}-{Stealth[length=3.5pt,width=3pt]}] (2.2,0.7) -- (2.2,-0.7);
    \node[fc lab,anchor=west] at (2.25,0) {$\boldsymbol b$};
    \draw[line width=0.5pt] (-1.45,-1.43) -- (-1.45,-1.55) -- (1.15,-1.55) -- (1.15,-1.43) (-0.15,-1.55) -- (-0.15,-1.65);
    \node[fc note] at (-0.15,-1.93) {$n$次交换（含全部交叉排列）};
    \node[fc lab,anchor=west] at (2.85,0.3) {$\displaystyle=\sum_{n\ge0}\frac{\big(\ii\chi(\boldsymbol b)\big)^n}{n!}=\ee^{\ii\chi(\boldsymbol b)}$};
    \node[fc lab,anchor=west] at (2.85,-0.75) {$\displaystyle\chi(\boldsymbol b)=-\frac{g_Ag_B}{2\pi}\,K_0(m|\boldsymbol b|)$};
    \node[fc lab] at (2.3,-2.6) {(a)\ 程函旁观者之间$n$次有质量阿贝尔玻色子交换之和为相位};
  \end{scope}
  % =====================================================================
  % (b) the two cut placements at O(g^2)
  % =====================================================================
  \begin{scope}[yshift=-4.95cm]
    \begin{scope}[xshift=-1.2cm]
      \draw[wilson=0.76] (-1.6,0.5) -- (1.6,0.5);
      \draw[wilson=0.76] (-1.6,-0.5) -- (1.6,-0.5);
      \node[fc amp] at (-1.6,0) {$\mathcal M^{(0)}$};
      \node[fc amp] at (1.6,0) {$\mathcal M^{(0)*}$};
      \draw[glauber] (-0.62,-0.5) -- (-0.62,0.5);
      \node[vertex] at (-0.62,0.5) {};  \node[vertex] at (-0.62,-0.5) {};
      \draw[cut] (0,-1.15) -- (0,1.15);
      \node[fc note] at (-0.62,0.78) {$s_A$};  \node[fc note] at (-0.62,-0.8) {$s_B$};
      \node[fc lab] at (0,-1.55) {$+\ii\chi(\boldsymbol b)\,|\tilde{\mathcal M}^{(0)}(\boldsymbol b)|^2$};
    \end{scope}
    \node[fc lab] at (1.55,0) {$+$};
    \begin{scope}[xshift=4.3cm]
      \draw[wilson=0.3] (-1.6,0.5) -- (1.6,0.5);
      \draw[wilson=0.3] (-1.6,-0.5) -- (1.6,-0.5);
      \node[fc amp] at (-1.6,0) {$\mathcal M^{(0)}$};
      \node[fc amp] at (1.6,0) {$\mathcal M^{(0)*}$};
      \draw[glauber] (0.62,-0.5) -- (0.62,0.5);
      \node[vertex] at (0.62,0.5) {};  \node[vertex] at (0.62,-0.5) {};
      \draw[cut] (0,-1.15) -- (0,1.15);
      \node[fc note] at (0.62,0.78) {$s_A$};  \node[fc note] at (0.62,-0.8) {$s_B$};
      \node[fc lab] at (0,-1.55) {$-\ii\chi(\boldsymbol b)\,|\tilde{\mathcal M}^{(0)}(\boldsymbol b)|^2$};
    \end{scope}
    \node[fc lab,anchor=west] at (6.55,0) {$=0$};
    \node[fc lab] at (2.3,-2.3) {$\order{g^{2n}}$：$\displaystyle\sum_{j=0}^{n}\frac{(\ii\chi)^j}{j!}\,\frac{(-\ii\chi)^{n-j}}{(n-j)!}=0\ (n\ge1)$；全阶：$|\ee^{\ii\chi(\boldsymbol b)}|^2=1$};
    \node[fc lab] at (2.3,-3.2) {(b)\ 割线位置之和（对旁观者相对横动量$\boldsymbol\Delta$积分，即在$\boldsymbol b$空间中）};
  \end{scope}
\end{tikzpicture}
\caption{阿贝尔模型中旁观者之间的 Glauber 交换（\cref{prop:dy-abelian-model}）。两个旁观者$s_A$、$s_B$近似为沿$n_A$、$n_B$（方向相反）的无限长程函线，荷为$g_A$、$g_B$，横向间距为$\boldsymbol b$；暗红色螺线为质量为$m$的阿贝尔矢量玻色子（Feynman 规范）；时间自左向右，双线上的箭头为荷的流向。
(a)~对$n$次交换（含全部交叉排列）求和给出冲量参数空间的相位：$n$次项为$(\ii\chi)^n/n!$，总和为$\ee^{\ii\chi(\boldsymbol b)}$，$\chi(\boldsymbol b)=-\frac{g_Ag_B}{2\pi}K_0(m|\boldsymbol b|)$为实数（$m\to0$时红外发散，但始终是相位）。
(b)~$\order{g^2}$的两个割图：交换位于割线左侧（振幅）给出$+\ii\chi$，位于右侧（共轭振幅）给出$-\ii\chi$；对旁观者相对横动量$\boldsymbol\Delta$单举的截面中（Parseval 恒等式化为$\boldsymbol b$空间）二者逐阶相消，全阶时$|\ee^{\ii\chi}|^2=1$，但对固定的$\boldsymbol\Delta$逐点并不相消（VF-DY-13，VF-GLA-3）。}
\label{fig:dy-abelian-model}
\end{figure}
%%% END fig:dy-abelian-model

\begin{example}[旁观者模型的光锥序表示]\label{ex:dy-spectator-lcopt}
把\cref{lem:dy-poles-orientation}用到\cref{prop:dy-glauber-pinch}的旁观者模型。去掉与 $k^-$ 无关的因子 $1/(2p'^+\,2l^+)$，$J_A$ 对一个附着于 $s_A$ 的软插入的依赖为
\begin{equation}\label{eq:dy-JA-two-orderings}
J_A(k^-)=\frac{\ii}{k^--d+\ii0}\cdot\frac{\ii}{c-k^-+\ii0}
=\frac{\ii}{c-d}\Big[\underbrace{\frac{\ii}{k^--d+\ii0}}_{\text{插入在 }V_H\text{ 之前}}+\underbrace{\frac{\ii}{c-k^-+\ii0}}_{\text{插入在 }V_H\text{ 之后}}\Big],
\end{equation}
第二个等号是部分分式 $\frac{1}{(x-d)(c-x)}=\frac{1}{c-d}\big[\frac{1}{x-d}+\frac{1}{c-x}\big]$。左边是两个协变传播子之积，右边是两个光锥序项之和：$\ii/(c-d)$ 是插入之前初态 $(a,s_A)$ 的能量分母 $D_1=c-d$，第一项是插入在 $V_H$ 之前的初态分母（零点在下半平面），第二项是插入在 $V_H$ 之后的末态分母（零点在上半平面）。记 $J^{\rm in}$、$J^{\rm fi}$ 为这两项，二者在 $|k^-|\gg|c|,|d|$ 时分别趋于 $\ii\mathcal J_-/k^-$ 与 $-\ii\mathcal J_+/k^-$，且 $\mathcal J_-=\mathcal J_+=\ii/(c-d)$：旁观者插入在 $V_H$ 两侧的等时值相等。设外因子 $R(k^-)=r/(k^--z)$ 只有一个远处的极点（连接线的 $k^-$ 极点，$|z|\gg|c|,|d|$）。若 $\Imag z<0$，$J^{\rm in}R$ 在上半平面解析且衰减为 $|k^-|^{-2}$，其积分为零；$J^{\rm fi}R$ 在下半平面闭合只包围 $z$，得 $\ii r/[(c-d)(c-z)]=R(0)\,\frac{\ii}{c-d}\,[1+\order{c/z}]$。$\Imag z>0$ 时角色互换。所以
$\int\frac{\dd k^-}{2\pi}J_A(k^-)R(k^-)=R(0)\,\frac{\ii}{c-d}\,\big[1+\order{\max(|c|,|d|)/|z|}\big]$：
没有额外的能量分母，交换发生在 $\tau=0$，只转移横动量——这就是等光锥时刻局域化（\cref{prop:dy-deformation}第 (2) 步）的最简单形式。单个光锥序项则依赖于正规化：若把连接线换成严格的 Glauber 传播子 $-1/\mathbf k_T^2$ 并在 $|k^-|<K$ 处对称截断，每一项各给出权重 $\frac12$；若保留完整传播子（$k^+>0$ 时极点在下半平面），末态项给出全部权重、初态项给出零，即权重 $\theta(\pm k^+)$。两项之和与正规化无关。数值上，取 $d=-0.011$、$c=0.017$ 与三个远处极点 $z$，积分与 $R(0)\ii/(c-d)$ 的相对偏差为 $0.7\%$--$1.9\%$，与 $\max(|c|,|d|)/|z|$ 相符；对称截断 $K=50$ 时两项的权重都等于 $0.5000$。与\cref{prop:dy-abelian-model}的核 $\chi(\boldsymbol b)$ 合起来，每次交换在冲击参数空间中贡献一个因子 $\ii\chi(\boldsymbol b)$，对相对横动量积分后相消。
\uses{lem:appO-lcopt,lem:dy-poles-orientation,prop:dy-abelian-model}
\end{example}

\begin{example}[阿贝尔 Glauber 模型的数值演示]\label{ex:dy-glauber-model}
\cref{prop:dy-abelian-model}(G-ii) 中的相消只发生在对旁观者相对横动量 $\boldsymbol\Delta$ 积分之后，逐点并不成立。取一个复的玩具振幅（$\boldsymbol\Delta$ 以其宽度为单位）
$\mathcal M_0(\boldsymbol\Delta)=\ee^{-[(\Delta_x-0.4)^2+\Delta_y^2]/2}\,(1+0.9\ii\Delta_x+0.3\Delta_y^2)\,\ee^{0.7\ii\Delta_y}$，
$g_Ag_B=2$、$m=0.3$，在 $512\times512$ 的网格（边长 $40$）上用快速 Fourier 变换实现 $\boldsymbol\Delta\leftrightarrow\boldsymbol b$，于是 $\tilde{\mathcal M}_1=\ii\chi\tilde{\mathcal M}_0$、$\tilde{\mathcal M}_2=\frac12(\ii\chi)^2\tilde{\mathcal M}_0$。结果：
(i) $\order{g^2}$ 的两个割图之和 $2\Real[\mathcal M_1\mathcal M_0^*](\boldsymbol\Delta)$ 逐点不为零，其绝对值的积分是两个割图各自绝对值积分之和的 $21\%$，而其积分本身为 $4\times10^{-17}$（以同一量为单位）；
(ii) $\order{g^4}$ 的割线位置之和 $|\mathcal M_1|^2+2\Real[\mathcal M_0^*\mathcal M_2]$ 的积分为 $\int|\mathcal M_1|^2$ 的 $1.5\times10^{-16}$，逐点则为其 $18\%$；
(iii) 全阶，$\int|\ee^{\ii\chi}\otimes\mathcal M_0|^2/\int|\mathcal M_0|^2=1.000000000000000$；
(iv) 然而相对横动量的分布确实被改变：$\langle\boldsymbol\Delta^2\rangle$ 由 $1.738$ 变为 $1.874$。
这说明相消不是平凡的：它要求对旁观者的末态求和；任何约束旁观者相对横动量的测量都会看到 Glauber 相位的效应。
\uses{prop:dy-abelian-model}
\end{example}

\paragraph{非阿贝尔情形：两次交换的色结构。}在非阿贝尔理论中，一次旁观者—旁观者交换的振幅是 $-\ii\chi X$，$X=\mathbf T_A\cdot\mathbf T_B$ 是 Hermite 的色算符。两次交换时，非交叉图（梯形）与交叉图的色因子分别是 $X^2$ 与 $X^2+\frac{C_A}2X$，二者不同，阿贝尔情形中把两者在被积函数层次合并为 $\delta(k^+)\delta(k^-)$ 的做法不再可用\cite[第~5.2.1~节]{RothsteinStewart2016}。若两个喷注中两次插入的次序以权重 $p_{\rm same}$、$p_{\rm opp}$（$p_{\rm same}+p_{\rm opp}=1$）相同或相反，二阶振幅为 $\mathcal A_2=\frac{(-\ii\chi)^2}2\big[X^2+p_{\rm opp}\frac{C_A}2X\big]$，幺正性的残差为 $\mathcal A_2+\mathcal A_2^\dagger+\mathcal A_1^\dagger\mathcal A_1=-\frac12C_A\chi^2(\Real p_{\rm opp})X$（\cref{lem:appO-glauber-operator}，\tagR，$R_{\mathrm{def}}$）。如果像“严格 Glauber 窗口”那样把连接线换成 $-1/\mathbf k_T^2$ 并对称截断，每个喷注的次序各以 $\frac12$ 的权重独立出现，$p_{\rm opp}=\frac12$，残差为 $-\frac14C_A\chi^2X$，于是一个色单态 $q\bar q$ 对（$X=-C_F$）的“概率”是 $1+\frac14C_AC_F\chi^2>1$。下面的命题表明，这个残差是截断的假象。

\begin{proposition}[两条程函线之间的双胶子交换：$O(g^4)$ 幺正性与 Regge 相位 \tagR]\label{prop:dy-eikonal-twogluon}
（$R_{\mathrm{PT}}$：在重整化微扰论中对红外正规化的部分子态严格。）本命题是 $\order{g^4}$ 的显式陈述，只涉及程函线模型。两条无限长的类时 Wilson 线，速度 $v_A$、$v_B$（$v^2=1$，$v_A\cdot v_B=\cosh Y$，快度差 $Y>0$），横向间距 $\boldsymbol b$，色生成元 $\mathbf T_A$、$\mathbf T_B$（任意表示），交换两个质量为 $m$（红外正规子）的胶子，使用完整的 Feynman 规范传播子（保留 $k^+k^-$），横动量固定；记 $a=\mathbf k_T^2+m^2$、$b=(\mathbf K-\mathbf k)_T^2+m^2$（$\mathbf K$ 为总横动量转移），$x_A=v_A\cdot k$、$x_B=v_B\cdot k$，$\mathcal Q\equiv k_\parallel^2=\big(2\cosh Y\,x_Ax_B-x_A^2-x_B^2\big)/\sinh^2Y$，$g(\mathcal Q)=[(\mathcal Q-a+\ii0)(\mathcal Q-b+\ii0)]^{-1}$。
\begin{enumerate}[label=(\alph*),leftmargin=2.2em,itemsep=1pt,topsep=2pt]
\item 纵向积分
\begin{equation}\label{eq:dy-Jbox}
\begin{aligned}
J_{\rm box}&\equiv\int\dd x_A\dd x_B\,\frac{\ii}{x_A+\ii0}\,\frac{\ii}{-x_B+\ii0}\,g=\frac{2\pi}{ab}(\pi+\ii Y),\\
J_{\rm cross}&\equiv\int\dd x_A\dd x_B\,\frac{\ii}{x_A+\ii0}\,\frac{\ii}{x_B+\ii0}\,g=-\frac{2\pi\ii Y}{ab};
\end{aligned}
\end{equation}
以 $X=\mathbf T_A\cdot\mathbf T_B$，非交叉图与交叉图的色因子为 $c_{\rm box}=X^2$、$c_{\rm cross}=X^2+\frac{C_A}2X$，于是
$J_{\rm box}c_{\rm box}+J_{\rm cross}c_{\rm cross}=\frac{2\pi}{ab}\big[\pi X^2-\ii Y\frac{C_A}2X\big]$。
% VF-REF-9
\item 以单胶子交换 $\mathcal A_1=-\ii\chi X$，$\chi(\boldsymbol b)=g^2\coth Y\int\frac{\dd^2\mathbf k_T}{(2\pi)^2}\frac{\ee^{\ii\mathbf k_T\cdot\boldsymbol b}}{\mathbf k_T^2+m^2}=\frac{g^2\coth Y}{2\pi}K_0(m|\boldsymbol b|)$，在横向积分之后双胶子交换为
$\mathcal A_2=\frac12(-\ii\chi X)^2+(-\ii\chi X)\big(-\frac{C_AY}{4\pi}\chi\big)$：第一项是指数化项，第二项是 $\ii\times$实数$\times$Hermite 色矩阵，即正比于 $C_AY$ 的单圈 Regge 化相位。
\item 在\cref{lem:appO-glauber-operator}的参数化 $\mathcal A_2=\frac{(-\ii\chi)^2}{2}\big[X^2+p_{\rm opp}\frac{C_A}2X\big]$ 中，$p_{\rm opp}=-\ii Y/\pi$，$\Real p_{\rm opp}=0$，因而 $\mathcal A_2+\mathcal A_2^\dagger+\mathcal A_1^\dagger\mathcal A_1=0$：双胶子交换图给出 $\mathcal V^\dagger\mathcal V=1+\order{\chi^3}$，且没有实发射（对无限长直线，实割线从 $\order{g^6}$ 才开始）。
\end{enumerate}
\uses{prop:prelim-wilson-rules,lem:qcd-colour-identities,lem:appO-glauber-operator}
\end{proposition}

\begin{proof}
完整的计算见\cref{sec:appO-wilson}；这里写出它的四个步骤与全部关键中间式。

\emph{(1) 化为纵向积分。}把两条线的路径排序指数各展开到 $\order{g^2}$，四个场两两缩并到不同的线上：$A$ 上较晚的顶点与 $B$ 上较晚的顶点相连（非交叉）或与较早的相连（交叉）。直线的参数积分给出 $(2\pi)^2\delta(v_A\cdot K)\delta(v_B\cdot K)$ 与两个程函传播子（\cref{prop:prelim-wilson-rules}），即\cref{eq:dy-Jbox}中的程函因子；$K$ 是纯横向的，故两条胶子传播子给出 $g(\mathcal Q)$。由 $\dd k^0\dd k^3=\dd x_A\dd x_B/\sinh Y$、$\int\dd^2K_\parallel\,\delta(v_A\cdot K)\delta(v_B\cdot K)=1/\sinh Y$ 与 $(\ii g)^4(v_A\cdot v_B)^2(-\ii)^2=-g^4\cosh^2Y$，
\[
\mathcal A_2=-g^4\coth^2Y\int\frac{\dd^2\mathbf K}{(2\pi)^2}\ee^{\ii\mathbf K\cdot\boldsymbol b}\int\frac{\dd^2\mathbf k_T}{(2\pi)^4}\big[J_{\rm box}c_{\rm box}+J_{\rm cross}c_{\rm cross}\big].
\]
非交叉配对的色因子是 $\sum_{a,b}T^a_AT^b_A\otimes T^a_BT^b_B=X^2$；交叉配对是 $X^2+\sum T^a_AT^b_A\otimes[T^b_B,T^a_B]=X^2+\ii f^{bac}T^a_AT^b_AT^c_B$，而 $\ii f^{bac}T^a_AT^b_A=\frac\ii2f^{bac}[T^a_A,T^b_A]=\frac{C_A}2T^c_A$（\cref{lem:qcd-colour-identities}），故 $c_{\rm cross}=X^2+\frac{C_A}2X$。

\emph{(2) $x_B$ 的留数。}作为 $x_B$ 的函数，$g=\sinh^4Y\big/\prod_{c=a,b}(x_B-x^c_+)(x_B-x^c_-)$，$x^c_\pm=\cosh Y\,x_A\pm\sinh Y\,R_c$，$R_c=\sqrt{x_A^2-c+\ii0}$（$\Imag R_c\ge0$），$x^c_+$ 在上半平面、$x^c_-$ 在下半平面；又 $(x^a_+-x^b_+)(x^a_+-x^b_-)=(b-a)\sinh^2Y$。交叉图（程函极点在下半平面）在上半平面闭合，非交叉图在下半平面闭合：
\[
\begin{aligned}
I_{\rm cross}(x_A)&\equiv\int\dd x_B\frac{\ii\,g}{x_B+\ii0}=\frac{\pi\sinh Y}{a-b}\big[h^+_a-h^+_b\big],\\
I_{\rm box}(x_A)&\equiv\int\dd x_B\frac{\ii\,g}{-x_B+\ii0}=\frac{\pi\sinh Y}{a-b}\big[h^-_b-h^-_a\big],
\end{aligned}
\]
其中 $h^\pm_c\equiv\frac{1}{R_cx^c_\pm}=\frac{\cosh Y\,x_A}{R_c(x_A^2+c\sinh^2Y)}\mp\frac{\sinh Y}{x_A^2+c\sinh^2Y}$。

\emph{(3) $x_A$ 积分。}以 $\frac{\ii}{x_A+\ii0}=\ii\,{\rm PV}\frac1{x_A}+\pi\delta(x_A)$，$J=\ii\,{\rm PV}\!\int\frac{I(x_A)}{x_A}\dd x_A+\pi I(0)$。$h^\pm_c$ 的第二项为偶函数，主值积分为零；第一项给出唯一需要的积分
\[
{\rm PV}\!\int\frac{h^\pm_c}{x_A}\dd x_A=\frac{\cosh Y}{c}\int_{-\infty}^{\infty}\frac{\dd u}{\sqrt{u^2-1+\ii0}\,(u^2+\sinh^2Y)}=\frac{2Y-\ii\pi}{c\sinh Y},
\]
其中 $|u|<1$ 段（$u=\sin\varphi$）给出 $-\frac{\ii\pi}{\sinh Y\cosh Y}$，$|u|>1$ 段（$u=\cosh t$，$w=\tanh t$）给出 $\frac{2Y}{\sinh Y\cosh Y}$；又 $h^\pm_c(0)=\mp\frac1{c\sinh Y}$。代入得
\[
J_{\rm cross}=-\frac{\pi(2\ii Y+\pi)}{ab}+\frac{\pi^2}{ab}=-\frac{2\pi\ii Y}{ab},\qquad
J_{\rm box}=\frac{\pi(2\ii Y+\pi)}{ab}+\frac{\pi^2}{ab}=\frac{2\pi(\pi+\ii Y)}{ab},
\]
即 (a)。作为检验，$J_{\rm box}+J_{\rm cross}=2\pi^2/(ab)$ 正是阿贝尔情形所要求的：$\frac{\ii}{-x_B+\ii0}+\frac{\ii}{x_B+\ii0}=2\pi\delta(x_B)$。

\emph{(4) 横向积分与幺正性。}单胶子交换同样给出
\[
\mathcal A_1=(\ii g)^2(v_A\cdot v_B)\frac{1}{\sinh Y}\int\frac{\dd^2\mathbf k_T}{(2\pi)^2}\,\ee^{\ii\mathbf k_T\cdot\boldsymbol b}\,\frac{-\ii}{-\mathbf k_T^2-m^2}\,X=-\ii\chi X ;
\]
由卷积定理，$g^4\coth^2Y\int\frac{\dd^2\mathbf K}{(2\pi)^2}\ee^{\ii\mathbf K\cdot\boldsymbol b}\int\frac{\dd^2\mathbf k_T}{(2\pi)^2}\frac{1}{ab}=\chi^2$，于是 $\mathcal A_2=-\frac12\chi^2X^2+\frac{\ii C_AY}{4\pi}\chi^2X$，即 (b)；与参数化比较，$-\frac{\chi^2}{2}p_{\rm opp}\frac{C_A}{2}=\frac{\ii C_AY}{4\pi}\chi^2$，$p_{\rm opp}=-\ii Y/\pi$。$\chi$ 为实数、$X$ 为 Hermite 矩阵，故 $\mathcal A_2+\mathcal A_2^\dagger=-\chi^2X^2=-\mathcal A_1^\dagger\mathcal A_1$。最后，一条无限长直线只吸收或放出满足 $v\cdot k=0$ 的动量，而质量为 $m$ 的在壳胶子在 $v$ 的静止系中 $k^0\ge m>0$，故 $\order{g}$ 与 $\order{g^2}$ 的实发射振幅为零，实割线最早出现在 $\order{g^3}$ 振幅的平方即 $\order{g^6}$。
\end{proof}

我们独立地核对了证明中的每一个公式：$u$ 积分的闭式与数值积分在 $Y=0.3,1,2.5,5$ 处相差小于 $3\times10^{-17}$；第 (2) 步的留数公式与对 $x_B$ 的直接数值积分（以 $\ii\varepsilon$、$\varepsilon=10^{-7}$ 正规化）在六个 $x_A$ 值上相差 $\order{\varepsilon}$；第 (3) 步的外层积分数值计算在四组 $(a,b,Y)$ 上重现 (a) 的闭式到 $2\times10^{-16}$；在全部分母中保留有限 $\ii\varepsilon$ 的二维直接积分外推到 $\varepsilon\to0$，与闭式相差 $3\times10^{-4}$；色因子在 SU(2)、SU(3) 的基础、反基础与伴随表示的各种组合中精确成立。

这一结果是已知的：Meggiolaro 计算了两条程函线之间的 $\order{g^4}$ 振幅\cite[第~3~节，式~(3.23)--(3.26)、(3.34)、(3.36)]{Meggiolaro:1996hf}，Giordano 与 Meggiolaro 给出了逐图的形式\cite[式~(4.24)]{Giordano:2006vv}；上面是本书给出的推导。高能极限下的盒图与交叉盒图（$\ln(-s)$ 与 $\ln s$）是经典结果\cite{Cheng:1969eh,Korchemskaya:1994qp,Korchemsky:1993hr,Caron-Huot:2013fea}。$\order{g^4}$ 的其余程函图（单胶子交换的单圈修正）给出树图振幅的实数倍，其形状因子对 $t<0$ 为实数\cite[式~(3.28)--(3.34)]{Meggiolaro:1996hf}；它们对 $\mathcal A+\mathcal A^\dagger$ 没有贡献，所以对完整的程函振幅 $\order{g^4}$ 的幺正性也成立——这一部分依据所引的计算，本书没有重新推导。\cref{fig:dy-two-gluon-exchange}画出了两个图及其纵向积分与色因子。

%%% BEGIN fig:dy-two-gluon-exchange
\begin{figure}[htbp]
\centering
\begin{tikzpicture}[fd,
  fc lab/.style={font=\small},
  fc note/.style={font=\footnotesize},
  fc tiny/.style={font=\scriptsize}]
  % =====================================================================
  % (a) box: the two gluons are attached in the same order on both lines
  % =====================================================================
  \begin{scope}
    \draw[gluon] (-0.8,-0.85) -- (-0.8,0.85);
    \draw[gluon] (0.8,-0.85) -- (0.8,0.85);
    \draw[wilson=0.14] (-2.1,0.85) -- (2.1,0.85);
    \draw[wilson=0.14] (-2.1,-0.85) -- (2.1,-0.85);
    \foreach \p in {(-0.8,0.85),(0.8,0.85),(-0.8,-0.85),(0.8,-0.85)} \node[vertex] at \p {};
    \path[mom={$k$},mom sep=0.27cm] (-0.8,-0.85) -- (-0.8,0.85);
    \path[mom'={$K-k$},mom sep=0.27cm] (0.8,-0.85) -- (0.8,0.85);
    \node[fc lab,anchor=east] at (-2.15,0.85) {$v_A,\,\mathbf T_A$};
    \node[fc lab,anchor=east] at (-2.15,-0.85) {$v_B,\,\mathbf T_B$};
    \node[fc tiny] at (0,1.33) {$\dfrac{\ii}{x_A+\ii0}$};
    \node[fc tiny] at (0,-1.35) {$\dfrac{\ii}{-x_B+\ii0}$};
    \node[fc lab] at (0,-2.25) {$J_{\rm box}=\dfrac{2\pi}{ab}\,(\pi+\ii Y)$};
    \node[fc lab] at (0,-3.15) {$c_{\rm box}=X^2$};
    \node[fc lab] at (0,-3.95) {(a)\ 非交叉图（box）};
  \end{scope}
  \node[fc lab] at (3.35,0) {$+$};
  % =====================================================================
  % (b) crossed box
  % =====================================================================
  \begin{scope}[xshift=6.7cm]
    \draw[gluon] (-0.8,-0.85) -- (0.8,0.85);
    % the gluon k passes over the gluon K-k
    \draw[white,line width=7.5pt] (0.2,-0.21) -- (-0.2,0.21);
    \draw[gluon] (0.8,-0.85) -- (-0.8,0.85);
    \draw[wilson=0.14] (-2.1,0.85) -- (2.1,0.85);
    \draw[wilson=0.14] (-2.1,-0.85) -- (2.1,-0.85);
    \foreach \p in {(-0.8,0.85),(0.8,0.85),(-0.8,-0.85),(0.8,-0.85)} \node[vertex] at \p {};
    \path[mom={$k$},mom pos=0.78,mom sep=0.27cm] (0.8,-0.85) -- (-0.8,0.85);
    \path[mom'={$K-k$},mom pos=0.78,mom sep=0.27cm] (-0.8,-0.85) -- (0.8,0.85);
    \node[fc lab,anchor=east] at (-2.15,0.85) {$v_A,\,\mathbf T_A$};
    \node[fc lab,anchor=east] at (-2.15,-0.85) {$v_B,\,\mathbf T_B$};
    \node[fc tiny] at (0,1.33) {$\dfrac{\ii}{x_A+\ii0}$};
    \node[fc tiny] at (0,-1.35) {$\dfrac{\ii}{x_B+\ii0}$};
    \node[fc lab] at (0,-2.25) {$J_{\rm cross}=-\dfrac{2\pi\ii Y}{ab}$};
    \node[fc lab] at (0,-3.15) {$c_{\rm cross}=X^2+\dfrac{C_A}{2}\,X$};
    \node[fc lab] at (0,-3.95) {(b)\ 交叉图（crossed box）};
  \end{scope}
  % =====================================================================
  % sum
  % =====================================================================
  \node[fc lab] at (3.35,-5.15) {$J_{\rm box}\,c_{\rm box}+J_{\rm cross}\,c_{\rm cross}
     =\dfrac{2\pi}{ab}\Big[\underbrace{\pi X^2}_{\mbox{\footnotesize 指数化项}}\;\underbrace{-\,\ii Y\,\tfrac{C_A}{2}\,X}_{\mbox{\footnotesize Regge 相位}}\Big]$};
\end{tikzpicture}
\caption{两条程函线之间的双胶子交换（\cref{prop:dy-eikonal-twogluon}）。两条无限长类时 Wilson 线的速度为$v_A$、$v_B$（快度差$Y>0$，$v_A\cdot v_B=\cosh Y$），色生成元为$\mathbf T_A$、$\mathbf T_B$，时间自左向右；两个胶子取质量$m$（红外正规子）并保留完整传播子，横动量$\mathbf k_T$、$(\mathbf K-\mathbf k)_T$固定，$a=\mathbf k_T^2+m^2$，$b=(\mathbf K-\mathbf k)_T^2+m^2$。
线段旁为两顶点之间的程函传播子（$x_{A,B}=v_{A,B}\cdot k$），其下为纵向积分$J$与色因子$c$（$X=\mathbf T_A\cdot\mathbf T_B$）。
全部$s$道权重在非交叉图中，交叉图只贡献单圈 Regge 相位$\propto C_AY$：横向积分后第一项是指数化项$\frac12(-\ii\chi)^2X^2$，第二项是$\ii\times$实数$\times$厄米色矩阵；因而$p_{\rm opp}=-\ii Y/\pi$，${\rm Re}\,p_{\rm opp}=0$，$\order{g^4}$的幺正性没有残差（VF-REF-9，VF-REF-10）。}
\label{fig:dy-two-gluon-exchange}
\end{figure}
%%% END fig:dy-two-gluon-exchange

\begin{remark}[共同时序是物理结果]\label{rem:dy-common-ordering}
(i) 精确的图把全部 $s$ 道权重（$J_{\rm box}$ 中的 $\pi$）放进非交叉图，交叉图只携带 Regge 化相位。“严格窗口”去掉了连接线近质壳 $2k^+k^-\simeq\mathbf k_T^2$ 的区域，而正是这一区域把两个喷注中插入的次序关联起来；去掉它之后，两个喷注中的次序彼此独立，剩下的权重只由正规子决定：对称正规子给出 $p_{\rm opp}=\frac12$，一般形状给出 $p_{\rm opp}=\frac12+2\delta_A\delta_B$（\cref{prop:appO-strict-window-shape}，\tagR，$R_{\mathrm{PT}}$）。所以 $\frac12$ 不是普适值，残差 $-\frac14C_A\chi^2X$ 是截断的产物。Rothstein 与 Stewart 的快度正规子与精确结果一致：交叉的 Glauber 盒图为零，梯形图指数化为一个 Hermite 的色相位\cite[式~(5.30)、(9.1)、(9.18)、(9.19)]{RothsteinStewart2016}，即共同时序 $p_{\rm opp}=0$。
(ii) 对任意阶的程函旁观者，幺正性来自 ${\rm T}\exp\{\ii g\int\dd t\,[v_A\cdot A(v_At)\mathbf T_A+v_B\cdot A(v_Bt+\boldsymbol b)\mathbf T_B]\}$ 的幺正性，其形式水平与\cref{lem:method-soft-unitarity}相同；它在所有阶没有作为定理写出。
(iii) 真正开放的问题是非程函（复合的、相互作用的）旁观者与活动部分子的 Glauber、软、共线模式之间的相互作用（\cref{asm:appO-mixing}(c)）。对复合旁观者，附录~O 在一个全局两时模型中证明了：在两个喷注的末态能量差都为负的构型上逐点地、$\order{g^4}$，完全重合的模式中\emph{未割}直接线的侧求和相消（\cref{prop:appO-twoexch-spectator}，模型定理，作为 QCD 陈述为 \tagP；其物理解读还依赖共线端点的匹配假设\cref{asm:appO-collinear-ends}）；被割的直接线（\cref{rem:appO-cut-ss}）与 $\order{g^6}$（\cref{prop:appO-g6-ladder}，\tagA）仍然开放。所以，“非阿贝尔双 Glauber 交换留下非幺正残差”是对严格截断的误读，而不是 DY 因子化的障碍。
\end{remark}

% =====================================================================
\section{主定理与其地位}\label{sec:dy-theorem}
% =====================================================================

现在可以把前几节的结果装配成单举 DY 的因子化定理。装配本身是一个演绎：把模板定理\cref{thm:method-template}的假设 (H1)--(H5)（该定理本身是全阶论证，\tagA）逐条对应到本章的陈述。这里的难点不在演绎，而在于如实交代每一条假设的地位。我们按四个层次给出这一交代：文献中的论证；把它写成的精确条件定理链；以标签列出的、尚未证明的输入；以及反例的情况。

\begin{theorem}[单举 DY 的共线因子化 \tagA]\label{thm:dy-factorization}
对单举过程 $A+B\to\gamma^*+X$（$Q\gg\Lqcd$，$\tau$ 与 $y$ 固定），
\begin{equation}\label{eq:dy-factorization}
\frac{\dd\sigma}{\dd Q^2\,\dd y}=\sum_{i,j}\int_{x_A}^1\frac{\dd\xi_A}{\xi_A}\int_{x_B}^1\frac{\dd\xi_B}{\xi_B}\,f_{i/A}(\xi_A;\mu)\,f_{j/B}(\xi_B;\mu)\,\mathcal H_{ij}\Big(\frac{x_A}{\xi_A},\frac{x_B}{\xi_B},\frac Q\mu,\as(\mu)\Big)+\order{\Big(\frac{\Lambda}{Q}\Big)^p},
\end{equation}
其中 $f_{i/A}$、$f_{j/B}$ 是\cref{def:pdf-quark,def:pdf-gluon}中的 $\MSbar$ 部分子分布，与 DIS 中的是同一组函数；$\mathcal H_{ij}$ 与外态无关，可用部分子态逐阶计算；余项的幂次 $p>0$ 不由证明确定（\tagC，预期 $p=2$）。\\
非微扰地位：\cref{asm:intro-bridge}（MA1、MA2、MA4）。
% errtype: E-S
\uses{thm:method-template,prop:dy-regions,lem:dy-fsi,prop:dy-deformation,lem:dy-soft,prop:dy-universality,thm:appO-main,cor:appO-interface,prop:appN-dy,def:pdf-quark,def:pdf-gluon,asm:intro-bridge,asm:appO-lcopt,asm:appO-envelope,asm:appO-mixing,asm:appO-deformation,asm:appO-rewritten-cs,asm:appN-collins,asm:appN-cutsum}
\end{theorem}

\begin{proofsketch}
由\cref{thm:method-template}（全阶论证，\tagA），只需对 DY 验证其假设 (H1)--(H5)，并把各步的精确形式接到附录的条件定理上。
(H1) 领头区域由\cref{prop:dy-regions}给出，其中包含旁观者之间的 Glauber 子区域；积分范数下的幂次计数是模板定理 (H1) 的前提。
(H2) $H$ 与 $J$ 的接口处的共线近似与纵向胶子的退耦同 DIS：$A$ 共线的纵向胶子附着在 $H$ 或 $J_B$ 上，在 $A$ 的共线因子中给出指向过去的 $W[\cdot,-\infty n_B]$；由\cref{prop:dy-universality}，它与 DIS 中的 $W[\cdot,\infty n_B]$ 给出同一个算符。
(H3) Glauber 区：割图求和之后把软动量沿 $e_G$ 形变出 Glauber 区，软胶子退耦为指向过去的 Wilson 线（\cref{lem:dy-fsi,prop:dy-deformation}）。
(H4) 不受测量约束的末态相互作用由幺正性相消，对 $q_T$ 积分的截面软因子为 $1$（\cref{lem:dy-fsi,lem:dy-soft}）。
(H5) 减除项与部分子分布算符的重整化矩阵元一致，同 DIS；DY 所需的积分形式是\cref{prop:appN-dy}：在\cref{asm:appN-collins}、\cref{asm:appN-cutsum}、\cref{asm:appO-rewritten-cs}、\cref{asm:appO-deformation}(b)(c) 之下，$\big|\int F_G-\sum_R\int C_RF_G\big|\le C\lambda^aL^K$，并且在每个由 $R_0$ 主导的块上，减除项可以用单侧改写的形式计算（\cref{cor:appO-interface}(e)），其中 Wilson 线指向过去。
把 (H1)--(H5) 代入模板定理，软因子为 $1$，得 $\sigma=\sum_{ij}\mathcal H_{ij}\otimes f_{i/A}\otimes f_{j/B}$ 加上被幂次压低的余项；$\mathcal H_{ij}$ 是部分子截面减去部分子分布的卷积（质量因子化），因而与外态无关。各步的精确条件形式依次是\cref{thm:appO-main}、\cref{cor:appO-interface}、\cref{prop:appN-dy}（均为 \tagR，$R_{\mathrm{cond}}$），其前提为\cref{asm:appO-lcopt}、\cref{asm:appO-envelope}、\cref{asm:appO-mixing}、\cref{asm:appO-deformation}、\cref{asm:appO-rewritten-cs}以及上面的附录~N 输入。完整的论证见\cite[第~14.4~节]{Collins2011}；精确形式见\cref{sec:appO-main}。
\end{proofsketch}

\paragraph{地位与缺口}\tagA。我们分四层说明。

\emph{(1) 文献中的论证。}CSS（1985、1988）与 Bodwin（1985）给出了全阶论证的关键步骤\cite{CollinsSoperSterman1985b,CollinsSoperSterman1988,Bodwin1985}：Glauber 区的围道变形、末态相互作用的幺正性相消、软因子为 $1$；两圈上旁观者的非因子化项之和为零\cite{Lindsay:1982pu,Lindsay:1982gd}，程函“介子”模型中的两圈相消见\cite{Brandt:1988xt}（均据其摘要）。Collins 的教科书说明，其第 14.4 节的处理遵循 CSS 1988，并作了改进、改正了其中的错误（他没有具体列出这些错误）\cite[第~14.4~节，第~553~页]{Collins2011}。他在第 14.4.8 节“Issues remaining”中列出：(i) 证明依赖微扰论的动量分类，非微扰的“分级快度”Glauber 交换（多周边、类 Regge）没有被覆盖；(ii) Feynman 规范下领头区域不完备，额外的不连通硬散射须对图求和后才相消\cite{Labastida:1984gy}；(iii) “The current versions of genuine proofs are evidently formidable”（现有的真正证明显然是令人生畏的）\cite[第~14.4.8~节]{Collins2011}。Rothstein 与 Stewart 的 Glauber 算符框架复现了旁观者相位的相消，但他们自己说明这一相消“taken alone it is far from a proof of factorization, even in the abelian case”（单独来看远不是因子化的证明，即使在阿贝尔情形）\cite[第~11.1~节]{RothsteinStewart2016}；据我们截至 2026 年 10 月在 INSPIRE 上的文献检索，他们宣布的完整 SCET 证明尚未发表。

\begin{table}[!tbp]
\centering
\caption{DY 证明的逐步状态（附录~O 的步骤 (O-1)--(O-8)，据表~\ref{tab:appO-checklist}）}
\label{tab:dy-status}
\footnotesize
\setlength{\tabcolsep}{3pt}
\renewcommand{\arraystretch}{1.18}
\begin{tabularx}{\textwidth}{@{}>{\raggedright\arraybackslash}p{1.95cm}>{\raggedright\arraybackslash}X>{\raggedright\arraybackslash}p{2.1cm}>{\raggedright\arraybackslash}p{1.25cm}>{\raggedright\arraybackslash}p{2.75cm}>{\raggedright\arraybackslash}p{2.0cm}@{}}
\toprule
步骤 & 陈述（标签） & 地位 & 误差类型 & 所用假设 & 文献\\
\midrule
(O-1) 区域、定义域与窗口 & \cref{prop:dy-regions}；\cref{def:appO-leading-data}、\cref{def:appO-windows}、\cref{lem:appO-glauber-pc} & 定义；幂次计数为标度陈述 & E-S；E-P & \cref{asm:appN-envelope}、\cref{asm:appN-separation} & \cite{CollinsSoperSterman1985b}；\cite[14.3~节]{Collins2011}\\
(O-2) 光锥序中的领头幂次形式 & \cref{lem:appO-lcopt}、\cref{lem:appO-hard-locality}、\cref{lem:dy-poles-orientation} & 在假设下完整 & 精确 & \cref{asm:appO-lcopt} & \cite{Kogut:1969xa,Chang:1973qi,Ligterink:1994tm,AybatSterman2009}\\
(O-3) 割图求和与单侧改写 & \cref{lem:dy-fsi}（\cref{lem:appO-unitarity}）、\cref{lem:appO-equal-time}、\cref{lem:appO-joint-localization}、\cref{prop:appO-ss-cancel}、\cref{prop:appO-one-sided} & 在假设下完整 & 精确；E-F；E-W & \cref{asm:appO-lcopt}、\cref{asm:appO-envelope}、\cref{asm:appO-mixing} & \cite[14.4.1~节]{Collins2011}；\cite[4~节]{Diehl:2015bca}\\
(O-4) 减除项与改写相容 & \cref{prop:appO-one-sided}（链项） & 在假设下完整 & E-W & \cref{asm:appN-separation}(ii)、\cref{asm:appO-envelope} & ---\\
(O-5) 形变目标 & \cref{prop:dy-deformation}；\cref{lem:appO-no-crossing}、\cref{prop:appO-target} & 对 $\mathcal G_0$ 与分级构型完整，一般情形为假设 & 精确（几何） & \cref{asm:appO-deformation}(a)--(d) & \cite{CollinsSoperSterman1988,AybatSterman2009}\\
(O-6) 程函近似与指向过去的 Wilson 线 & \cref{prop:appO-eikonal}、\cref{prop:appO-past-wilson} & 程函界完整；Wilson 线为 \tagA & E-P & \cref{asm:appO-rewritten-cs}、\cref{asm:appN-ward-product} & \cite{GrammerYennie1973}；\cite[14.4~节]{Collins2011}\\
(O-7) 软因子为 1；共线因子为 DIS 的 PDF & \cref{lem:dy-soft}（\cref{lem:method-soft-unitarity}）、\cref{prop:dy-universality} & 完整（$R_{\mathrm{def}}$） & 精确 & --- & \cite[14.4.6~节]{Collins2011}\\
(O-8) 与附录~N 的接口 & \cref{thm:appO-main}、\cref{cor:appO-interface}、\cref{prop:appN-dy} & 在假设下完整（$R_{\mathrm{cond}}$） & E-W$\to$E-N & 以上全部，及\cref{asm:appN-collins}、\cref{asm:appN-cutsum} & ---\\
\bottomrule
\end{tabularx}
\end{table}

\emph{(2) 精确形式。}\cref{app:appO}把 (H3)(H4) 写成条件定理\cref{thm:appO-main}（\tagR，$R_{\mathrm{cond}}$），经接口\cref{cor:appO-interface}接到\cref{app:appN}的\cref{prop:appN-dy}（\tagR，$R_{\mathrm{cond}}$），由此得到因子化形式。因此，DY 的共线因子化等于“割图求和后的 Glauber 相消”（\cref{app:appO}）加上减除方法中组合与幂次计数的部分，而后者与 $e^+e^-$、DIS 中的完全相同。这条链的内容等级是 \tagA。

\emph{(3) 以标签列出的开放输入。}\cref{asm:appO-lcopt}（光锥序微扰论的适用性）；\cref{asm:appO-envelope}（纤维积分误差的可积包络）；\cref{asm:appO-mixing}(c)（重合连接的完整图集相消；对程函线之间的双胶子交换在 $\order{g^4}$ 成立，见\cref{prop:dy-eikonal-twogluon}；对复合旁观者的未割线，在全局两时模型中有模型证据，见\cref{prop:appO-twoexch-spectator}）；\cref{asm:appO-deformation}(b)(c)(d)（\emph{即使对阿贝尔图也需要}）；\cref{asm:appO-rewritten-cs}；附录~N 的一般输入\cref{asm:appN-collins}、\cref{asm:appN-cutsum}，以及经\cref{thm:method-template}进入的\cref{asm:appN-envelope}、\cref{asm:appN-separation}、\cref{asm:appN-ward-product}、\cref{asm:appN-values}、\cref{asm:appN-uv}、\cref{asm:appN-multren}；非微扰的桥接\cref{asm:intro-bridge}。此外，该论证中被明确假定的一步（\cref{asm:appN-brst-nonlocal}，\tagC，多定域光线算符的 BRST 分类）是\cref{thm:method-template}中点名的具名缺口，在此重申；它不是本定理的前提。除约定\cref{asm:appN-values}之外，全部技术性输入都是 \tagA，所以本定理是 \tagA。各步的误差类型见\cref{tab:dy-status}：标度陈述（E-S）、逐点界（E-P）、纤维积分界（E-F）与弱等价（E-W）。

\emph{(4) 反例。}没有已知的领头幂次反例；文献中“严格窗口”的非幺正残差是截断的假象（\cref{prop:dy-eikonal-twogluon}）。混合光锥序与“钳夹变为初、末态形变方向的不匹配”的看法应归功于 Aybat 与 Sterman\cite{AybatSterman2009}（这是一个论证纲要，不是证明）；Diehl 等人给出了一个光锥微扰论中逐图进行的论证，其中不出现次序权重\cite[第~4~节]{Diehl:2015bca}（与\cref{app:appO}框架的比较见\cref{rem:appO-literature-comparison}）。余项的幂次：重整化子分析表明，重求和的 DY 截面中领头的红外重整化子相消，高扭度修正预期为 $\Lambda^2/[Q^2(1-z)^2]$ 量级（$Q(1-z)>\Lambda$ 时）\cite{BenekeBraun1995b}，这是证据而不是证明；$1/Q^2$ 处的广义因子化是\cref{thm:dy-higher-twist}，其外的幂次上因子化被破坏（\cref{prop:dy-twist6-failure}）。要把本定理提升到条件链完全成立的形式，需要证明 (3) 中的各项；要把它变成关于强子的陈述，还需要\cref{asm:intro-bridge}。



\cref{tab:dy-status}中，“在假设下完整”指该步在所列编号假设之下有完整证明；Glauber 特有的未证明输入是\cref{asm:appO-mixing}(c)、\cref{asm:appO-deformation}(a)--(d)、一般软子图的\cref{asm:appO-mixing}(a$'$)(b)、\cref{asm:appO-envelope}与\cref{asm:appO-rewritten-cs}，而附录~N 的输入与 DIS、$e^+e^-$ 共享。

\begin{remark}[固定阶部分子截面不检验旁观者 Glauber 相消]\label{rem:dy-nnlo-not-glauber}
DY 的部分子截面已经算到 NNLO\cite{HambergVanNeervenMatsuura1991}与 N$^3$LO（中性流\cite{Duhr:2020seh}与带电流\cite{Duhr:2020sdp}），质量因子化之后都是有限的，并与数据相符。但部分子截面里没有强子的旁观者：其有限性检验的是活动部分子的 Glauber（Coulomb）贡献在该阶的相消——例如单圈形状因子的虚部，即活动 $q\bar q$ 对之间红外发散的 Coulomb 相位，在与 Born 振幅的干涉 $2\Real F$ 中消去（\cref{sec:dy-nlo}）——而不检验旁观者之间的交换；后者是\cref{prop:dy-abelian-model}与\cref{app:appO}的内容。所以高阶计算与数据的符合是对整个框架的证据，而不是对旁观者 Glauber 相消的证明；“NNLO/N$^3$LO 的 DY 与数据一致，说明 Glauber 相消已被检验”是一个误解。
\end{remark}

\begin{remark}[Doria--Frenkel--Taylor 非相消的前提]\label{rem:dy-dft-premises}
Doria、Frenkel 与 Taylor 给出了一个反应，其中两个有质量、颜色平均的初态夸克的两软胶子非领头红外发散在实、虚图之间不相消\cite{DoriaFrenkelTaylor1980}；这一结论也被\cite{DiLietoEtAl1981}（他们在若干中间结果上与前者不一致）与\cite{Andrasi:1980qw}得到。这些非相消项在高能下被幂次压低：Frenkel、Gatheral 与 Taylor 表明到 $\order{g^4}$ 为止除了幂次压低的项之外是有限的\cite{Frenkel:1983di}，Catani 等人称之为“higher twist type”\cite{Catani:1985xt}（以上均据各文摘要）；对 DY，CSS 指出两圈的非相消出现在相对 $m^4/s^2$ 的量级\cite[第~10.2~节]{CollinsSoperSterman1989}。这一情形在\cref{thm:appO-main}的前提之外：后者处理由色单态顶点产生的强子（其旁观者携带补偿的色荷，并可与软胶子耦合）中无质量的活动部分子，并且只断言领头幂次；DFT 的外态则是有质量、有色的入射夸克，对其颜色求平均并不等价于色单态强子。在有质量的入射线上，$p_A=(p_A^+,\frac{m^2}{2p_A^+},\mathbf 0_T)$，
\begin{equation}\label{eq:dy-dft}
(p_A-k)^2-m^2=-2p_A^+\Big(k^-+\frac{m^2}{2(p_A^+)^2}\,k^+\Big)+k^2 ,
\end{equation}
% VF-REF-8
其中 $\frac{m^2}{2(p_A^+)^2}k^+$ 一项把 $k^+$ 与 $k^-$ 绑在一起，破坏了“$J_A$ 只通过 $k^-$ 感知软动量”这一使两个喷注可以分别按各自的光锥时间排序的性质（\cref{lem:dy-poles-orientation}所依据的领头幂次近似）；它是幂次压低的。对强子的 DY，因子化在 $m^4/s^2$ 处失效：在标量夸克模型中，因子化在 $m^2/s$ 成立而在 $m^4/s^2$ 失效\cite{Basu:1984ba}（\cref{prop:dy-twist6-failure}）。比较\cref{prop:ir-bn-failure}（\tagV）。
\end{remark}

% =====================================================================
\section{\texorpdfstring{$\order{\as}$}{O(alpha_s)} 系数函数}\label{sec:dy-nlo}
% =====================================================================

\cref{thm:dy-factorization}中唯一可以逐阶算出的因子是硬函数。本节在 $\MSbar$ 方案中把它算到 $\order{\as}$。这一计算只涉及部分子态，与\cref{thm:dy-factorization}的地位无关：取无质量部分子作初态、用维数正规化正规化其红外发散，我们将看到，软发散在实、虚修正之间相消，剩下的初态共线发散恰好由两个 $\MSbar$ 部分子分布的抵消项消去——这本身就是因子化在单圈上的一个非平凡检验。对 $y$ 积分之后硬函数只依赖于 $z=Q^2/\hat s$；按\cref{def:dy-kin}的归一化，因子化公式写成
\begin{multline}\label{eq:dy-fact-nlo}
\frac{\dd\sigma}{\dd Q^2}=\frac{\sigma_0}{s}\sum_qe_q^2\int_0^1\!\dd x_1\,\dd x_2\,\dd z\;\delta(x_1x_2z-\tau)\\
\times\Big\{\big[q\bar q+\bar qq\big]\Delta_{q\bar q}(z)+\big[g\,(q+\bar q)+(q+\bar q)\,g\big]\Delta_{qg}(z)\Big\}+\order{\as^2},
\end{multline}
% VF-DY-8
方括号中第一个因子取 $x_1$（强子 $A$）、第二个取 $x_2$（强子 $B$），部分子分布都取在标度 $\mu_F$；$\Delta_{qg}$ 是一个味为 $q$ 的夸克（或反夸克）与一个胶子的系数。等价地，$\hat s\,\dd\hat\sigma_{ij}/\dd Q^2=\sigma_0e_q^2\Delta_{ij}(z)$ 是扣除共线发散之后的部分子截面。记 $a\equiv\as(\mu_R)/(2\pi)$，$\mathcal D_k(z)\equiv[\ln^k(1-z)/(1-z)]_+$（\cref{def:prelim-plus}）；光滑函数乘加号分布一律按乘积约定理解，$g\mathcal D_k=g(1)\mathcal D_k+\frac{g-g(1)}{1-z}\ln^k(1-z)$。图形见\cref{fig:dy-oneloop}。

%%% BEGIN fig:dy-oneloop
\begin{figure}[htbp]
\centering
\begin{tikzpicture}[fd,
  t2 lab/.style={font=\footnotesize}]
  % ======================= row 1:  q qbar channel =======================
  % (a) q qbar -> gamma* g, gluon off the quark
  \begin{scope}
    \coordinate (V) at (0,0); \coordinate (E1) at (-2.0,1.15); \coordinate (E2) at (-2.0,-1.15);
    \coordinate (G) at ($(V)!0.45!(E1)$);
    \draw[quark] (E1) -- (G); \draw[quark] (G) -- (V);
    \draw[quark] (V) -- (E2);
    \draw[photon] (V) -- (1.45,0);
    \draw[gluon] (G) -- (0.55,1.2);
    \foreach \p in {V,G} \node[vertex] at (\p) {};
    \path[mom={$p_1$},mom pos=0.45] (E1) -- (G);
    \path[mom'={$p_2$},mom pos=0.6] (E2) -- (V);
    \path[mom'={$q$},mom sep=0.28cm] (V) -- (1.45,0);
    \path[mom={$k$},mom pos=0.55] (G) -- (0.55,1.2);
    \node[t2 lab,anchor=east] at (E1) {$q$};  \node[t2 lab,anchor=east] at (E2) {$\bar q$};
    \node[t2 lab,anchor=west] at (1.45,0) {$\gamma^*$}; \node[t2 lab,anchor=west] at (0.55,1.2) {$g$};
    \node at (-0.3,-1.8) {(a)};
  \end{scope}
  % (b) q qbar -> gamma* g, gluon off the antiquark
  \begin{scope}[xshift=5.0cm]
    \coordinate (V) at (0,0); \coordinate (E1) at (-2.0,1.15); \coordinate (E2) at (-2.0,-1.15);
    \coordinate (G) at ($(V)!0.45!(E2)$);
    \draw[quark] (E1) -- (V);
    \draw[quark] (V) -- (G); \draw[quark] (G) -- (E2);
    \draw[photon] (V) -- (1.45,0);
    \draw[gluon] (G) -- (0.55,-1.2);
    \foreach \p in {V,G} \node[vertex] at (\p) {};
    \path[mom={$p_1$},mom pos=0.6] (E1) -- (V);
    \path[mom'={$p_2$},mom pos=0.45] (E2) -- (G);
    \path[mom'={$q$},mom sep=0.28cm] (V) -- (1.45,0);
    \path[mom'={$k$},mom pos=0.55] (G) -- (0.55,-1.2);
    \node[t2 lab,anchor=east] at (E1) {$q$};  \node[t2 lab,anchor=east] at (E2) {$\bar q$};
    \node[t2 lab,anchor=west] at (1.45,0) {$\gamma^*$}; \node[t2 lab,anchor=west] at (0.55,-1.2) {$g$};
    \node at (-0.3,-1.8) {(b)};
  \end{scope}
  % (c) virtual vertex correction
  \begin{scope}[xshift=10.0cm]
    \coordinate (V) at (0,0); \coordinate (E1) at (-2.0,1.15); \coordinate (E2) at (-2.0,-1.15);
    \coordinate (G1) at ($(V)!0.56!(E1)$); \coordinate (G2) at ($(V)!0.56!(E2)$);
    \draw[quark] (E1) -- (G1); \draw[quark] (G1) -- (V);
    \draw[quark] (V) -- (G2); \draw[quark] (G2) -- (E2);
    \draw[photon] (V) -- (1.45,0);
    \draw[gluon] (G1) -- (G2);
    \foreach \p in {V,G1,G2} \node[vertex] at (\p) {};
    \path[mom={$p_1$},mom pos=0.45] (E1) -- (G1);
    \path[mom'={$p_2$},mom pos=0.45] (E2) -- (G2);
    \path[mom'={$q$},mom sep=0.28cm] (V) -- (1.45,0);
    \node[t2 lab,anchor=east] at (E1) {$q$};  \node[t2 lab,anchor=east] at (E2) {$\bar q$};
    \node[t2 lab,anchor=west] at (1.45,0) {$\gamma^*$};
    \node at (-0.3,-1.8) {(c)};
  \end{scope}
  % ======================= row 2:  q g channel =======================
  % (d) s-channel
  \begin{scope}[xshift=2.3cm,yshift=-4.1cm]
    \coordinate (V1) at (-0.55,0); \coordinate (V2) at (0.65,0);
    \draw[quark] (-2.0,1.05) -- (V1);
    \draw[gluon] (-2.0,-1.05) -- (V1);
    \draw[quark] (V1) -- (V2);
    \draw[photon] (V2) -- (1.95,1.05);
    \draw[quark] (V2) -- (1.95,-1.05);
    \foreach \p in {V1,V2} \node[vertex] at (\p) {};
    \path[mom={$p_1$},mom pos=0.45] (-2.0,1.05) -- (V1);
    \path[mom'={$p_2$},mom pos=0.45] (-2.0,-1.05) -- (V1);
    \path[mom={$q$},mom pos=0.6] (V2) -- (1.95,1.05);
    \path[mom'={$k$},mom pos=0.6] (V2) -- (1.95,-1.05);
    \node[t2 lab,anchor=east] at (-2.0,1.05) {$q$};  \node[t2 lab,anchor=east] at (-2.0,-1.05) {$g$};
    \node[t2 lab,anchor=west] at (1.95,1.05) {$\gamma^*$}; \node[t2 lab,anchor=west] at (1.95,-1.05) {$q$};
    \node at (0,-1.8) {(d)};
  \end{scope}
  % (e) t-channel
  \begin{scope}[xshift=8.0cm,yshift=-4.1cm]
    \coordinate (V1) at (0,0.72); \coordinate (V2) at (0,-0.72);
    \draw[quark] (-1.85,0.72) -- (V1);
    \draw[photon] (V1) -- (1.85,0.72);
    \draw[quark] (V1) -- (V2);
    \draw[gluon] (-1.85,-0.72) -- (V2);
    \draw[quark] (V2) -- (1.85,-0.72);
    \foreach \p in {V1,V2} \node[vertex] at (\p) {};
    \path[mom={$p_1$},mom pos=0.5] (-1.85,0.72) -- (V1);
    \path[mom'={$p_2$},mom pos=0.5] (-1.85,-0.72) -- (V2);
    \path[mom={$q$},mom pos=0.5] (V1) -- (1.85,0.72);
    \path[mom'={$k$},mom pos=0.5] (V2) -- (1.85,-0.72);
    \node[t2 lab,anchor=east] at (-1.85,0.72) {$q$};  \node[t2 lab,anchor=east] at (-1.85,-0.72) {$g$};
    \node[t2 lab,anchor=west] at (1.85,0.72) {$\gamma^*$}; \node[t2 lab,anchor=west] at (1.85,-0.72) {$q$};
    \node at (0,-1.8) {(e)};
  \end{scope}
\end{tikzpicture}
\caption{Drell--Yan 过程 $\order{\as}$ 修正的部分子图（振幅层次，时间自左向右）。上排为 $q\bar q$ 道：(a)、(b)~$q(p_1)\,\bar q(p_2)\to\gamma^*(q)\,g(k)$ 的两个实辐射图，
(c)~单圈顶点修正（与 Born 振幅干涉；无质量外线自能在维数正规化中为零）。下排为 $qg$ 道 $q(p_1)\,g(p_2)\to\gamma^*(q)\,q(k)$ 的两个图：(d)~$s$ 道（内夸克线动量 $p_1+p_2$），(e)~$t$ 道（内夸克线动量 $p_1-q$）
（$\bar q g$ 道的图由电荷共轭得到）。入射反夸克线上的费米子数箭头背离顶点，动量箭头另行标出。}
\label{fig:dy-oneloop}
\end{figure}
%%% END fig:dy-oneloop

\begin{proposition}[$\order{\as}$ 的 $\MSbar$ 系数函数 \tagR]\label{prop:dy-nlo}
（$R_{\mathrm{PT}}$：在重整化微扰论中对红外正规化的部分子态严格。）本命题是单圈（$\order{\as}$）的显式陈述。（本命题只涉及部分子态。）在常规维数正规化（$D=4-2\varepsilon$，外胶子按 $D-2$ 个极化态平均）与 $\MSbar$ 方案下，\cref{eq:dy-fact-nlo}中的系数为
\begin{align}
\Delta_{q\bar q}&=\delta(1-z)+a\,C_F\Big[4(1+z^2)\mathcal D_1(z)-2\,\frac{1+z^2}{1-z}\ln z+\Big(\frac{2\pi^2}{3}-8\Big)\delta(1-z)\Big]\nonumber\\
&\qquad+2a\,P_{qq}(z)\ln\frac{Q^2}{\mu_F^2}+\order{\as^2},\label{eq:dy-Dqq}\\
\Delta_{qg}&=a\,T_F\Big[\big(z^2+(1-z)^2\big)\ln\frac{(1-z)^2}{z}+\frac12+3z-\frac72z^2\Big]+a\,P_{qg}(z)\ln\frac{Q^2}{\mu_F^2}+\order{\as^2},\label{eq:dy-Dqg}
\end{align}
% VF-DY-8
其中 $P_{qq}=C_F\big[\frac{1+z^2}{1-z}\big]_+=C_F\big[2\mathcal D_0-(1+z)+\frac32\delta(1-z)\big]$ 与 $P_{qg}=T_F\big[z^2+(1-z)^2\big]$ 是领头阶分裂函数（\cref{prop:pdf-lo-kernels}，\tagR，$R_{\mathrm{PT}}$），$4(1+z^2)\mathcal D_1=8\mathcal D_1-4(1+z)\ln(1-z)$（乘积约定）。系数在此阶不依赖于 $\mu_R$。
\uses{def:dy-kin,prop:dy-born,def:prelim-plus,prop:pdf-lo-kernels,prop:sudakov-oneloop,lem:prelim-plus-expansion,prop:pdf-partonic-dr,lem:appC-phase-space}
\end{proposition}

\begin{proof}
分六步。记 $L\equiv\ln(\mu^2/Q^2)$，$\mu$ 是维数正规化的标度，在此阶等于 $\mu_R$。

\emph{第一步：质量因子化的结构。}对部分子初态，\cref{eq:dy-bare}的裸系数含有初态共线发散。因子化公式对部分子“靶”的形式为
\[
\Delta^{\rm bare}_{ij}(z,\varepsilon)=\sum_{k,l}\big(\Gamma_{ki}\otimes\Gamma_{lj}\otimes\Delta_{kl}\big)(z),\qquad
\Gamma_{ki}(x)=\delta_{ki}\delta(1-x)-\frac{a}{\varepsilon}\Big(\frac{\mu^2}{\mu_F^2}\Big)^{\varepsilon}P_{ki}(x)+\order{a^2},
\]
$\otimes$ 为 Mellin 卷积。$\Gamma_{ki}$ 是部分子 $i$ 中部分子 $k$ 的 $\MSbar$ 分布：在纯维数正规化中裸分布为 $\delta_{ki}\delta(1-x)$，$\MSbar$ 重整化留下的纯极点是红外极点（\cref{prop:pdf-partonic-dr}，\tagR，$R_{\mathrm{PT}}$）；因子 $(\mu^2/\mu_F^2)^\varepsilon$ 在此阶把 $\as(\mu_F)$ 换成 $\as(\mu)$。展开到 $\order{a}$：$q\bar q$ 道的两条入射夸克线各有一个抵消项；$qg$ 道中胶子分裂出的反夸克与夸克湮灭，抵消项为 $f_{\bar q/g}\otimes\hat\sigma_{q\bar q}$（$P_{\bar qg}=P_{qg}$）。以上标 $(1)$ 记 $a^1$ 的系数，
\begin{equation}\label{eq:dy-mf}
\Delta^{(1)}_{q\bar q}=\Delta^{{\rm bare},(1)}_{q\bar q}+\frac2\varepsilon\Big(\frac{\mu^2}{\mu_F^2}\Big)^{\varepsilon}P_{qq},\qquad
\Delta^{(1)}_{qg}=\Delta^{{\rm bare},(1)}_{qg}+\frac1\varepsilon\Big(\frac{\mu^2}{\mu_F^2}\Big)^{\varepsilon}P_{qg}.
\end{equation}
% VF-DY-9
抵消项乘的是 $\Delta^{(0)}_{q\bar q}=\delta(1-z)$，即 $D$ 维 Born 项除以它自身；这正是\cref{eq:dy-bare}除以 $(1-\varepsilon)$ 的原因（\cref{prop:dy-born}）。余下的任务是：精确到 $\varepsilon$ 的所有阶算出 $\order{a}$ 的 $\Delta^{\rm bare}$，证明它的极点恰被\cref{eq:dy-mf}抵消，并读出有限部分。

\emph{第二步：实辐射矩阵元。}实辐射图是\cref{fig:dy-oneloop}(a)(b)（$q\bar q\to\gamma^*g$）与 (d)(e)（$qg\to\gamma^*q$）。$q\bar q$ 道记 $\hat t=(p_1-k)^2$、$\hat u=(p_2-k)^2$，$qg$ 道记 $\hat t=(p_2-k)^2$、$\hat u=(p_1-k)^2$，都有 $\hat s+\hat t+\hat u=Q^2$。光子指标的极化求和取 $-g_{\mu\nu}$（流守恒）；唯一外胶子的极化求和也可取 $-g_{\alpha\beta}$，因为只有一个外胶子时 $k_\alpha M^{\mu\alpha}=0$ 与阿贝尔情形相同，不需要鬼场。在 $D$ 维中求迹得 $\sum|H|^2=e_q^2g_0^2C_FN_c\,T_{q\bar q}$ 与 $\sum|H|^2=e_q^2g_0^2T_F(N_c^2-1)\,T_{qg}$，$g_0^2=4\pi\as\mu^{2\varepsilon}S_\varepsilon^{-1}$，
\begin{align}
T_{q\bar q}&=8(1-\varepsilon)\Big[(1-\varepsilon)\Big(\frac{\hat u}{\hat t}+\frac{\hat t}{\hat u}\Big)+\frac{2\hat sQ^2}{\hat t\hat u}-2\varepsilon\Big],\label{eq:dy-Tqq}\\
T_{qg}&=-8(1-\varepsilon)\Big[(1-\varepsilon)\Big(\frac{\hat s}{\hat t}+\frac{\hat t}{\hat s}\Big)+\frac{2\hat uQ^2}{\hat s\hat t}-2\varepsilon\Big],\label{eq:dy-Tqg}
\end{align}
% VF-DY-3
二者由交叉相联系：$T_{qg}(\hat s,\hat t,\hat u)=-T_{q\bar q}(\hat s\to\hat u,\hat t\to\hat s,\hat u\to\hat t)$。代入\cref{eq:dy-bare}，计入初态平均（$q\bar q$：$\frac14\cdot\frac1{N_c^2}$；$qg$：$\frac12\cdot\frac{1}{2(1-\varepsilon)}\cdot\frac{1}{N_c(N_c^2-1)}$），得
\[
\Delta^R_{q\bar q}=\frac{g_0^2C_F}{8\pi(1-\varepsilon)}\int\dd\Phi_2\,T_{q\bar q},\qquad
\Delta^R_{qg}=\frac{g_0^2T_F}{8\pi(1-\varepsilon)^2}\int\dd\Phi_2\,T_{qg}.
\]

\emph{第三步：相空间与 $v$ 积分。}在部分子质心系中以 $v=(1-\cos\theta)/2$ 参数化（$\theta$ 为末态胶子或夸克相对于 $p_1$ 的极角），$D$ 维两体相空间为（\cref{lem:appC-phase-space}）
\[
\dd\Phi_2=\frac{1}{8\pi}\Big(\frac{4\pi}{\hat s}\Big)^{\varepsilon}\frac{(1-z)^{1-2\varepsilon}}{\Gamma(1-\varepsilon)}\,[v(1-v)]^{-\varepsilon}\dd v,
\]
% VF-DY-4
$q\bar q$ 道 $\hat t=-\hat s(1-z)v$、$\hat u=-\hat s(1-z)(1-v)$，$qg$ 道 $\hat u=-\hat s(1-z)v$、$\hat t=-\hat s(1-z)(1-v)$。于是 $q\bar q$ 道的被积函数为
$T_{q\bar q}=8(1-\varepsilon)\big[(1-\varepsilon)\big(\frac{1-v}{v}+\frac{v}{1-v}\big)+\frac{2z}{(1-z)^2v(1-v)}-2\varepsilon\big]$，
对 $v$ 的积分都是 Beta 函数；记 $G_\varepsilon\equiv\Gamma(1-\varepsilon)^2/\Gamma(1-2\varepsilon)$，
\[
\int_0^1\!\frac{(1-v)^{1-\varepsilon}}{v^{1+\varepsilon}}\,\dd v=-\frac{(1-\varepsilon)\,G_\varepsilon}{\varepsilon(1-2\varepsilon)},\qquad
\int_0^1\!\frac{\dd v}{[v(1-v)]^{1+\varepsilon}}=-\frac{2G_\varepsilon}{\varepsilon},\qquad
\int_0^1\![v(1-v)]^{-\varepsilon}\dd v=\frac{G_\varepsilon}{1-2\varepsilon},
\]
从而
\[
\int_0^1\!\dd v\,[v(1-v)]^{-\varepsilon}\,T_{q\bar q}=-\frac{16(1-\varepsilon)}{\varepsilon}\,G_\varepsilon\,\frac{1+z^2+2\varepsilon^2(1-z)^2/(1-2\varepsilon)}{(1-z)^2}.
\]
$v\to0,1$ 的共线奇异性成为 $1/\varepsilon$，$z\to1$ 的软奇异性与相空间的 $(1-z)^{1-2\varepsilon}$ 合成 $(1-z)^{-1-2\varepsilon}$。再用 $\hat s=Q^2/z$ 与 $g_0^2(4\pi/\hat s)^\varepsilon=4\pi\as(\mu^2/Q^2)^\varepsilon z^\varepsilon\ee^{\varepsilon\gamma_E}$，并记 $K_r(\varepsilon)\equiv\ee^{\varepsilon\gamma_E}\Gamma(1-\varepsilon)/\Gamma(1-2\varepsilon)$，得
\begin{equation}\label{eq:dy-real-qq}
\Delta^R_{q\bar q}=-a\,\frac{2C_F}{\varepsilon}\Big(\frac{\mu^2}{Q^2}\Big)^{\varepsilon}K_r(\varepsilon)\,z^{\varepsilon}(1-z)^{-1-2\varepsilon}\Big[1+z^2+\frac{2\varepsilon^2(1-z)^2}{1-2\varepsilon}\Big].
\end{equation}
% VF-DY-5
$qg$ 道的同样计算（所需积分为 $B(1-\varepsilon,-\varepsilon)$、$B(1-\varepsilon,2-\varepsilon)$、$B(2-\varepsilon,-\varepsilon)$ 与 $B(1-\varepsilon,1-\varepsilon)$）给出
\begin{equation}\label{eq:dy-real-qg}
\Delta^R_{qg}=-a\,\frac{T_F}{2\varepsilon}\Big(\frac{\mu^2}{Q^2}\Big)^{\varepsilon}K_r(\varepsilon)\,\frac{z^{\varepsilon}(1-z)^{-2\varepsilon}}{(1-\varepsilon)(1-2\varepsilon)}\Big[2\big(z^2+(1-z)^2\big)-\varepsilon(5z^2-6z+7)+\varepsilon^2(1+z)^2\Big];
\end{equation}
% VF-DY-5
多出的 $\frac1{1-\varepsilon}$ 来自胶子极化平均。这一道只有 $\hat t\to0$（$v\to1$）的共线极点，没有软奇异性。

\emph{第四步：虚修正。}单圈虚修正只有顶点修正（\cref{fig:dy-oneloop}(c)）：无质量外线的自能是无标度积分，在维数正规化中为零，$Z_1=Z_2$，耦合的重整化在此阶不出现。\cref{prop:sudakov-oneloop}（\tagR，$R_{\mathrm{PT}}$）给出精确到 $\varepsilon$ 的所有阶的形状因子
\[
F=1+\frac{\as}{4\pi}C_F\Big(\frac{\mu^2}{-q^2-\ii0}\Big)^{\varepsilon}f(\varepsilon),\qquad
f(\varepsilon)=-K_v(\varepsilon)\,\frac{2-\varepsilon+2\varepsilon^2}{\varepsilon^2(1-2\varepsilon)},
\]
其中 $K_v(\varepsilon)\equiv\ee^{\varepsilon\gamma_E}\Gamma(1+\varepsilon)\Gamma^2(1-\varepsilon)/\Gamma(1-2\varepsilon)$。
对 DY 的类时 $q^2=Q^2>0$，$(-q^2-\ii0)^{-\varepsilon}=(Q^2)^{-\varepsilon}\ee^{\ii\pi\varepsilon}$。对无质量在壳夸克，顶点修正向 Born 结构的投影是精确的；与 Born 振幅的干涉给出
\begin{equation}\label{eq:dy-virt}
\Delta^V=2\Real F^{(1)}\,\delta(1-z)=a\,C_F\Big(\frac{\mu^2}{Q^2}\Big)^{\varepsilon}\cos(\pi\varepsilon)\,f(\varepsilon)\,\delta(1-z).
\end{equation}
% VF-DY-6
虚部 $\Imag F^{(1)}=\frac{\as C_F}{4\pi}\big(\frac{\mu^2}{Q^2}\big)^{\varepsilon}\sin(\pi\varepsilon)f(\varepsilon)=-\frac{\as C_F}{2}\big(\frac1\varepsilon+L+\frac32\big)+\order{\varepsilon}$ 是入射 $q\bar q$ 对之间红外发散的 Coulomb 相位；它不进入此阶的 $2\Real F$，到 $\order{\as^2}$ 才通过 $|F|^2$ 进入截面。

\emph{第五步：按分布展开。}由\cref{lem:prelim-plus-expansion}（取 $a=2$），在 $[0,1]$ 上
\[
(1-z)^{-1-2\varepsilon}=-\frac{\delta(1-z)}{2\varepsilon}+\sum_{k\ge0}\frac{(-2\varepsilon)^k}{k!}\,\mathcal D_k(z);
\]
又 $K_r=1-\frac32\zeta_2\varepsilon^2+\order{\varepsilon^3}$，$K_v=1-\frac12\zeta_2\varepsilon^2+\order{\varepsilon^3}$，$\cos(\pi\varepsilon)=1-3\zeta_2\varepsilon^2+\order{\varepsilon^4}$，$f=-\frac2{\varepsilon^2}-\frac3\varepsilon-8+\zeta_2+\order{\varepsilon}$。把\cref{eq:dy-real-qq,eq:dy-real-qg,eq:dy-virt}写成 $\delta(1-z)$、$\mathcal D_k$ 与正规函数的组合，即得\cref{tab:dy-laurent}。以实辐射 $q\bar q$ 项为例：$\delta(1-z)$ 只来自展开式的第一项，在 $z=1$ 处 $z^\varepsilon$ 与方括号合为 $2$，给出 $\frac{2C_F}{\varepsilon^2}(\frac{\mu^2}{Q^2})^\varepsilon K_r=\frac{2C_F}{\varepsilon^2}+\frac{2C_FL}{\varepsilon}+C_F(L^2-3\zeta_2)+\order{\varepsilon}$，而 $3\zeta_2=\pi^2/2$；$\mathcal D_0$ 项用 $(1+z^2)\mathcal D_0=2\mathcal D_0-(1+z)$；$\order{\varepsilon^0}$ 的 $\mathcal D_1$ 项来自 $-\frac{2C_F}{\varepsilon}(-2\varepsilon)(1+z^2)\mathcal D_1=4C_F(1+z^2)\mathcal D_1$；$z^\varepsilon=1+\varepsilon\ln z+\dots$ 给出正规的 $-2C_F\frac{1+z^2}{1-z}\ln z$。

\begin{table}[htbp]
\centering
\caption{$\order{\as}$ 各部分的 Laurent 展开（单位 $a=\as/2\pi$；$L=\ln(\mu^2/Q^2)$，$\delta\equiv\delta(1-z)$）}
\label{tab:dy-laurent}
\footnotesize
\setlength{\tabcolsep}{4pt}
\begin{tabular}{@{}llll@{}}
\toprule
 & $\varepsilon^{-2}$ & $\varepsilon^{-1}$ & $\varepsilon^0$\\
\midrule
实辐射 $q\bar q$ & $2C_F\delta$ & $2C_FL\delta-4C_F\mathcal D_0+2C_F(1+z)$ & $C_F(L^2-\frac{\pi^2}{2})\delta-4C_FL\mathcal D_0+8C_F\mathcal D_1+2C_F(1+z)L$\\
 & & & $\ -4C_F(1+z)\ln(1-z)-2C_F\frac{1+z^2}{1-z}\ln z$\\
虚修正 & $-2C_F\delta$ & $-C_F(3+2L)\delta$ & $C_F\big(-L^2-3L-8+\frac{7\pi^2}{6}\big)\delta$\\
实辐射 $qg$ & --- & $-P_{qg}$ & $T_F\big[(z^2+(1-z)^2)\ln\frac{(1-z)^2}{z}+\frac{1+6z-7z^2}{2}\big]-P_{qg}L$\\
\bottomrule
\end{tabular}
\end{table}
% VF-DY-7

\emph{第六步：求和与质量因子化。}$q\bar q$ 道实、虚之和的 $\varepsilon^{-2}$ 项相消（软双极点）；$\varepsilon^{-1}$ 项为
\[
2C_FL\delta-4C_F\mathcal D_0+2C_F(1+z)-C_F(3+2L)\delta=-2C_F\Big[2\mathcal D_0-(1+z)+\frac32\delta\Big]=-2P_{qq},
\]
即剩下的极点恰为 $-\frac2\varepsilon P_{qq}$，每条入射部分子线一个共线极点；它与\cref{eq:dy-mf}的抵消项
\[
\frac2\varepsilon\Big(\frac{\mu^2}{\mu_F^2}\Big)^{\varepsilon}P_{qq}=\frac2\varepsilon P_{qq}+2\Big(L+\ln\frac{Q^2}{\mu_F^2}\Big)P_{qq}+\order{\varepsilon}
\]
相消。有限部分中，$\delta$ 的系数为 $C_F(L^2-\frac{\pi^2}2)+C_F(-L^2-3L-8+\frac{7\pi^2}6)+3C_FL=C_F(\frac{2\pi^2}3-8)$；$L\mathcal D_0$ 与 $(1+z)L$ 项在实辐射与抵消项之间相消；剩下的 $8C_F\mathcal D_1-4C_F(1+z)\ln(1-z)-2C_F\frac{1+z^2}{1-z}\ln z+2P_{qq}\ln\frac{Q^2}{\mu_F^2}$ 就是\cref{eq:dy-Dqq}。$qg$ 道的极点 $-\frac1\varepsilon P_{qg}$ 与抵消项 $\frac1\varepsilon(\frac{\mu^2}{\mu_F^2})^\varepsilon P_{qg}$ 相消，$L$ 项相消，剩下\cref{eq:dy-Dqg}。两道的结果都不含 $L$，故此阶没有 $\mu_R$ 依赖。
\end{proof}

结果中有两点值得注意。其一，$\delta$ 项的 $\frac{2\pi^2}3=-\frac{\pi^2}2+\frac{\pi^2}6+\pi^2$，三项分别来自实辐射的 $K_r$、虚修正的 $K_v$ 与类时延拓的 $\cos\pi\varepsilon$；最后一项是单圈时类时与类空形状因子之比 $|F(q^2)/F(-q^2)|^2-1=aC_F\pi^2$，即“$\pi^2$ 增强”。其二，$\Delta_{qg}$ 的有限部分依赖于胶子极化平均的约定（\cref{ex:dy-nlo}(c)）。结果与 Duhr、Dulat、Mistlberger 附件中的 NLO 项（换算其展开参数 $\as/\pi$ 后）逐项一致，包括 $\mu_F$ 项\cite{Duhr:2020sdp}；形状因子与\cite{MochVermaserenVogt2005}一致到 $\varepsilon^4$，阈值常数与\cite{Moch:2005ky}一致（\cref{ex:dy-threshold}）。早期的 $\order{\as}$ 计算见\cite{AltarelliEllisMartinelli1979}与\cite{KubarAndrePaige1979}：前者用由轻子产生定出的部分子分布表示截面，即 DIS 方案（\cref{rem:dy-dis-scheme}）；后者指出红外与质量奇异性可以一致地吸收进夸克分布（均据其摘要）；上面的 $\MSbar$ 形式是按现代约定整理的，不应直接归于这两篇文章。完整的 $\order{\as^2}$ 结果见\cite{HambergVanNeervenMatsuura1991}，$\order{\as^3}$ 见\cite{Duhr:2020seh}（中性流）与\cite{Duhr:2020sdp}（带电流）。

\begin{example}[$\order{\as}$ 系数函数的完整推导]\label{ex:dy-nlo}
\cref{prop:dy-nlo}的证明写出了推导的每一步；这里把它落实为可以逐项手算或直接数值重复的检验，并指出三个改变有限部分的约定。取 $\mu_F=Q$，$C_F=\frac43$，$T_F=\frac12$。

\emph{(a) 逐项的一阶矩。}对试探函数 $h\equiv1$，乘积约定给出
\[
\big\langle4(1+z^2)\mathcal D_1,1\big\rangle=4\int_0^1(z^2-1)\frac{\ln(1-z)}{1-z}\dd z=-4\int_0^1(1+z)\ln(1-z)\dd z=7;
\]由 $\int_0^1\frac{\ln z}{1-z}\dd z=-\zeta_2$ 与 $\int_0^1\frac{z^2\ln z}{1-z}\dd z=\frac54-\zeta_2$，$-2\int_0^1\frac{1+z^2}{1-z}\ln z\,\dd z=4\zeta_2-\frac52$；$\delta$ 项贡献 $4\zeta_2-8$。三项相加，
\[
\int_0^1\Delta^{(1)}_{q\bar q}\,\dd z=C_F\Big(8\zeta_2-\frac72\Big)=12.879297\ldots,
\]
$qg$ 道三项为 $2\int_0^1(z^2+(1-z)^2)\ln(1-z)\dd z=-\frac{13}9$、$-\int_0^1(z^2+(1-z)^2)\ln z\,\dd z=\frac{13}{18}$、$\int_0^1(\frac12+3z-\frac72z^2)\dd z=\frac56$，和为 $\frac{T_F}9$。同样，二阶矩为 $C_F(8\zeta_2+\frac{13}{18})$ 与 $-\frac{11}{18}T_F$。
% VF-DY-14

\emph{(b) 在有限 $\varepsilon$ 上的数值核对。}对 $\varepsilon<0$（$D>4$），\cref{eq:dy-real-qq,eq:dy-real-qg,eq:dy-virt}与抵消项作用在光滑试探函数上的积分全部收敛，可以直接计算。取 $h(z)=z^3$、$L=0.37$、$\ln(Q^2/\mu_F^2)=0.21$，在 $\varepsilon=-0.001,\dots,-0.008$ 上算出 $\langle\Delta^{{\rm bare},(1)}+\text{抵消项},h\rangle$ 并拟合 Laurent 级数：$\varepsilon^{-2}$、$\varepsilon^{-1}$ 的系数小于 $2\times10^{-14}$，有限部分 $28.279889305624$ 与\cref{eq:dy-Dqq}的 $28.279889305640$ 相差 $2\times10^{-11}$；$qg$ 道同样一致到 $10^{-12}$。有限部分与 $L$ 无关，正是“此阶没有 $\mu_R$ 依赖”的数值表现。

\emph{(c) 改变有限部分的三个约定。}(i) 若把加号作用于整个 $\frac{(1+z^2)\ln(1-z)}{1-z}$，由 $-\int_0^1(1+z)\ln(1-z)\dd z=\frac74$ 得 $(1+z^2)\mathcal D_1=\big[\frac{(1+z^2)\ln(1-z)}{1-z}\big]_++\frac74\delta(1-z)$，$\delta$ 项的系数变为 $C_F(\frac{2\pi^2}3-1)$，差 $7C_F$。(ii) 若胶子极化按 $\frac12$ 而不按 $\frac1{2(1-\varepsilon)}$ 平均，$\Delta_{qg}\to\Delta_{qg}+P_{qg}$，这是另一个方案。(iii) 若裸系数按四维 Born 项归一化（不除以 $1-\varepsilon$）而抵消项仍取 $\frac2\varepsilon P_{qq}$，则因 $(1-\varepsilon)(-\frac2\varepsilon P_{qq})=-\frac2\varepsilon P_{qq}+2P_{qq}$，$\Delta_{q\bar q}$ 会多出 $2P_{qq}$。
% VF-DY-12, VF-CONV-6
\uses{prop:dy-nlo}
\end{example}

\begin{proposition}[$\mu_F$ 无关性 \tagR]\label{prop:dy-scale}
（$R_{\mathrm{PT}}$：在重整化微扰论中对红外正规化的部分子态严格。）本命题是 $\order{\as}$ 的显式陈述。设部分子分布 $f_i(x;\mu_F)$ 满足 DGLAP 方程（\cref{thm:pdf-dglap}，\tagR，$R_{\mathrm{PT}}$），则以\cref{prop:dy-nlo}的系数构成的\cref{eq:dy-fact-nlo}满足
\[
\frac{\dd}{\dd\ln\mu_F^2}\,\frac{\dd\sigma}{\dd Q^2}=\order{\as^2}.
\]
% VF-DY-11
非微扰地位：\cref{asm:intro-composite}（MA0；用于强子态的 DGLAP 方程，部分子态无需）。
\uses{prop:dy-nlo,thm:pdf-dglap,lem:prelim-mellin,asm:intro-composite}
\end{proposition}

\begin{proof}
固定 $Q^2$ 与 $\mu_F$，把\cref{eq:dy-fact-nlo}中 $s/\sigma_0$ 乘以截面的那部分看作 $\tau$ 的函数 $\mathcal F(\tau)$，取 Mellin 矩 $g_N\equiv\int_0^1x^{N-1}g(x)\dd x$。由\cref{lem:prelim-mellin}，$\int_0^1\dd\tau\,\tau^{N-1}\int\dd x_1\dd x_2\dd z\,\delta(x_1x_2z-\tau)f(x_1)g(x_2)h(z)=f_Ng_Nh_N$，故
\[
\mathcal F_N=\sum_qe_q^2\Big\{\big(q_N^A\bar q_N^B+\bar q_N^Aq_N^B\big)\Delta_{q\bar q,N}+\big[g_N^A(q_N^B+\bar q_N^B)+(q_N^A+\bar q_N^A)g_N^B\big]\Delta_{qg,N}\Big\}.
\]
记 $\ell_F\equiv\ln(\mu_F^2/Q^2)$ 与 $\gamma_{ij}(N)=\int_0^1z^{N-1}P_{ij}(z)\dd z$。由\cref{prop:dy-nlo}，$\partial_{\ell_F}\Delta_{q\bar q,N}=-2a\gamma_{qq}+\order{a^2}$，$\partial_{\ell_F}\Delta_{qg,N}=-a\gamma_{qg}+\order{a^2}$；由 DGLAP 方程，$\partial_{\ell_F}q_N=a(\gamma_{qq}q_N+\gamma_{qg}g_N)+\order{a^2}$（$\bar q_N$ 同，$\gamma_{\bar q\bar q}=\gamma_{qq}$，$\gamma_{\bar qg}=\gamma_{qg}$），而 $\partial_{\ell_F}g_N=\order{a}$ 只在乘以 $\Delta_{qg,N}=\order{a}$ 时出现。于是
\begin{multline*}
\partial_{\ell_F}\mathcal F_N=a\sum_qe_q^2\Big\{(\gamma_{qq}q^A+\gamma_{qg}g^A)\bar q^B+q^A(\gamma_{qq}\bar q^B+\gamma_{qg}g^B)\\
+(\gamma_{qq}\bar q^A+\gamma_{qg}g^A)q^B+\bar q^A(\gamma_{qq}q^B+\gamma_{qg}g^B)\\
-2\gamma_{qq}\big(q^A\bar q^B+\bar q^Aq^B\big)-\gamma_{qg}\big[g^A(q^B+\bar q^B)+(q^A+\bar q^A)g^B\big]\Big\}+\order{a^2}=\order{a^2}
\end{multline*}
（指标 $N$ 省去）：$\gamma_{qq}$ 项在 $q\bar q$ 道内相消；$\gamma_{qg}$ 项（夸克或反夸克由胶子产生）与 $qg$ 道相消。所有矩在 $\order{a}$ 都与 $\mu_F$ 无关，由 Mellin 反演的唯一性（\cref{lem:prelim-mellin}），$\mathcal F(\tau)$ 亦然。
\end{proof}

在 $x$ 空间用试探分布直接计算，这一相消也成立，残差为 $3\times10^{-13}$。应当强调\cref{prop:dy-scale}说了什么：残余的 $\mu_F$ 依赖属于 $\order{\as^2}$；$\mu_R$ 依赖在 $\order{\as^2}$ 才出现，正比于 $\beta_0\ln(\mu_R^2/\mu_F^2)\,\Delta^{(1)}$。实践中常把 $\mu_F$、$\mu_R$ 在 $Q/2$ 与 $2Q$ 之间变化所得的变化当作“理论误差”，这只是一种约定：因子化只保证残余依赖属于更高阶，并不给出被截掉的高阶项的大小。

\begin{remark}[DIS 方案中的 DY 与 PDF 的普适性]\label{rem:dy-dis-scheme}
DIS 方案用 $F_2$ 定义夸克分布：$q^{\rm DIS}=q+a\,[C_2^q\otimes q+C_2^g\otimes g]$，其中 $C_2^{q,g}$ 是 $\MSbar$ 的 DIS 系数函数（\cref{prop:disall-nlo}，\tagR，$R_{\mathrm{PT}}$；归一化为 $F_2/x=\sum_qe_q^2\{q+\bar q+a[C_2^q\otimes(q+\bar q)+2C_2^g\otimes g]\}$；方案变换见\cref{prop:disall-dis-scheme}）。把\cref{eq:dy-fact-nlo}中的 $\MSbar$ 分布换成 DIS 方案的分布，$\order{a}$ 的系数变为（取 $\mu_F=Q$）
\[
\begin{aligned}
\Delta^{\rm DIS}_{q\bar q}&=C_F\Big[\Big(1+\frac{4\pi^2}3\Big)\delta(1-z)+2(1+z^2)\mathcal D_1+3\mathcal D_0-6-4z\Big],\\
\Delta^{\rm DIS}_{qg}&=T_F\Big[\big(z^2+(1-z)^2\big)\ln(1-z)+\frac32-5z+\frac92z^2\Big],
\end{aligned}
\]
% VF-DY-10
并且精确地有 $\Delta^{\MSbar}_{q\bar q}-2C_2^q=\Delta^{\rm DIS}_{q\bar q}$、$\Delta^{\MSbar}_{qg}-C_2^g=\Delta^{\rm DIS}_{qg}$（我们用符号运算与 50 位精度的数值作了核对；这正是常引的 DIS 方案形式）。这两个恒等式是部分子分布普适性在单圈上的显式检验：DY 的两条入射线与 DIS 的一条入射线的共线发散由同一组分布吸收，换方案时 DY 系数恰好为每条入射线改变一个 $C_2$。比较文献中的公式时还须注意加号分布的约定（\cref{ex:dy-nlo}(c)(i)，差 $7C_F\delta(1-z)$）；这只是约定的说明。
\end{remark}

\begin{example}[大 $N$ 极限]\label{ex:dy-threshold}
取 $\mu_F=Q$ 与 Mellin 矩 $\Delta_{q\bar q,N}=\int_0^1z^{N-1}\Delta_{q\bar q}(z)\dd z$。由\cref{lem:prelim-mellin}，$\langle\mathcal D_0,z^{N-1}\rangle=-S_1(N-1)$，$\langle\mathcal D_1,z^{N-1}\rangle=\frac12\big[S_1(N-1)^2+S_2(N-1)\big]$，其中 $S_k(n)=\sum_{j=1}^nj^{-k}$；又 $\int_0^1z^{N-1}\frac{\ln z}{1-z}\dd z=-\psi'(N)=\order{1/N}$。乘积约定给出 $\langle(1+z^2)\mathcal D_1,z^{N-1}\rangle=\langle\mathcal D_1,z^{N-1}\rangle+\langle\mathcal D_1,z^{N+1}\rangle$，于是
\[
\langle4(1+z^2)\mathcal D_1,z^{N-1}\rangle=2\big[S_1(N-1)^2+S_2(N-1)+S_1(N+1)^2+S_2(N+1)\big].
\]
用 $S_1(N)=\ln\bar N+\order{1/N}$（$\bar N=N\ee^{\gamma_E}$）与 $S_2(N)=\zeta_2+\order{1/N}$，并加上 $\delta$ 项的 $4\zeta_2-8$，
\[
\Delta_{q\bar q,N}=1+a\,C_F\big[4\ln^2\bar N+8\zeta_2-8\big]+\order{\frac{\ln N}{N}} .
\]
% VF-DY-14
数值上（$C_F=\frac43$），$N=10,\,100,\,1000,\,10^4$ 时精确的 $\Delta^{(1)}_{q\bar q,N}$ 为 $52.540$、$150.391$、$305.718$、$517.798$，渐近式给出 $51.110$、$150.117$、$305.678$、$517.793$；差值的领头渐近形式是 $4C_F\ln\bar N/N$（来自 $S_1(N-1)=\ln\bar N-\frac1{2N}+\dots$ 与 $S_1(N+1)=\ln\bar N+\frac3{2N}+\dots$），二者之比依次为 $0.931$、$0.991$、$0.9990$、$0.9999$，符合 $\order{\ln N/N}$。$qg$ 道不含 $\mathcal D_k$，$\Delta_{qg,N}=\order{\ln N/N}$。双对数 $4\ln^2\bar N$ 来自阈值 $z\to1$ 处的软胶子辐射，在 $\as^n$ 阶为 $\ln^{2n}\bar N$，其重求和见第~13~章（\cref{prop:thr-oneloop}）。按 $a_s=\as/4\pi$ 换算，$2\Delta^{(1)}_{q\bar q,N}=C_F\big[8\ln^2N+16\gamma_E\ln N-16+16\zeta_2+8\gamma_E^2\big]$，常数项与\cite{Moch:2005ky}的阈值常数一致。
\uses{prop:dy-nlo,lem:prelim-mellin}
\end{example}

% =====================================================================
\section{其他强子—强子过程}\label{sec:dy-other}
% =====================================================================

\cref{thm:dy-factorization}的论证用到了被测末态的三个性质：(i) 被测系统在硬散射中产生，本身是色单态，其衰变产物不参与强相互作用；(ii) 截面对其余末态（旁观者、未观测的喷注、软辐射）完全单举；(iii) 测量不约束旁观者的末态，在割线两侧作为恒等算符出现（\cref{rem:dy-one-time-per-jet}）。凡具有这三个性质的过程都继承 DY 的地位；被测末态带色的过程在 (iii) 处失效，其因子化至今只是猜想。下面分别陈述。

\begin{theorem}[色单态末态 \tagA]\label{thm:dy-colour-singlet}
设被测的只是一个在硬散射中产生的色单态系统 $V$：$W^\pm$、$Z$（含其轻子衰变）；Higgs 玻色子（经胶子融合，用重顶夸克极限下的有效顶点或保留顶夸克圈的完整理论，或经其他色单态道产生），其衰变产物不带色；电弱规范玻色子对 $W^+W^-$、$ZZ$、$W^\pm Z$（轻子衰变）。对其余末态单举，则 $V$ 的质量与快度分布具有\cref{thm:dy-factorization}的形式：部分子分布与 DY、DIS 中的相同，只是硬函数 $\mathcal H_{ij}$ 换成相应过程的硬函数，余项同样是 $\order{(\Lambda/Q)^p}$，$Q$ 为硬标度。\\
非微扰地位：\cref{asm:intro-bridge}（MA1、MA2）。
% errtype: E-S
\uses{thm:dy-factorization,prop:dy-regions,lem:dy-fsi,prop:dy-deformation,def:pdf-gluon,prop:appI-larin,asm:intro-bridge}
\end{theorem}

\begin{proofsketch}
(1) \cref{thm:dy-factorization}的证明中，虚光子只以“在硬顶点 $V_H$ 处产生的色单态系统”出现：区域分析、割图求和、单侧改写与形变（\cref{prop:dy-regions,lem:dy-fsi,prop:dy-deformation}）都只涉及两个强子的共线子图、软子图与 Glauber 交换，不涉及 $V$ 本身；所测的 $V$ 的变量不约束强子末态。(2) 对 $W$、$Z$ 与规范玻色子对，改变的只是 $H$ 中的电弱顶点；对无质量夸克，$\order{\as}$ 的 QCD 修正与虚光子相同，只需替换电弱耦合（维数正规化中轴矢流的 $\gamma_5$ 按\cref{prop:appI-larin}（\tagR，$R_{\mathrm{PT}}$）处理）。(3) 对胶子融合产生 Higgs，活动部分子是胶子，共线因子是胶子分布（\cref{def:pdf-gluon}），其 Wilson 线在伴随表示中；\cref{app:appO}的分析只用到活动部分子无质量、强子为色单态，与活动部分子是夸克还是胶子无关。重顶夸克有效顶点的匹配是一个额外的、逐阶可算的硬步骤。完整论证与 DY 相同\cite[第~14.4~节]{Collins2011}；硬函数见\cite{AltarelliEllisMartinelli1979,HambergVanNeervenMatsuura1991}（DY 与 $W$、$Z$）及\cite{Dawson1991,DjouadiSpiraZerwas1991,HarlanderKilgore2002,AnastasiouEtAl2015}（胶子融合，NLO 到 N$^3$LO）。
\end{proofsketch}

\paragraph{地位与缺口}\tagA。与\cref{thm:dy-factorization}完全相同：Glauber 区的输入\cref{asm:appO-mixing}(c)、\cref{asm:appO-deformation}(b)--(d)、\cref{asm:appO-rewritten-cs}等，附录~N 的一般输入，具名缺口\cref{asm:appN-brst-nonlocal}（\tagC，见\cref{thm:method-template}），余项幂次 $p$ 不由证明确定（\tagC）。没有额外的非微扰输入，但要求 $V$ 的衰变产物不带色；若测量带色的衰变产物（例如 $H\to b\bar b$ 的 $b$ 喷注），就落入\cref{conj:dy-coloured-final}的情形。对 $V$ 的横动量 $q_T\ll Q$ 的微分截面，适用的是 TMD 因子化（\cref{thm:tmd-factorization}），软因子不再为 $1$。

\begin{theorem}[极化 Drell--Yan \tagA]\label{thm:dy-polarized}
对纵向极化（螺旋度）与横向极化（横自旋）强子的 DY，自旋相关的截面差（例如两个强子螺旋度相同与相反之差）具有\cref{thm:dy-factorization}的因子化形式，其中的部分子分布换成螺旋度分布与横度分布（\cref{def:pdf-polarized}），硬函数换成相应的自旋投影硬函数。\\
非微扰地位：\cref{asm:intro-bridge}（MA1、MA2、MA4）。
% errtype: E-S
\uses{thm:dy-factorization,def:pdf-polarized,def:appI-hvbm,asm:intro-bridge}
\end{theorem}

\begin{proofsketch}
把\cref{thm:dy-factorization}的证明逐步重复，只在共线因子与硬部分的连接处插入自旋投影：纵向极化取 $\gamma^+\gamma_5$ 投影，得到螺旋度分布；横向极化取 $\gamma^+\gamma^j_\perp\gamma_5$ 型投影，得到横度分布。Glauber 区与软胶子的处理与自旋无关，因为领头幂次的软与 Glauber 耦合只经过共线线的大分量，不改变其自旋结构。领头阶时，矢量流保持手征，入射夸克与反夸克的螺旋度相反，部分子层次的双螺旋度不对称为 $-1$。横度分布是手征奇的，在 DIS 中不出现于领头扭度，在两个横向极化强子的 DY 中则以两个手征奇分布之积的形式出现\cite{RalstonSoper1979,JaffeJi1991}。极化硬函数的计算需要维数正规化中 $\gamma_5$ 的处理（\cref{def:appI-hvbm}）。
\end{proofsketch}

\paragraph{地位与缺口}\tagA。与\cref{thm:dy-factorization}相同；此外极化部分子分布的 $\MSbar$ 定义依赖 $\gamma_5$ 方案（HVBM 加 Larin 有限重整化），这是方案约定而不是缺口。

\begin{conjecture}[强子对撞中有色末态的共线因子化 \tagC]\label{conj:dy-coloured-final}
单举的大横动量强子与喷注以及重夸克对的强子—强子产生满足标准的共线因子化。例如，对 $A+B\to h(p_h)+X$（$p_{h,T}\sim\sqrt s$），
\[
E_h\frac{\dd\sigma}{\dd^3p_h}=\sum_{i,j,k}\int\dd\xi_A\,\dd\xi_B\,\frac{\dd z}{z^2}\,f_{i/A}(\xi_A;\mu)\,f_{j/B}(\xi_B;\mu)\,d_{h/k}(z;\mu)\,E_k\frac{\dd\hat\sigma_{ij\to k}}{\dd^3p_k}+\order{\Big(\frac{\Lambda}{p_{h,T}}\Big)^{p}},
\]
其中 $d_{h/k}$ 是碎裂函数（\cref{def:pdf-ff}），$p_h=zp_k$，因子 $z^{-2}$ 来自 $\dd^3p_h/E_h=z^2\,\dd^3p_k/E_k$；对红外安全的喷注截面与重夸克对总截面，类似地有 $\sigma=\sum_{ij}f_i\otimes f_j\otimes\hat\sigma_{ij}$。
\uses{thm:dy-factorization,def:pdf-ff}
\end{conjecture}

\paragraph{地位与缺口}\tagC。这些公式被普遍使用，但没有与 DY 同等细节的全阶证明。Collins 的教科书写道，这些情形中存在“relative lack of detailed proofs”，并把 Nayak、Qiu 与 Sterman 的分析\cite{NayakQiuSterman2005b}列为例外\cite[第~14.7~节，第~570--571~页]{Collins2011}；Diehl 等人指出，对末态带色的单次散射过程，“to the best of our knowledge the absence of Glauber gluon effects has not yet been proven”\cite[第~5~节]{Diehl:2015bca}。对重夸克产生，CSS 从低阶图论证了标准公式，并强调需要所有阶的分析才能可靠地确立这一猜想\cite{Collins:1985gm}（据其摘要）；Aybat 与 Sterman 给出了有色末态软与 Glauber 相消的纲要\cite{AybatSterman2009}；Catani、de~Florian 与 Rodrigo 发现类空共线极限中严格共线因子化在振幅层次被破坏，并指出它可能从 N$^3$LO 起影响质量奇异性的因子化\cite[第~8~节]{CataniDeFlorianRodrigo2012}。在\cref{app:appO}的框架中，论证在单侧改写一步（(O-3)）失效：被测的有色末态使测量投影不与光锥演化对易，相关的钳夹区是相对于方向对 $(n_A,w_1)$ 的 Glauber 区（\cref{prop:break-measurement}）。另一方面，NLO 的重夸克\cite{NasonDawsonEllis1988}与喷注\cite{Kang:2016mcy}等固定阶计算都以这种形式组织，与数据的符合是证据而不是证明；对同一类过程的 TMD 版本（背对背强子），因子化不成立（第~18~章\cref{thm:break-tmd}），但共线的、对 $q_T$ 积分的截面没有已知的反例；类空共线极限中严格共线因子化的失效（\cref{prop:break-spacelike}）与超出 NNLO 的质量奇异性因子化（\cref{prop:break-mass-singularity}）也在第~18~章讨论。要把本猜想提升为 \tagA，需要对有色末态完成 (O-3) 的替代论证。

\begin{conjecture}[极化强子对撞中的 $A_{LL}$（喷注、$\pi^0$） \tagC]\label{conj:dy-polarized-jets}
纵向极化强子对撞中喷注与 $\pi^0$ 产生的双自旋不对称 $A_{LL}$ 以螺旋度分布按\cref{conj:dy-coloured-final}的形式因子化。
\uses{conj:dy-coloured-final,thm:dy-polarized}
\end{conjecture}

\paragraph{地位与缺口}\tagC。继承\cref{conj:dy-coloured-final}：末态带色，Glauber 相消没有证明；自旋投影本身不引入新的困难（\cref{thm:dy-polarized}）。

% =====================================================================
\section{高扭度：\texorpdfstring{$1/Q^2$}{1/Q2} 处的广义因子化及其在更高幂次的失效}\label{sec:dy-twist}
% =====================================================================

\cref{thm:dy-factorization}的余项 $\order{(\Lambda/Q)^p}$ 中 $p$ 不由证明确定。对 DIS，幂次修正由 OPE 的高扭度算符给出（\cref{prop:dis-twist4}，\tagR，$R_{\mathrm{OPE}}$）；DY 不是 OPE 可及的量，幂次修正只能在动量空间中逐区域分析。已知的情形是：相对领头幂次的第一阶修正（$1/Q^2$，扭度 4）仍有广义的因子化形式，第二阶（$m^4/s^2$，扭度 6）则没有。

\begin{theorem}[$1/Q^2$ 处的广义因子化 \tagA]\label{thm:dy-higher-twist}
DY 截面的领头幂次修正是相对 $1/Q^2$ 的扭度 4 项，具有广义因子化形式：它们是一个可计算的短程硬部分、一个强子的扭度 2 矩阵元（部分子分布）与另一个强子的扭度 4 多部分子关联矩阵元的卷积，
\[
\frac{\dd\sigma}{\dd Q^2}=\frac{\dd\sigma^{(\rm tw\,2)}}{\dd Q^2}+\frac{1}{Q^2}\sum_{i,j}\Big[f_{i/A}\otimes T^{(4)}_{j/B}\otimes H^{(4)}_{ij}+(A\leftrightarrow B)\Big]+\order{\frac{1}{Q^4}},
\]
其中卷积是示意的写法（$T^{(4)}$ 依赖于多个动量分数，卷积作用于全部动量分数），$T^{(4)}_{j/B}$ 是含四个场（例如一对夸克场与一对场强，或两对夸克场）的光锥算符在强子 $B$ 中的矩阵元\cite{QiuSterman1991a,QiuSterman1991b}。同样的结构用于原子核中的双散射\cite{Kastella:1989ux}。\\
非微扰地位：\cref{asm:intro-bridge}（MA1、MA2）。
% errtype: E-S
\uses{thm:dy-factorization,asm:intro-bridge}
\end{theorem}

\begin{proofsketch}
(1) 在次领头幂次上，额外的一对共线部分子（夸克或物理极化的胶子）从一个强子的共线子图连到硬部分，按幂次计数每多一对代价为 $\Lambda^2/Q^2$；纵向胶子仍由 Ward 恒等式收成 Wilson 线，使多部分子关联函数规范不变。(2) 末态相互作用的幺正性相消与软胶子的处理在这一幂次上重复领头幂次的论证；在标量夸克模型中，单圈软胶子效应在 $m^2/s$ 处确实因子化\cite{Basu:1984ba}。(3) 由此得到式中的三重卷积，硬部分 $H^{(4)}$ 可逐阶计算。Qiu 与 Sterman 在 $\as$ 的领头阶算出了这些修正，并论证了这类 $1/Q^2$ 项可以由广义因子化定理处理\cite{QiuSterman1991a,QiuSterman1991b}（据其摘要）。（各幂次计数步骤只给出 E-S 型陈述。）
\end{proofsketch}

\paragraph{地位与缺口}\tagA。扭度 4 的全阶论证见\cite{QiuSterman1991b}；\cref{app:appO}的 Glauber 分析只在领头幂次进行，次领头幂次上的单侧改写与形变没有写成条件定理，这是本定理比\cref{thm:dy-factorization}更弱的地方。CSS 的综述指出，微扰论在 $m^2/s$ 处遵守因子化，这使得因子化定理在这一幂次上对强子—强子散射可能成立\cite[第~10.2~节]{CollinsSoperSterman1989}。还要注意，领头扭度微扰级数的红外重整化子不确定性与 $1/Q^2$ 修正同阶（\cref{asm:intro-bridge}(Br3)；对 DY 见\cite{BenekeBraun1995b}），因此扭度 4 矩阵元只有在规定了微扰级数的求和方式之后才有确定的意义。

\begin{proposition}[更高幂次的因子化失效 \tagV]\label{prop:dy-twist6-failure}
在 $m^4/s^2$（相对领头幂次的第二阶幂次修正，扭度 6）上，Drell--Yan 截面不具有广义因子化形式：在微扰论中，这一幂次上存在既不在实、虚修正之间相消，也不能吸收进强子矩阵元的红外敏感贡献；而在 $m^2/s$ 上因子化仍成立（\cref{thm:dy-higher-twist}）。
\uses{thm:dy-higher-twist}
\end{proposition}

\paragraph{地位与缺口}\tagV。证据有三类。(i) 模型中的显式微扰反例：在 Drell--Yan 过程的标量夸克 QCD 模型中，单圈软胶子效应在 $m^2/s$ 处因子化，在 $m^4/s^2$ 处不因子化，而且非因子化项不被形状因子压低\cite{Basu:1984ba}（据其摘要）。(ii) QCD 中的显式计算：CSS 的综述总结道，在高扭度上强子—强子散射的红外发散超出单圈即不相消，在 Drell--Yan 中两圈的非相消发散出现在 $m^4/s^2$ 量级\cite[第~10.2~节]{CollinsSoperSterman1989}（另见\cref{rem:dy-dft-premises}）。(iii) Bodwin、Braaten 与 Lepage 的总结：在 $1/Q^2$ 的次领头阶之外，“factorization is known to fail”\cite[第~VI~节]{BodwinBraatenLepage1995}。这些都是微扰论的陈述。CSS 同时强调，如何解释这种不相消“is not quite clear”，包括非微扰效应在内的实际情况可能比微扰计算所显示的更好或更坏\cite[第~10.2~节]{CollinsSoperSterman1989}；因此 \tagV 指的是微扰论中广义因子化形式的失效，它对强子数据的数值影响并未确定。

% =====================================================================
\section{本章小结}\label{sec:dy-summary}
% =====================================================================

本章的结论按严格性分级为：Born 截面、$\order{\as}$ 系数函数与 $\mu_F$ 无关性、单个割图中的 Glauber 钳夹、中间态取向决定 $k^-$ 极点、程函线间双胶子交换的 $\order{g^4}$ 幺正性在重整化微扰论中严格，单个割图给出指向未来的 Wilson 线与 DY、DIS 部分子分布的同一性由定义直接推出（\tagR）；旁观者 Glauber 相位的相消在阿贝尔程函模型中完整证明（\tagP）；区域、幺正性相消、同时形变、软因子为 $1$ 以及单举 DY、色单态末态、极化 DY 的共线因子化与 $1/Q^2$ 处的广义因子化给出全阶论证（\tagA），其精确形式是附录中以编号假设为前提的条件定理（\cref{tab:dy-correspondence}）；有色末态的因子化是猜想（\tagC），$m^4/s^2$ 上的广义因子化被破坏（\tagV）。

\begin{table}[htbp]
\centering
\caption{本章陈述与附录~O、附录~N 中精确形式的对应}
\label{tab:dy-correspondence}
\footnotesize
\setlength{\tabcolsep}{4pt}
\renewcommand{\arraystretch}{1.2}
\begin{tabularx}{\textwidth}{@{}>{\raggedright\arraybackslash}p{3.0cm}>{\raggedright\arraybackslash}X>{\raggedright\arraybackslash}p{5.2cm}@{}}
\toprule
本章陈述 & 精确形式 & 差别（假设、误差类型）\\
\midrule
\cref{lem:dy-fsi}（\tagA） & \cref{lem:appO-unitarity}（\tagR，$R_{\mathrm{cond}}$）；割位置恒等式\cref{lem:appO-cut-position} & 前提\cref{asm:appO-lcopt}显式化；在 E-S 型近似之下为精确恒等式\\
\cref{lem:dy-poles-orientation}（\tagR） & \cref{lem:appO-orientation}（\cref{sec:appO-orientation}） & 内容相同；附录给出复化后 $\Imag D$ 的下界及与\cref{lem:method-i0-compat}的对照\\
\cref{ex:dy-future-vs-past}（\tagR） & \cref{prop:appO-past-wilson}证明的第 (2) 步 & 单个割图给出未来指向的线；过去指向只在单侧改写之后\\
\cref{prop:dy-deformation}（\tagA） & \cref{prop:appO-one-sided}（E-W）、\cref{prop:appO-target}（精确几何）、\cref{prop:appO-eikonal}（E-P）、\cref{prop:appO-past-wilson}（\tagA） & 已证明的步骤与假设\cref{asm:appO-mixing}、\cref{asm:appO-deformation}、\cref{asm:appO-envelope}、\cref{asm:appO-rewritten-cs}分开\\
\cref{lem:dy-soft}（\tagA） & \cref{prop:appO-past-wilson}(b)、\cref{lem:method-soft-unitarity} & 退耦为 Wilson 线之后，软因子为 $1$ 是 $R_{\mathrm{def}}$ 的\\
\cref{prop:dy-abelian-model}（\tagP） & \cref{prop:appO-ss-cancel}（\tagR，$R_{\mathrm{cond}}$） & 推广到任意 $J_A$、$J_B$ 与交换族；只含非重合连接\\
\cref{prop:dy-eikonal-twogluon}（\tagR） & 证明见\cref{sec:appO-wilson}；参数化\cref{lem:appO-glauber-operator}；\cref{prop:appO-strict-window-shape} & 附录补充严格窗口对正规子形状的依赖\\
\cref{rem:dy-common-ordering}(iii) & \cref{prop:appO-twoexch-spectator}（\tagP，模型定理） & 复合旁观者、未割线、$\order{g^4}$；被割的直接线开放\\
\cref{thm:dy-factorization}（\tagA） & \cref{thm:appO-main}$\to$\cref{cor:appO-interface}$\to$\cref{prop:appN-dy}（均 \tagR，$R_{\mathrm{cond}}$） & 全部前提逐条列出；误差类型由 E-W 经接口化为 E-N\\
\bottomrule
\end{tabularx}
\end{table}

\begin{tcolorbox}[enhanced,breakable,title=本章涉及的常见误解,fonttitle=\bfseries,
  colback=gray!3!white,colframe=gray!55!black,boxrule=0.6pt,arc=2pt,
  left=5pt,right=5pt,top=3pt,bottom=3pt,before skip=6pt,after skip=6pt]
\small
\begin{itemize}[leftmargin=1.2em,itemsep=0pt,topsep=0pt,parsep=0pt]
\item “Glauber 胶子在单举截面中总是相消”：相消要求对旁观者的末态求和，并且测量在割线两侧作为恒等算符出现（\cref{lem:dy-fsi}的地位段落）；对 TMD 的 $pp\to h_1h_2X$ 与强子—强子硬衍射，它不成立（\cref{thm:break-tmd,prop:break-hh-diffraction}）。
\item “NNLO/N$^3$LO 与数据一致，说明 Glauber 相消已被检验”：部分子截面里没有旁观者（\cref{rem:dy-nnlo-not-glauber}）。
\item “把 $k^-$ 围道推到 $|k^-|\sim|\mathbf k_T|$ 就离开了 Glauber 区”：$|k^+k^-|$ 仍是 $\order{\lambda^3}$，须沿 $e_G$ 同时形变两个分量（\cref{prop:dy-deformation}）。
\item “非阿贝尔双 Glauber 交换留下非幺正残差”：那是严格截断的假象，精确的图给出 $\Real p_{\rm opp}=0$（\cref{prop:dy-eikonal-twogluon}，\cref{rem:dy-common-ordering}）。
\item “单个割图中的软 Wilson 线已经是指向过去的”：单个割图给出指向未来的线（\cref{ex:dy-future-vs-past}）。
\item “对 $q_T$ 微分的 DY，软因子也是 $1$”：只对共线单举截面成立；$q_T$ 谱的软因子进入 TMD 分布（\cref{lem:dy-soft}的地位段落）。
\item “把 $\mu_F$、$\mu_R$ 变化两倍得到的就是理论误差”：这是约定，因子化只保证残余依赖是高阶的（\cref{prop:dy-scale}之后的讨论）。
\item “余项是 $\Lambda^2/Q^2$，因为高扭度算符的维数高二”：这是 OPE 的陈述；对 DY，余项的幂次是猜想，重整化子分析只是证据（\cref{thm:dy-factorization}的地位段落，\cref{sec:dy-twist}）。
\item “所有强子—强子过程都像 DY 一样因子化”：色单态末态为 \tagA，有色末态只是 \tagC（\cref{thm:dy-colour-singlet}，\cref{conj:dy-coloured-final}）。
\end{itemize}
\end{tcolorbox}

后文的用途：阈值行为是重求和的出发点（\cref{prop:thr-oneloop}）；对 $q_T$ 微分的 DY，Glauber 分析给出指向过去的订书钉形 Wilson 线（\cref{thm:tmd-factorization,thm:tmd-sivers-sign}）；测量相容性判据与因子化的破坏见\cref{prop:break-measurement}，Glauber 有效理论见\cref{prop:front-syz}。

